\def \titolo {Appunti di Fisica 2}
\documentclass[11pt,oneside]{scrbook} % KOMA: pulito e flessibile

% ---------- LINGUA & CODIFICA ----------
\usepackage[italian]{babel}
\usepackage[T1]{fontenc}
\usepackage[utf8]{inputenc}

% ---------- TIPOGRAFIA CONSIGLIATA ----------
\usepackage{lmodern}     % font semplice e diffuso
\usepackage{microtype}   % micro-tipografia, migliore resa

% ---------- GEOMETRIA PAGINA ----------
% Scegli *una* delle due geometrie qui sotto.
% (1) Impostazioni consigliate per A4 (Italia) con ampio margine esterno:
\usepackage[
  a4paper,
  top=24mm,
  bottom=28mm,
  inner=24mm,   % margine interno (verso rilegatura)
  outer=30mm,   % margine esterno: ampio per le note a margine
  heightrounded
]{geometry}
\setlength{\headheight}{17pt} % evita sovrapposizioni con l'header
\setlength{\headsep}{6mm}     % distanza tra header e corpo

\usepackage{scrlayer-scrpage}
\clearpairofpagestyles
\ihead{\leftmark}
\ofoot*{\pagemark}
\automark{chapter}

% (2) Se vuoi replicare il formato del PDF caricato (US Letter 8.5x11"):
% \usepackage[
%   letterpaper,
%   top=25mm,
%   bottom=30mm,
%   inner=30mm,
%   outer=55mm,
%   heightrounded
% ]{geometry}

% ---------- STRUMENTI DI STRUTTURA ----------
\setcounter{secnumdepth}{3}
\setcounter{tocdepth}{3}

% ---------- STILE TITOLI IN STILE KOMA ----------
\setkomafont{chapter}{\normalfont\huge\bfseries}
\setkomafont{section}{\normalfont\Large\bfseries}
\setkomafont{subsection}{\normalfont\large\bfseries}
\setkomafont{subsubsection}{\normalfont\normalsize\bfseries}

% ---------- STILE FRONTESPIZIO ----------
\setkomafont{title}{\sffamily\bfseries\Huge}
\setkomafont{author}{\sffamily\large}
\setkomafont{publishers}{\sffamily\normalsize}
\setkomafont{date}{\large}

\RedeclareSectionCommand[beforeskip=1.5ex,afterskip=1.5ex]{chapter}
\RedeclareSectionCommand[beforeskip=1ex,afterskip=0.8ex]{section}
\RedeclareSectionCommand[beforeskip=0.8ex,afterskip=0.6ex]{subsection}
\RedeclareSectionCommand[beforeskip=0.6ex,afterskip=0.4ex]{subsubsection}

% ---------- AMBIENTI UTILI ----------
\usepackage{amsmath,amssymb,amsthm}
\usepackage{booktabs}
\usepackage{array}
\usepackage{graphicx}
\usepackage{hyperref}
\hypersetup{  % stile dei riferimenti
  colorlinks=true,
  linkcolor=black,
  citecolor=blue,
  urlcolor=blue
}
\usepackage{xcolor}
\usepackage{tcolorbox}
\usepackage{subcaption} 
\usepackage{upgreek} % \uptau

\usepackage{mathrsfs} % \mathscr{}

\usepackage{tikz}
\usetikzlibrary{intersections}
\usepackage{pgfplots} % grafici, creato in tikZ
\pgfplotsset{compat = 1.18}
\usetikzlibrary{
arrows.meta,
positioning,
calc,
decorations.markings,
angles,
quotes,
3d,
patterns
}

\usepackage{subcaption}

\usepackage[european]{circuitikz} % circuiti

\usepackage{comment} % \begin{comment}  \end{comment}

\usepackage{mathtools} % \coloneqq

\usepackage{siunitx}
\sisetup{separate-uncertainty=true}
%----------- TEOREMI E STRUTTURE VARIE (KOMA, AMSTHM)

% Stile "teorema" (corsivo)
\theoremstyle{plain}
\newtheorem{theorem}{Teorema}[chapter]
\newtheorem{proposition}[theorem]{Proposizione}
\newtheorem{lemma}[theorem]{Lemma}
\newtheorem{corollary}[theorem]{Corollario}

% Stile "definizione" (testo dritto)
\theoremstyle{definition}
\newtheorem{definition}[theorem]{Definizione}

% Stile "nota" (più discorsivo)
\theoremstyle{remark}
\newtheorem{remark}[theorem]{Osservazione}
\newtheorem{example}[theorem]{Esempio}


% ---------- NOTE A MARGINE (parole-chiave/sidenote) ----------
% \marginpar stampa nel margine esterno con oneside.
% Definiamo un comando comodo e graficamente coerente:
\newcommand{\mkeyword}[1]{%
  \marginpar{\raggedright\footnotesize\emph{#1}}%
}

% Se preferisci le note a sinistra, attiva la riga seguente:
% \reversemarginpar

% ---------- PARAGRAFI E SPAZI ----------
\setlength{\parindent}{0.8em}  % rientro tradizionale
\setlength{\parskip}{0.25em}   % piccolo spazio extra tra paragrafi

%----------- COMANDI AGGIUNTIVI -----------
% --- Comando per vettori colonna di lunghezza variabile ---

\newcommand{\colvec}[1]{%
  \begin{pmatrix}
    #1
  \end{pmatrix}%
}

% --- comando derivata --- parziale \der[ordine]{funzione}{variabile}
\newcommand{\der}[3][]{%
  \frac{\partial^{#1} #2}{\partial #3^{#1}}%
}

% --- comando versore --- 
\newcommand{\vers}[1]{
  \hat{\boldsymbol{#1}}
}

% --- epsilon ---
\newcommand{\eps}[0]{
  \varepsilon
}

% --- norma ---
\newcommand{\norm}[1]{
  \left\lVert #1 \right\rVert
}

% --- prodotto scalare ---
\newcommand{\sc}[1]{
  \langle #1 \rangle
}

% --- vettore ---
\renewcommand{\v}[1]{
  \mathbf{#1}
}

% ---------- BIBLIOGRAFIA ----------
\usepackage{csquotes} % \enquote{}
\usepackage[backend=biber,style=numeric,sorting=nyt]{biblatex}
\addbibresource{bibliografia.bib} % nome del file .bib

% ---------- FRONTESPIZIO ----------
\title{\textbf{\titolo}}
\author{Cosimo Pascarella}
\date{Anno Accademico 2025--2026}

\publishers{
  Email: \text{c.pascarella2@studenti.unipi.it}\\
  Versione: v0.1
}

\begin{document}
\frontmatter
\maketitle
\tableofcontents
\cleardoublepage
\markboth{}{}
\mainmatter

% \chapter*{Prefazione}

% ---------- INIZIO DOCUMENTO ----------

\chapter{Elettrostatica}
Questo capitolo è stato aggiunto a posteriori rispetto alla trattazione degli appunti soltanto per completezza e ripasso, è da considerare dunque una bozza 
volutamente incompleta e rapida. Tuttavia credo che tornare sugli argomenti base una volta essersi addentrati nella materia sia utile per comprendere le 
nozioni più profonde che all'inizio non vengono adeguatamente indagate.

È noto che oggetti carichi elettricamente esercitano delle forze fra loro, determinate dalla carica, dal segno, dalla distanza. Sperimentalmente sappiamo 
che la loro interazione è descritta dalla cosiddetta \textit{forza di Coulomb}. Siano due cariche $ q_1, q_2 $ con posizioni $ \v{r_1}, \v{r_2} $, allora 
la forza che la carica $ q_2 $ esercita sulla carica $ q_1 $ è 
\begin{align}
  \v{F}_{1,2} = k_e \frac{q_1 \, q_2 \, \left( \v{r}_1 - \v{r}_2 \right)}{|\v{r}_1 - \v{r}_2|^3}.
\end{align}
Analizziamo questa espressione: la costante $ k_e $ vale 
\begin{align*}
  k_e = 
  \begin{cases}
    \DS \frac1{4\pi \eps_0} \qquad & \textnormal{[SI]} \\
    1 \qquad & \textnormal{[cgs]}
  \end{cases}
\end{align*}
dove SI indica il Sistema internazionale di unità di misura, mentre cgs indica il sistema Centimetro-Grammo-Secondo, anche detto sistema Gaussiano. La 
forza è inoltre antisimmetrica per le due cariche, infatti si osserva facilmente che 
\begin{align}
  \v{F}_{2,1} = - \v{F}_{1,2}
\end{align}
Notiamo la dipendenza quadratica al denominatore della distanza fra le due cariche, esattamente come la forza di attrazione gravitazionale, ma che a 
differenza di quest'ultima possiamo avere anche delle forze repulsive dovute alla possibilità di avere segni delle cariche discordi, cosa che nella forza 
gravitazionale non è possibile.

\section{Campo elettrostatico}
Abbiamo discusso la forza di Coulomb nel caso di due cariche, proviamo quindi ad aggiungerne una terza e otteniamo che per il principio di 
sovrapposizione la forza risultante sulla carica $ q_1 $ è data dalla somma dei due contributi:
\begin{align}
  \v{F}_{1, 2\,3} = k_e \frac{q_1 \, q_2 \, \left( \v{r}_1 - \v{r}_2 \right)}{|\v{r}_1 - \v{r}_2|^3} + 
  k_e \frac{q_1 \, q_3 \, \left( \v{r}_1 - \v{r}_3 \right)}{|\v{r}_1 - \v{r}_3|^3}
\end{align}
Generalizziamo considerando un sistema di $ q,q_1,\dots,q_N $ cariche, la forza esercitata sulla carica $ q $ dalle altre è
\begin{align}
  \v{F}_q = k_e \, q\, \sum_{i = 1}^N  \frac{q_i \, \left( \v{r}_1 - \v{r}_i \right)}{|\v{r}_1 - \v{r}_i|^3}.
\end{align}
Definiamo campo elettrico $ \v{E(r)} $ il rapporto
\begin{align}
  \v{E(r)} \coloneq \frac{\v{F}_q}{q}
\end{align}
la cui generalizzazione a un sistema continuo di cariche si basa sul sostituire la carica i-esima con una densità di carica moltiplicata per un volume 
infinitesimo $ q_i \mapsto \rho(\v{r}) \,d^3r' $ e integrare sullo spazio (dove la densità di carica è non nulla):
\begin{align}
  \v{E(r)} = k_e\, \int \frac{\rho(\v{r'}) \, \left( \v{r}_1 - \v{r'} \right)}{|\v{r}_1 - \v{r'}|^3} \,d^3r'
\end{align}
Nella gravitazione classica il campo gravitazionale viene considerato come uno ``artificio'' matematico costruito a partire dalla forza e il motivo 
di questa considerazione è dovuto principalmente a due ragioni: la prima è storica, lo sviluppo della teoria gravitazionale da parte di Newton precede 
la maggior parte dei lavori che hanno costruito la teoria elettromagnetica non disponendo quindi degli strumenti matematici necessari; la seconda è 
dovuta all'intensità della forza infatti $ k_e \sim 10^9 $ mentre $ G_N \sim 10^{-11} $, dunque fare misure che riguardano il campo risulta difficile senza 
degli strumenti sofisticati. Il punto del discorso è che l'elettrodinamica è la prima teoria considerabile come \textit{teoria di campo}, dove questo 
termine è dilagato ampiamente nella fisica più avanzata ma a noi basta sapere che è considerata tale perché i campi sono delle entità fisiche tanto quanto 
lo sono le forze e le cariche, e il campo elettrico in particolare è considerato il mediatore dell'interazione fra cariche.

\section{Teorema di Gauss}
Dopo aver introdotto il campo elettrostatico vogliamo calcolarne il flusso attraverso una superficie chiusa che racchiude le cariche, per farlo sfrutteremo 
il teorema della divergenza:
\begin{align}
  \Phi(\v{E}(\v{r}))_S = \oint_S \v{E} \cdot \,d\v{S} = \int_V \nabla\!\cdot\!\v{ E } \,dV
\end{align}
adesso dobbiamo trovare la divergenza del campo elettrostatico e per farlo sfruttiamo il risultato noto che 
\begin{align}
  \nabla \cdot \left( \frac{\v{r' - r''}}{|\v{r'-r''}|^3} \right) = 4\pi \delta(\v{r' - r''})
\end{align}
dove $ \delta(\v{x}) $ è la delta di Dirac. In questo modo abbiamo
\begin{align}
  \nabla\!\cdot\!\v{ E } &= 4\pi k_e  \, \rho(\v{r'}) = \frac{\rho(\v{r'})}{\eps_0} 
  \\
  \implies \Phi(\v{E(r)}) &= \int_V \frac{\rho(\v{r'})}{\eps_0} \,dV = \frac{Q_{\text{int}}}{\eps_0}
\end{align}
dove $ Q_{\text{int}} $ è la carica interna al volume $ V $. Il risultato che abbiamo ottenuto è noto come teorema di Gauss ed è un risultato utile al 
fine di svolgere integrali lunghi con alta probabilità di errore. Noi abbiamo mostrato come l'espressione sperimentale della forza di Coulomb implica 
il  teorema di Gauss, ma può essere dimostrato vale anche il viceversa, ovvero che assumendo il teorema allora l'unica forza possibile è l'espressione 
di Coulomb.

Per quanto riguarda il flusso attraverso una superficie qualsiasi non necessariamente chiusa sfruttiamo la definizione di angolo solido:
\begin{align}
  d\Omega \coloneq \frac{dA}{R^2} = \sin\theta \,d\theta\, d\varphi
\end{align}
cioè l'angolo solido corrispondente ad una superficie infinitesima circolare $ dA $ è il differenziale d'area diviso per il raggio al quadrato dal centro, 
per una superficie generica dato che sono infinitesime basta moltiplicare per il $ \cos\theta $. Tramite questa definizione possiamo trovare l'elemento 
infinitesimo di flusso su una superficie infinitesima come
\begin{align}
  &d\Phi = \v{E}\cdot \vers{n} dA = E \, \cos\theta \, dA = E\, R^2 \,d\Omega
  \\
  &\implies \Phi_S(\v{E}) = \int_\Omega E \, R^2 \,d\Omega
\end{align}
L'intero angolo solido vale $ 4\pi $ quindi quando facciamo il flusso su una superficie chiusa del campo di una carica otteniamo
\begin{align}
  \Phi_S(\v{E}) = \oint_\Omega \frac{q}{4\pi \eps_0} \,d\Omega = \frac{q}{\eps_0}
\end{align}
e per la linearità dell'integrale con un sistema di cariche abbiamo che il flusso del campo elettrico totale su una superficie chiusa è la somma degli 
integrali sugli angoli solidi visti da ogni carica, i quali sono nulli nel caso la carica sia esterna alla superficie, dunque è immediato ottenere il 
teorema di Gauss.


\section{Potenziale elettrostatico}
Definendo il campo elettrico abbiamo visto che questo dipende dalla distanza fra le cariche e il punto di osservazione, che è indipendente dal tempo e 
che è centrale, queste caratteristiche sono sufficienti per dire che è conservativo, dato che il lavoro lungo una linea chiusa $ \gamma $ è nullo. Questo 
implica l'esistenza di una funzione potenziale che indicheremo con $ V(\v{r}) $ \footnote{In alcuni testi è indicato con $ \phi(\v{r}) $ o $ \Phi(\v{r}) $}:
\begin{align}
  \oint_\gamma \v{E} \cdot d\v{l} = 0 \iff \exists\,V:\;\v{E(r)} = - \nabla V(\v{r})
\end{align}
Per ottenere una espressione per il potenziale consideriamo l'integrale curvilineo del campo (che è un integrale solo radiale) cioè
\begin{align}
  V(\v{r}_B) - V(\v{r}_A) = - \int_{\v{r}_A}^{\v{r}_B} \v{E} \cdot d\v{l} = - \int_{r_A}^{r_B} \v{E} \cdot \vers{r} \,dr
\end{align}
Nel caso del campo elettrostatico di una distribuzione di carica abbiamo che dalla definizione di potenziale riconosciamo una primitiva della funzione 
$ \frac{|\v{r - r'}|}{|\v{r - r'}|^3} = - \nabla \left( \frac{1}{|\v{r - r'}|} \right) $, così il potenziale diventa
\begin{align} \label{eq:_potenziale_elettrostatico}
  V(\v{r}) = k_e \int \frac{\rho(\v{r'})}{|r - r'|} \,d^3r'
\end{align}
il potenziale è chiaramente definito a meno di una costante e in questo caso abbiamo posto $ V(\v{r} \to \infty) = 0 $.
Nel definire il potenziale abbiamo trovato due importanti relazioni del campo elettrostatico cioè
\begin{align}
  &\oint_S \v{E} \cdot d\v{S} = \frac{Q_{\text{int}}}{\eps_0}
  \\
  &\oint_\gamma \v{E} \cdot d\v{l} = 0
\end{align}
che sono note come equazioni di Maxwell dell'elettrostatica in forma globale. Possiamo esprimerle in forma locale tramite il teorema della divergenza e il 
teorema del rotore:
\begin{align}
  &(\nabla \cdot \v{E})(\v{r}) = \frac{\rho(\v{r})}{\eps_0}
  \\
  &\nabla \times \v{E} = 0 
\end{align}

\section{Equazione di Poisson}
Dalla definizione di potenziale come la funzione scalare il cui gradiente (cambiato di segno) è il campo, possiamo applicare l'operatore divergenza dato che 
la prima equazione di Maxwell mette in relazione la divergenza del campo con la densità di carica, quindi
\begin{align}
  \nabla \cdot \nabla V(\v{r}) = - \nabla \cdot \v{E(r)} = -\frac{\rho(\v{r})}{\eps_0}
\end{align}
e l'operatore ``divergenza di gradiente'' è chiamato \textit{operatore di Laplace}, anche detto \textit{laplaciano}, che indicheremo con il simbolo 
$ \nabla^2 $. In questi termini la relazione appena ricavata diventa 
\begin{align}
  \nabla^2 V(\v{r}) = -\frac{\rho(\v{r})}{\eps_0}
\end{align}
che è nota come \textit{equazione di Poisson}. La corrispettiva omogenea dell'equazione di Poisson è detta \textit{equazione di Laplace}:
\begin{align}
  \nabla^2V(\v{r}) = 0
\end{align}
Entrambe queste due equazioni sono delle cosiddette PDE (partial differential equation), cioè delle equazioni alle derivate parziali lineari del secondo 
ordine, la cui soluzione è tutt'altro che banale. L'equazione di Poisson è la diretta applicazione della matematica alle equazioni di Maxwell, ma 
il motivo per cui abbiamo introdotto anche l'equazione di Laplace è perché nel caso di una distribuzione di carica (indipendente dal tempo) abbiamo già 
una espressione del potenziale trovata in equazione \eqref{eq:_potenziale_elettrostatico}, che consiste in una soluzione particolare. Quindi se troviamo 
una soluzione all'equazione omogenea possiamo scrivere la soluzione all'equazione di Poisson come somma della particolare e della soluzione omogenea, il 
che ci permette di trovare il campo elettrico applicando semplicemente il gradiente (meno dispendioso di fare un integrale di volume). Tuttavia risolvere 
l'equazione di Laplace non è sempre semplice e dipende molto dalle simmetrie del problema, infatti troveremo dei modi di scrivere le soluzioni ogni volta 
diversi in base al sistema di coordinate che scegliamo di impiegare per risolvere il problema.

\section{Energia elettrostatica}
Dalla conservatività del campo elettrico abbiamo potuto definire un potenziale elettrostatico, adesso vogliamo sfruttare ciò per definire una energia 
associata alla forza elettrostatica.

Il lavoro per portare una carica $ q $ dall'`infinito'' a un punto $ \v{R} $ è dato da
\begin{align}
  W_q = \int_{+\infty}^\v{R} \v{F} \,d\v{l} = q \int_{+\infty}^\v{R} \v{E}\cdot \vers{r} \,dr = q \left( V(+\infty) - V(\v{R}) \right)
\end{align}
se vediamo il potenziale come se generato da una distribuzione discreta di $ N $ cariche otteniamo che il lavoro sulla carica $ q $ è 
\begin{align}
  W_q = \frac{q}{4\pi \eps_0} \sum_{i= 1}^N \frac{q_i}{|\v{r_i - r}|}
\end{align}
Interpretiamo questo lavoro come l'energia potenziale della carica $ q $; calcoliamo l'energia potenziale di tutto il sistema di $ N $ cariche, cioè 
considerando la mutua interazione di ognuna:
\begin{align}
  U = k_e \sum_{i} U_i = k_e \sum_{j > i} \frac{q_i\, q_j}{|\v{r_i - r_j}|} 
  = k_e \cdot \frac{1}{2} \sum_{j \neq i} \frac{q_i\, q_j}{|\v{r_i - r_j}|}
\end{align}
dove abbiamo sfruttato la simmetria dell'interazione fra le singole cariche. Il passaggio al continuo è diretto ed è
\begin{align}
  U = \frac{1}{8\pi\eps_0} \int \int \frac{\rho(\v{x}) \rho(\v{x'})}{|\v{x - x'}|} \,d^3x \,d^3x'
\end{align}
e notiamo che un integrale può essere espresso come il potenziale elettrostatico di una distribuzione di cariche quindi otteniamo
\begin{align}
  U = \frac{1}{2} \int \rho(\v{x}) V(\v{x}) \,d^3x
\end{align}
Da questa forma possiamo sfruttare l'equazione di Poisson appena ottenuta per eliminare la $ \rho(\v{x}) $ dall'integrale:
\begin{align}
  U = -\frac1{8\pi k_e} \int V\, \nabla^2V \,d^3x = -\frac1{8\pi k_e} \int \nabla \cdot (V \nabla V) - |\nabla V|^2 \,d^3x 
  \\
  = -\frac1{8\pi k_e} \left( \oint_S V\nabla V \,d^2x -\int |\v{E}|^2 \,d^3x \right) = \frac1{8\pi k_e} \int |\v{E}|^2 \,d^3x
\end{align}
dove abbiamo integrato per parti e sfruttato l'identità vettoriale
\begin{align}
  V \nabla^2V = \nabla \cdot (V \nabla V) - |\nabla V|^2
\end{align}
Nella forma che abbiamo ottenuto riconosciamo una densità di energia elettrostatica $ u_e $:
\begin{align}
  u_e \coloneq \frac1{8\pi k_e} |\v{E}|^2 \; \implies \;
  U = \int u_e \,d^3x
\end{align}
dove l'integrale è svolto stavolta in tutto lo spazio essendo limitato al campo elettrostatico e non più alla distribuzione di carica $ \rho(x) $.
Il problema è che in questa forma l'energia potenziale $ U $ contiene sia l'energia di mutua interazione che l'auto-energia (\textit{self energy}), 
la quale per cariche puntiformi chiaramente diverge. Nella distribuzione discreta avevamo risolto questo problema semplicemente non contando i contributi 
in cui $ i = j $, ma adesso per ottenere la consueta forma dell'energia elettrostatica occorre fare una distinzione: 
$ U = U_{\textit{self}} + U_{\text{int}} $. Sarà per ora nostra accortezza eliminare accuratamente i termini di self energy.

\section{Teorema della media per le funzioni armoniche}
Introduciamo ora un risultato che ci sarà utile in seguito che è noto come il teorema del valor medio per le funzioni armoniche. 

Una funzione $ f :\Omega \subseteq \mathbb{R}^n \to \mathbb{R} $ differenziale fino al secondo ordine, cioè $ f \in C^2(\Omega) $, si dice armonica 
in $ \Omega $ se:
\begin{align}
  \nabla^2f = 0
\end{align}
Il teorema afferma che se $ f \in C^2(S) $ è armonica, dove $ S $ è una superficie sferica (o un volume), allora la sua media su $ S $ è uguale al valore 
nel centro, cioè
\begin{align}
  \left. \nabla^2 f \right|_S = 0 \implies \langle f \rangle_S = f(c)
\end{align}
Questo risultato è particolarmente importante per il caso dei conduttori: per definizione in un conduttore all'equilibrio il campo elettrico è ortogonale 
alla superficie, questo perché se non lo fosse le cariche si sposterebbero ottenendo una nuova configurazione di equilibrio. Il potenziale sulla superficie 
è quindi costante, ma allora applicando il teorema della media abbiamo che il potenziale è costante anche all'interno del conduttore, che ci porta dire che 
il campo elettrico all'interno è nullo. In realtà si potrebbe obiettare che non possiamo applicare il teorema della media a meno che il conduttore non sia 
una sfera, tuttavia questo si può aggirare tramite una forma locale del teorema nota come \textit{principio del massimo}. Tale principio dice che se $ f $ 
è una funzione armonica in $ \Omega $ allora questa ammette massimi e minimi solo sul bordo $ \partial\Omega $, altrimenti è costante. Applicandolo a ogni 
punto del conduttore si mostra che il potenziale deve essere costante in ogni punto.

\section{Formulazione del potenziale tramite funzione di Green}
Nell'ultima sezione abbiamo introdotto i conduttori e le loro proprietà principali, fra cui il potenziale elettrostatico costante sul bordo. Descrivere da 
un punto di vista matematico questi sistemi è particolarmente complicato se si vuole considerare come si redistribuiscono le cariche libere all'interno del 
conduttore in modo da soddisfare le condizioni al bordo. Tuttavia possiamo sfruttare dei risultati matematici noti come \textit{identità di Green} che 
permettono di scrivere il potenziale in termini di una certa funzione di Green $ G $.


\subsection{Identità di Green}

Siano $f$ e $g$ due funzioni scalari sufficientemente regolari definite su un dominio
$\Omega \subseteq \mathbb{R}^3$ con bordo $\partial\Omega$.
Partendo dall'identità vettoriale
\begin{equation}
    \nabla \cdot (f \nabla g) = f \nabla^2 g + \nabla f \cdot \nabla g,
\end{equation}
e integrando su $\Omega$ con il teorema della divergenza, si ottiene la
\emph{prima identità di Green}:
\begin{equation}\label{eq:green1}
    \int_\Omega \left( f \nabla^2 g + \nabla f \cdot \nabla g \right) d^3x
    = \oint_{\partial\Omega} f \, \frac{\partial g}{\partial n} \, dS,
\end{equation}
dove $\partial/\partial n = \hat{\v{n}} \cdot \nabla$ indica la derivata lungo la
normale uscente a $\partial\Omega$.
Sottraendo la~\eqref{eq:green1} scritta con $f$ e $g$ scambiati, si ottiene la
\emph{seconda identità di Green}:
\begin{equation}\label{eq:green2}
    \int_\Omega \left( f \nabla^2 g - g \nabla^2 f \right) d^3x
    = \oint_{\partial\Omega} \left( f \, \frac{\partial g}{\partial n}
      - g \, \frac{\partial f}{\partial n} \right) dS.
\end{equation}
Questa identità è il punto di partenza per la costruzione della soluzione formale
del problema di Poisson.

\subsection{Definizione della funzione di Green}

Si consideri il problema di Poisson con condizioni al contorno su $\Omega$:
\begin{equation}\label{eq:poisson}
    \nabla^2 V(\v{x}) = -\frac{\rho(\v{x})}{\varepsilon_0}, \qquad
    \v{x} \in \Omega.
\end{equation}
La \emph{funzione di Green} $G(\v{x}, \v{x}')$ del dominio $\Omega$ è la soluzione del
problema ausiliario
\begin{equation}\label{eq:green_def}
    \nabla^2 G(\v{x}, \v{x}') = \delta^{(3)}(\v{x} - \v{x}'), \qquad \v{x} \in \Omega,
\end{equation}
con opportune condizioni al contorno su $\partial\Omega$ (specificate nelle
sottosezioni successive). Fisicamente, $G(\v{x},\v{x}')$ rappresenta il potenziale
(a meno di $\varepsilon_0$) generato da una carica puntiforme unitaria posta in
$\v{x}'$, in presenza del dominio $\Omega$.

La soluzione fondamentale dell'equazione di Laplace in $\mathbb{R}^3$,
\begin{equation}
    \phi(\v{x}, \v{x}') = -\frac{1}{4\pi |\v{x} - \v{x}'|},
\end{equation}
soddisfa già la~\eqref{eq:green_def} su tutto $\mathbb{R}^3$, ma non rispetta in
generale le condizioni al contorno del dominio $\Omega$.
Si decompone quindi $G$ come
\begin{equation}
    G(\v{x}, \v{x}') = -\frac{1}{4\pi|\v{x}-\v{x}'|} + H(\v{x}, \v{x}'),
\end{equation}
dove $H$ è una funzione regolare (armonica) in $\Omega$, scelta in modo che la
condizione al contorno su $G$ sia soddisfatta.

\subsection{Derivazione della formula integrale per il potenziale}

Si applica la seconda identità di Green~\eqref{eq:green2} con $f = V(\v{x})$ e
$g = G(\v{x}, \v{x}')$, trattando $\v{x}'$ come parametro:
\begin{equation}
    \int_\Omega \!\left( V(\v{x})\, \nabla^2 G(\v{x},\v{x}') - G(\v{x},\v{x}')\,
    \nabla^2 V(\v{x}) \right) d^3x
    = \oint_{\partial\Omega} \!\left( V \, \frac{\partial G}{\partial n}
    - G \, \frac{\partial V}{\partial n} \right) dS.
\end{equation}
Sostituendo le equazioni~\eqref{eq:poisson} e~\eqref{eq:green_def}:
\begin{equation}
    \int_\Omega V(\v{x})\, \delta^{(3)}(\v{x}-\v{x}')\, d^3x
    + \frac{1}{\varepsilon_0}\int_\Omega G(\v{x},\v{x}')\, \rho(\v{x})\, d^3x
    = \oint_{\partial\Omega} \left( V \, \frac{\partial G}{\partial n}
    - G \, \frac{\partial V}{\partial n} \right) dS.
\end{equation}
La delta di Dirac seleziona $V(\v{x}')$. Scambiando infine i nomi
$\v{x} \leftrightarrow \v{x}'$ si ottiene la \emph{formula integrale generale}
per il potenziale:
\begin{equation}\label{eq:V_green}
    \boxed{
    V(\v{x}) = \frac{1}{\varepsilon_0}\int_\Omega G(\v{x},\v{x}')\,\rho(\v{x}')\,d^3x'
    - \oint_{\partial\Omega} \left( V(\v{x}')\, \frac{\partial G}{\partial n'}(\v{x},\v{x}')
    - G(\v{x},\v{x}')\, \frac{\partial V}{\partial n'}(\v{x}') \right) dS'
    }
\end{equation}
I due contributi hanno interpretazione distinta: il primo termine è il potenziale
generato dalla distribuzione di carica $\rho$ nell'interno del dominio; il secondo
termine codifica l'effetto delle condizioni al contorno su $\partial\Omega$.

\subsection{Il problema della sovradeterminazione e le condizioni al contorno}

La formula~\eqref{eq:V_green} contiene, sull'integrale di superficie, sia il valore
di $V$ che quello di $\partial V / \partial n$ su $\partial\Omega$.
Specificare \emph{entrambe} le quantità costituisce un problema \emph{sovradeterminato}:
l'equazione di Poisson è una PDE del second'ordine, e per garantire esistenza e
unicità della soluzione è sufficiente assegnare \emph{una sola} condizione su
ciascuna porzione del bordo.
Imporre entrambe le condizioni è in generale incompatibile con la soluzione
dell'equazione, oppure — circolarmente — richiede di conoscere già la soluzione per
ricavare la derivata normale, vanificando l'utilità della formula stessa.

La scelta della condizione al contorno opportuna definisce due problemi classici:

\begin{itemize}
    \item \textbf{Condizioni di Dirichlet.} Si assegna il valore del potenziale sulla
    superficie:
    \begin{equation}
        V(\v{x})\big|_{\partial\Omega} = V_0(\v{x}'), \qquad \v{x}' \in \partial\Omega.
    \end{equation}
    È il caso fisico dei conduttori a potenziale fissato.
    Si elimina il termine con $\partial V/\partial n'$ dalla~\eqref{eq:V_green}
    imponendo $G(\v{x},\v{x}')\big|_{\v{x}' \in \partial\Omega} = 0$,
    così il secondo termine dell'integrale di superficie si annulla e la formula
    si semplifica in:
    \begin{equation}
        V(\v{x}) = \frac{1}{\varepsilon_0}\int_\Omega G(\v{x},\v{x}')\,\rho(\v{x}')\,d^3x'
        - \oint_{\partial\Omega} V_0(\v{x}')\, \frac{\partial G}{\partial n'}(\v{x},\v{x}')\, dS'.
    \end{equation}

    \item \textbf{Condizioni di Neumann.} Si assegna la derivata normale del potenziale
    sulla superficie:
    \begin{equation}
        \frac{\partial V}{\partial n}\bigg|_{\partial\Omega} = \Phi(\v{x}'),
        \qquad \v{x}' \in \partial\Omega,
    \end{equation}
    che fisicamente equivale ad assegnare la componente normale del campo elettrico
    (e quindi la densità di carica superficiale $\sigma = -\varepsilon_0\,\partial V/\partial n$).
    In questo caso si sceglie $G$ tale che
    $\partial G/\partial n'\big|_{\partial\Omega} = \text{cost}$
    (il valore costante è fissato da una condizione di compatibilità globale),
    e il termine con $V_0$ scompare dall'integrale di superficie.
\end{itemize}

\subsection{Unicità della soluzione}

\subsubsection{Enunciato e significato fisico}

Il problema di Poisson con condizioni al contorno di Dirichlet \emph{o} di Neumann
ammette una soluzione \emph{unica} (a meno di una costante additiva nel caso di
Neumann). L'unicità ha un ruolo fondamentale in elettrostatica: garantisce che il
metodo delle immagini, il metodo delle serie di Fourier o qualunque altra tecnica
che individui \emph{una} soluzione compatibile con le condizioni al contorno fornisca
automaticamente \emph{la} soluzione del problema, senza bisogno di verificare che
non ne esistano altre.

\subsubsection{Dimostrazione}

Supponiamo che $V_1$ e $V_2$ siano due soluzioni del problema di Poisson con le
stesse condizioni al contorno (di Dirichlet o di Neumann). Definiamo la differenza
\begin{equation}
    u(\v{x}) = V_1(\v{x}) - V_2(\v{x}).
\end{equation}
Per linearità dell'equazione di Poisson, $u$ soddisfa l'equazione di Laplace
\begin{equation}
    \nabla^2 u = 0 \qquad \text{in } \Omega,
\end{equation}
con condizioni al contorno omogenee:
\begin{align}
    u\big|_{\partial\Omega} &= 0 \qquad \text{(Dirichlet)}, \\
    \frac{\partial u}{\partial n}\bigg|_{\partial\Omega} &= 0 \qquad \text{(Neumann)}.
\end{align}
Applichiamo la prima identità di Green~\eqref{eq:green1} con $f = g = u$:
\begin{equation}
    \int_\Omega \left( u\, \nabla^2 u + |\nabla u|^2 \right) d^3x
    = \oint_{\partial\Omega} u\, \frac{\partial u}{\partial n}\, dS.
\end{equation}
Poiché $\nabla^2 u = 0$, il primo termine dell'integrale di volume scompare.
Il termine di superficie è nullo in entrambi i casi:
nel caso di Dirichlet perché $u = 0$ su $\partial\Omega$;
nel caso di Neumann perché $\partial u / \partial n = 0$ su $\partial\Omega$.
Rimane quindi:
\begin{equation}
    \int_\Omega |\nabla u|^2 \, d^3x = 0.
\end{equation}
Poiché l'integrando $|\nabla u|^2 \geq 0$ è continuo, l'unica possibilità è
\begin{equation}
    \nabla u = \v{0} \quad \text{in tutto } \Omega,
\end{equation}
ovvero $u$ è costante in $\Omega$. Nel caso di Dirichlet la condizione $u|_{\partial\Omega}=0$
fissa la costante a zero, dunque $V_1 = V_2$ ovunque. Nel caso di Neumann la costante
rimane arbitraria, il che è fisicamente atteso: il campo $\v{E} = -\nabla V$ è determinato
in modo unico, ma il potenziale è definito a meno di una costante additiva globale.

\chapter{Equazione di Laplace}
Dall'elettrostatica sappiamo che il campo elettrico in forma locale deve soddisfare
due equazioni differenziali, cioè
\begin{align}
  \nabla\cdot\v{ E } = \frac{\rho}{\varepsilon} \quad \v{\nabla}\times\v{ E } = 0
\end{align}
e abbiamo visto che la seconda ci permette di definire un funzione scalare detta potenziale tale che il suo 
gradiente sia il campo elettrico cambiato di segno
\begin{equation}
  \v{ E } = -\nabla V
\end{equation}
Applicando il fatto che $ \v{ E } $ sia conservativo alla prima equazione di Maxwell otteniamo
\begin{align} \label{eq:Poisson}
  && \nabla^2 V = -\frac{\rho}{\varepsilon} \tag*{[Equazione di Poisson]}
\end{align}
cioè
\begin{equation}
  \der[2]{V}{x} + \der[2]{V}{y} + \der[2]{V}{z}
\end{equation}
dove $ \nabla^2 $ è l'operatore di Laplace (laplaciano). Si tratta di una equazione differenziale alle 
derivate parziali ed, in generale, non omogenea. Dato che nell'equazione la funzione che fa da "termine noto" è la 
densità di carica (calcolata nel punto $ \v{ x } $, dove stiamo valutando le derivate di $ V $) nei casi in 
cui la densità sia nulla l'equazione diventa omogenea ed è chiamata Equazione di Laplace
\begin{align} \label{eq:Laplace}
  && \nabla^2 V = 0 \tag*{[Equazione di Laplace]}
\end{align}
che rispetto alla generica equazione di Poisson, possiamo trovare delle tecniche per risolverla in base al sistema di 
coordinate che scegliamo di utilizzare. Nei prossimi paragrafi risolveremo l'equazione di Laplace tramite la 
separazione delle variabili, ma è importante notare che le funzioni che otterremo non sono tutte le soluzioni 
dell'equazione ma soltanto una parte (cioè quelle scrivibili come prodotto di funzioni di una singola variabile)-

\section{Coordinate cartesiane}
In un sistema cartesiano $ xyz $ è possibile separare le variabili imponendo l'\textit{ansatz}:
\begin{equation}
  V = X(x) Y(y) Z(z)
\end{equation}
Calcolandone il laplaciano si ottiene
\begin{align}
  \nabla^2 \left( X Y Z \right) = Y Z X'' + X Z Y'' + X Y Z'' \overset{!}{=} 0
\end{align}
dividiamo per $ XYZ $
\begin{equation}
  \frac{X''}{X} + \frac{ Y''}{ Y} + \frac{ Z''}{ Z} = 0
\end{equation}
Adesso sfruttiamo il fatto che ognuno dei termini deve azzerarsi in modo indipendente, questo perché abbiamo 
imposto che $ V $ sia il prodotto di tre funzioni ognuna di una singola variabile (se non si annullassero in modo indipendente, 
quando una si annulla faccio variare l'altra e ottengo che non è più $ 0 $)
\begin{align}
  Z'' &= m^2 Z
  \\
  X'' &= -a^2 X
  \\
  Y'' &= -b^2 Y
\end{align}
dove per i coefficienti abbiamo imposto che
\begin{equation}
  -m^2 = -a^2 -b^2 = -(a^2 + b^2)
\end{equation}
\textbf{SE} vogliamo imporre $ V_{\infty} = 0 $ allora scartiamo l'esponenziale crescente in $ Z(z) $
\begin{align}
  Z (z) &= C e^{- \sqrt{a^2 + b^2} z}
  \\
  X(x) &= E \cos(ax) + F \sin (ax)
  \\
  Y(y) &= G \cos(by) + H \sin (by)
\end{align}
Per determinare i coefficienti vanno imposte le condizioni specifiche del problema.
\subsection{Caso bidimensionale}
Problema della striscia con potenziale $ 0 $ su ogni bordi tranne che sul bordo $ 0<x<a $ che vale $ \phi(x) $.
Allora il potenziale si scrive come 
\begin{equation}
  V(x,y) = \sum_{n = 0}^{\infty} A_n \sin\left( \frac{n\pi}{a} x \right) e^{-\frac{n\pi}{a} y}
\end{equation}
dove abbiamo già imposto la condizione di periodo $ T = 2a \implies \sin (\alpha (x + T)) = \sin(\alpha x + 2\pi) $ quindi $ \alpha T = 2\pi $
cioè $ \alpha = \frac{\pi}{a} $, e abbiamo già scartato il coseno che non si poteva annullare agli estremi.\\
Manca solo da imporre la condizione del potenziale fissato $ V (x,0) = \phi(x) $ 
\begin{equation}
  \left( V(x,0) = \sum_{n = 0}^{\infty} A_n \sin\left( \frac{n\pi}{a} x \right) = \phi(x) \right) \sin\left(  \frac{n'\pi}{a} x \right)
\end{equation}
e integrando sfruttiamo la ortogonalità dei seni:
\begin{equation}
  \int_{-T/2}^{T/2} \sin \left( \frac{2\pi n}{T} \right) \sin \left( \frac{2\pi n'}{T} \right)  \,dx = \delta_{nn'} \frac{T}{2}
\end{equation}


\section{Coordinate sferiche}
In coordinate sferiche per risolvere il laplaciano proviamo l'\textit{ansatz}
\begin{equation}
  V(r, \theta, \varphi) = \frac{U (r)}{r} P(\theta) Q(\varphi)
\end{equation}
ne calcoliamo il laplaciano e si ottiene
\begin{align}
  Q (\varphi) &\propto e^{im\varphi} \quad \text{con} \quad m \in \mathbb{Z}
  \\
  P(\theta) &= P_l(\cos\theta) \quad \text{polinomi di Legendre}
  \\
  \frac{U(r)}{r} &= A_l r^l + \frac{B_l}{r^{l+1}}
\end{align}
\textbf{SE} imponiamo simmetria azimutale, cioè $ m = 0 $ si ottiene
\begin{equation}
  V(r, \theta) = \sum_{l=0}^{\infty} \left( A_l \, r^l + \frac{B_l}{r^{l+1}} \right) P_l\left( \cos \theta \right)
\end{equation}

\textbf{SE NON} c'è simmetria azimutale, cioè $ m\neq0 $ allora
\begin{equation}
  P(\theta) Q(\varphi) \to Y_{l,m} (\theta, \varphi) = e^{im\varphi} P_l^m(\sin\theta,\cos\theta)
\end{equation}
dove $ P_l^m $ sono i polinomi di Legendre generalizzati.

Per i polinomi di Legendre classici è possibile ricavarli tramite la successione che si ottiene risolvendo Laplace, oppure esiste una formula diretta 
per il calcolo
\begin{equation}
  P_l(x) = \frac{1}{2^l \, l!} \,\frac{d^l}{dx^l} \left[ \left( x^2 -1 \right)^l \right] \tag*{[Formula di Rodrigues]}
\end{equation}
I polinomi di Legendre formano una base ortonormale nello spazio $ \mathbb{L}^2 $ delle funzioni a quadrato sommabile, la normalizzazione vale
\begin{equation}
  \langle P_l, P_{l'} \rangle = \int_{-1}^{1} P_l P_{l'}  \,dx = \delta_{ll'} \frac{2}{2l+1} 
\end{equation}

\subsection{Armoniche sferiche}
Il potenziale si scrive come
\begin{equation}
  V = \sum_{l=0}^{\infty} \sum_{m=-l}^{l} \left( A_l r^l + \frac{B_l}{r^{l+1}} \right) Y_{l,m}
\end{equation}
dove vale che
\begin{equation}
  Y_{l,m} = (-1)^m Y_{l,-m}^*
\end{equation}
mentre l'ortogonalità è normalizzata, integrando sull'angolo solido
\begin{equation}
  \int_{4\pi}  \,d\Omega \,Y_{l,m} Y_{l',m'}^*  = \int_0^{2\pi} \,d\varphi \int_0^{\pi} \sin\theta  \,d\theta \, Y_{l,m} Y_{l',m'}^* 
  = \delta_{ll'} \delta_{mm'} 
\end{equation}

\[
Y_\ell^m(\theta,\varphi) = N_{\ell m}\, P_\ell^m(\cos\theta)\, e^{i m \varphi}
\]
\\l=0
\[
Y_0^0 = \frac{1}{\sqrt{4\pi}}
\]
\\
l=1
\[
\begin{aligned}
Y_1^{-1} &= \sqrt{\frac{3}{8\pi}}\,\sin\theta\, e^{-i\varphi} \\
Y_1^{0}  &= \sqrt{\frac{3}{4\pi}}\,\cos\theta \\
Y_1^{+1} &= -\sqrt{\frac{3}{8\pi}}\,\sin\theta\, e^{i\varphi}
\end{aligned}
\]
\\
l=2
\[
\begin{aligned}
Y_2^{-2} &= \sqrt{\frac{15}{32\pi}}\,\sin^2\theta\, e^{-2i\varphi} \\
Y_2^{-1} &= \sqrt{\frac{15}{8\pi}}\,\sin\theta\cos\theta\, e^{-i\varphi} \\
Y_2^{0}  &= \sqrt{\frac{5}{16\pi}}\,(3\cos^2\theta - 1) \\
Y_2^{+1} &= -\sqrt{\frac{15}{8\pi}}\,\sin\theta\cos\theta\, e^{i\varphi} \\
Y_2^{+2} &= \sqrt{\frac{15}{32\pi}}\,\sin^2\theta\, e^{2i\varphi}
\end{aligned}
\]
\\
l=3:
\[
\begin{aligned}
Y_3^{-3} &= \sqrt{\frac{35}{64\pi}}\,\sin^3\theta\, e^{-3i\varphi} \\
Y_3^{-2} &= \sqrt{\frac{105}{32\pi}}\,\sin^2\theta\cos\theta\, e^{-2i\varphi} \\
Y_3^{-1} &= \sqrt{\frac{21}{64\pi}}\,\sin\theta\,(5\cos^2\theta - 1)\, e^{-i\varphi} \\
Y_3^{0}  &= \sqrt{\frac{7}{16\pi}}\,(5\cos^3\theta - 3\cos\theta) \\
Y_3^{+1} &= -\sqrt{\frac{21}{64\pi}}\,\sin\theta\,(5\cos^2\theta - 1)\, e^{i\varphi} \\
Y_3^{+2} &= \sqrt{\frac{105}{32\pi}}\,\sin^2\theta\cos\theta\, e^{2i\varphi} \\
Y_3^{+3} &= -\sqrt{\frac{35}{64\pi}}\,\sin^3\theta\, e^{3i\varphi}
\end{aligned}
\]

\chapter{Esercizi dell'elettrostatica}


\section{Distribuzione di carica}
Si consideri la distribuzione di carica seguente indipendente da $ y $:
\begin{align}
  \rho(x, |z| \leq \frac{L}{2}) = \rho_0 \cos^2(\frac{\pi z}{L}) \cos (\frac{2\pi x}{a}) \quad \text{e} \quad \rho(x, |z| > \frac{L}{2}) = 0
\end{align}
Trovare il potenziale elettrostatico in tutto lo spazio.

\textit{Soluzione}\\
Ragioniamo intanto sulle simmetrie del problema: $ \rho $ è simmetrica per inversione su $ x $ e su $ z $, inoltre c'è invarianza traslazionale 
lungo $ y $. Questo ci dice che il potenziale non può dipendere da $ y $, e per le simmetrie possiamo considerare il semispazio $ x > 0 $. 
Le equazioni da risolvere sono:
\begin{align}
  \nabla^2 V = 0 \quad &\text{per} \quad z > \frac{L}{2} \\
  \nabla^2 V = -\frac{\rho}{\varepsilon_0} \quad &\text{per} \quad z \leq \frac{L}{2} 
\end{align}
Le condizioni che abbiamo sono che
\begin{itemize}
  \item $ V(x, z \to \infty) = 0 $
  \item $ V(x,z) = V(x,-z) $ \hfill [parità in $ z $] 
  \item $ V(x, z \to \frac{L}{2}^{+}) = V(x, z \to \frac{L}{2}^{-}) $ \hfill [continuità del potenziale]
  \item $ (\partial_z V_F - \partial_z V_{in})|_{|z| = \frac{L}{2}} = 0 $ \hfill [continuità del campo elettrico]
\end{itemize}
Fuori abbiamo che $ V $ è armonico quindi la soluzione è della forma
\begin{equation}
  V(x,z) = \sum_{n = 0}^{\infty} C_n e^{-\alpha z} \cos(\alpha x)
\end{equation}
dove abbiamo già imposto che il potenziale vada a $ 0 $ in $ z $ (esponenziale negativo), e in proviamo una soluzione simile alla forma della $ \rho $, 
quindi proviamo il coseno.

Dentro dobbiamo risolvere l'equazione di Poisson; per farlo proviamo l'\textit{ansatz}
\begin{equation}
  V(x,z) = f(z) \cos (\frac{2\pi x}{a})
\end{equation}
calcoliamone il laplaciano:
\begin{align}
  \nabla^2 V &= (\nabla^2 f) \cos (\frac{2\pi x}{a}) + f(z) \nabla^2 \left( \cos (\frac{2\pi x}{a}) \right) \\
  &= \frac{\partial ^2 f(z) }{\partial ^2 z } \cos (\frac{2\pi x}{a}) - f(z) \frac{4\pi^2}{a^2} \cos (\frac{2\pi x}{a}) 
  \overset{!}{=} -\frac{\rho_0}{\varepsilon_0} \cos^2(\frac{\pi z}{L}) \cos (\frac{2\pi x}{a})
\end{align}
dobbiamo quindi risolvere
\begin{equation}
  f''(z) = f(z) \frac{4\pi^2}{a^2} -\frac{\rho_0}{\varepsilon_0} \cos^2(\frac{\pi z}{L})
\end{equation}
di cui la soluzione omogenea la sappiamo ed è un esponenziale, mentre la per la particolare possiamo sfruttare una proprietà del coseno e tentare 
un \textit{ansatz} di quella forma:
\begin{align}
  &\text{\textbullet} f_{om} = A e^{\frac{2\pi z}{a}} + B e^{-\frac{2\pi z}{a}}
  \\
  &\text{\textbullet} \cos^2(\frac{\pi z}{L}) = \frac{\cos\left( \frac{2\pi z}{L} \right) + 1}{2}  \implies \quad f_p = C \cos\! \left( \frac{2\pi z}{L} \right) + D \tag*{[ansatz]}
\end{align}
Troviamo i coefficienti $ C $ e $ D $:
\begin{align}
  &-C \frac{4\pi^2}{L^2} \cos \! \left( \frac{2\pi}{L} z \right)  = C \frac{4\pi^2}{L^2} \cos \! \left( \frac{2\pi}{L} z \right) + D \frac{4\pi^2}{a^2} 
  - \frac{\rho_0}{\varepsilon_0} \frac{ \cos\!\left( \frac{2\pi}{L}z \right) +1 }{2}
  \\
  &\cos\!\left( \frac{2\pi}{L} z \right) \left( -C \frac{4\pi^2}{L^2} -C \frac{4\pi^2}{a^2} + \frac{\rho_0}{2\varepsilon_0}\right) = D\frac{4\pi^2}{a^2} - \frac{\rho_0}{2\varepsilon_0}
\end{align}
che deve azzerarsi $ \forall z $, quindi
\begin{align}
  &\text{\textbullet} -C \frac{4\pi^2}{L^2} -C \frac{4\pi^2}{a^2} + \frac{\rho_0}{2\varepsilon_0} = 0 &&\implies C = \frac{\rho_0}{8\pi^2 \varepsilon_0} \frac{L^2 a^2}{L^2 + a^2} &&
  \\
  &\text{\textbullet} D\frac{4\pi^2}{a^2} - \frac{\rho_0}{2\varepsilon_0} = 0 &&\implies D = \frac{\rho_0}{8\pi^2 \varepsilon_0} a^2 &&
\end{align}
Abbiamo adesso la forma del potenziale interno:
\begin{equation}
  V_{in}(x,z) = \left[ A e^{\frac{2\pi z}{a}} + B e^{-\frac{2\pi z}{a}} 
  + \frac{\rho_0}{8\pi^2 \varepsilon_0} a^2 \left( \frac{L^2}{L^2 + a^2} \cos\! \left( \frac{2\pi z}{L} \right) + 1  \right)\right]
  \cos (\frac{2\pi x}{a})
\end{equation}
Imponiamo la condizione di parità su $ z $, che ci porta a $ A = B $, quindi
\begin{equation}
  V_{in}(x,z) = \left[ A \left( e^{\frac{2\pi z}{a}} + e^{-\frac{2\pi z}{a}} \right)
  + \frac{\rho_0}{8\pi^2 \varepsilon_0} a^2 \left( \frac{L^2}{L^2 + a^2} \cos\! \left( \frac{2\pi z}{L} \right) + 1  \right)\right]
  \cos (\frac{2\pi x}{a})
\end{equation}
Per determinare il potenziale manca da imporre le condizioni di continuità del potenziale e della sua derivata lungo $ z $. Siccome fuori il potenziale è 
scritto ancora come serie di Fourier, queste ultime condizioni ci dicono che fuori c'è soltanto il termine $ n = 1 $ (perché i termini sono una base), 
e per ragioni di simmetria deve avere lo stesso periodo lungo $ x $.
\begin{equation}
  V_F (x,z) = C_1 e^{-\frac{2\pi}{a}z} \cos \left( \frac{2\pi}{a} x \right)
\end{equation}
Per trovare $ A $ e $ C_1 $ dobbiamo risolvere il sistema; $ k := \frac{\pi L}{a} $
\begin{align}
  &\text{\textbullet} C_1 e^{-k} = 2A \cosh (k) + \frac{\rho_0}{8\pi^2 \varepsilon_0} a^2 \left( 1 - \frac{L^2}{L^2 + a^2} \right)
  \\
  &\text{\textbullet} -C_1 e^{-k} = 2A \sinh(k)
\end{align}


\section{Cilindro con carica}
Una carica $ Q $ è uniformemente distribuita nel volume di un cilindro di raggio $ a $ e altezza $ h $, inoltre al centro del cilindro è posta una 
carica puntiforme $ -Q $. Si determini il potenziale elettrostatico a grandissima distanza dal cilindro e si discuta il risultato ottenuto al variare 
del rapporto $ h/a $.

\textit{Soluzione, SISTEMA CGS}\\
Analizziamo prima le simmetrie del sistema: c'è simmetria cilindrica e per inversione di $ z $ (cioè riflessione rispetto $ xy $). Quindi possiamo 
scrivere il potenziale come espansione in polinomi di Legendre (non c'è dipendenza da $ \varphi $), ma solo di ordine pari perché deve rispettare 
$ V(z) = V(-z) $, che visto nei polinomi diventa
\begin{align}
  \theta \mapsto \pi - \theta \implies P_l(\cos\theta) = P_l(\cos(\pi-\theta)) = P_l (-\cos\theta) 
\end{align}
che è possibile solo con potenze pari di $ \cos\theta $. Il fatto che ci sia una carica $ Q $ e una $ -Q $ ci dice anche che il termine di monopolo è 
nullo, quindi il primo termine non nullo è il quadrupolo ($ l=2 $) e quindi della forma
\begin{equation}
  V_{\text{quad}}(r, \theta) = \frac{B_2}{r^3} \frac{3\cos^2\theta - 1}{2}
\end{equation}
Facciamo emergere questo fatto dai calcoli (per facilitare i conti mi posso mettere lungo l'asse $ z $, tanto il coefficiente che dobbiamo trovare è
sempre $ B_2 $):
\begin{align}
  V(z) = \int \frac{\rho(\v{ r }')}{\v{ |r-r'| } } \,d^3r' && \frac{1}{\v{ |r-r'| }} = \sum_{l=0}^{\infty} \frac{r_{<}^{l}}{r_{>}^{l+1}} \; P_l \left( \cos \theta \right)
\end{align}
espandendo
\begin{align}
  &r_< = r' , \, r_> = r \quad \text{e} \quad \rho(\v{ r' }) = \frac{Q}{\pi a^2 h} -Q \delta (\v{ r }) = \rho_0 - Q \delta(\v{ r })
  \\
  &V(r) = \frac{Q}{r} - \frac{Q}{r} + \int \rho_0 \frac{r'^2}{r^3} \frac{3\cos^2\gamma - 1}{2} \,d^3r' + \dots
\end{align}
dove $ \cos\gamma $ si riferisce all'angolo fra $ \v{ r }' $ e $ \v{ r } $, quindi $ r' \cos \gamma = \rho $ (variabile cilindrica, non densità 
di carica)
\begin{align}
  V(r) &= \frac{Q}{\pi a^2 h} \frac{2\pi}{r^3} \int_{-h/2}^{h/2} \int_{0}^{a}  \left( \frac{3}{2} z'^2 - \frac{\rho'^2 + z'^2}{2} \right) \rho' d\rho dz'
  \\
  &= \frac{Q}{a^2 h} \frac{2}{r^3}\, 2 \,\left[ \frac{a^2}{6} \left( \frac{h}{2} \right)^3 - \frac{a^4}{8} \frac{h}{2} \right] = \frac{Q}{12 r^3} \left( h^2 - 3a^2 \right)
\end{align}
questo lungo l'asse $ z $, quindi in generale vale
\begin{equation}
  V(r,\theta) = \frac{Q}{12} \left( h^2 - 3a^2 \right) \frac{1}{r^3} \frac{3\cos^2\theta - 1}{2}
\end{equation}


\chapter{Magnetostatica}
Da un punto di vista storico, il magnetismo era un fenomeno che consisteva in una attrazione e repulsione in punti diversi fra alcuni materiali 
detti per questo motivo "ferromagnetici". In una ottica più moderna, come il campo elettrico è il mediatore dell'interazione fra cariche, possiamo 
definire il campo magnetico $\v{B}$ come il mediatore dell'interazione fra correnti. Questo aspetto sarà fondamentale per casi dinamici.
L'equazione di continuità

\begin{equation} \label{eq:continuità}
  \nabla\cdot \v{J} + \frac{\partial\rho}{\partial t} = 0
\end{equation}
è l'espressione matematica del principio di conservazione della carica: se in volume la carica è diminuita, allora deve essere per forza uscita sotto forma 
di corrente, e questo si traduce in 
\begin{align}
  -\der{}{t} \left( \int_V \rho(\v{r'}) \,d^3r' \right) = - \der{Q}{t} = I = \oint_S \v{J} \cdot d\v{S} 
\end{align}
dove il segno meno è dovuto al fatto che se la carica diminuisce allora la corrente uscente è positiva (la normale $ \vers{n} $ alla superficie è uscente). 
Quella che abbiamo ottenuto è detta forma globale e per passare a quella locale ci basta osservare che la relazione
\begin{align}
  \int_V \nabla \cdot \v{J} + \der{\rho}{t} \,d^3r' = 0
\end{align}
deve valere per ogni volume $ V $, per questo motivo l'integrando deve essere identicamente nullo, ottenendo così la \eqref{eq:continuità}.
Nella statica abbiamo che
\begin{align}
  \frac{\partial\rho}{\partial t} = 0 \implies\nabla\cdot\v{J} = 0
\end{align}
e quindi le correnti che otteniamo sono dette \textit{stazionarie} perché integrando abbiamo
\begin{align}
  I = \int_S \v{J} \cdot d\v{S} = \int_V \nabla \cdot \v{J} \, d^3r' = 0
\end{align}
cioè che il bilancio totale delle correnti entranti e uscenti in un volume è nullo.

\section{Interazione fra correnti}
Nel 1820 Ørsted scopre che fili percorsi da corrente interagiscono con degli aghi magnetici facendoli ruotare.
Poco dopo Ampère osserva che le correnti interagiscono fra loro esercitando una forza sui fili che è proporzionale alle correnti e inversamente 
proporzionale alla distanza:
\begin{align}
  \frac{\Delta F}{\Delta L} = k\frac{I_1 I_2}{d}
\end{align}
dove la forza è attrattiva se le correnti scorrono nello stesso verso, altrimenti repulsiva.
La costante di proporzionalità vale $ k = \SI{2e-7}{H/m}$ e non fu misurata da Ampère ma soltanto in epoca moderna dopo che Maxwell introdusse il 
concetto di permeabilità magnetica $ \mu $. Fu deciso $ k = \frac{\mu_0}{2\pi}  $ in modo tale da ridefinire 1 Ampère come la corrente che scorrendo in 2 
fili alla distanza di 1 metro produce una forza di $ \SI{2e-7}{N} $ (2019).

Biot e Savart arrivano a una espressione del campo magnetico generato da un filo percorso da una corrente $ I $:
\begin{equation} \label{eq:B_integrale_di_dB}
  \v{B} =  k_m I \oint_C \frac{d\v{s}\times\v{u}_r}{r^2}
\end{equation}
l'integrale è su un circuito chiuso perché deve valere l'equazione di continuità. Nell'ultima espressione abbiamo introdotto la $ k_m $ che in analogia 
alla $ k_e $ definiamo come
\begin{align*}
  k_m = 
  \begin{cases}
    \DS \frac{\mu_0}{4\pi} \qquad & \textnormal{[SI]} \\
    \DS \frac1c \qquad & \textnormal{[cgs]}
  \end{cases}
\end{align*}
dove $ \mu_0 $ è la \textit{permeabilità magnetica del vuoto}, la quale per ora rappresenta soltanto una costante di proporzionalità ricavata in modo 
sperimentale. La legge di Biot-Savart ha anche una sua forma differenziale:
\begin{equation} \label{eq:dB}
  d \v{B} =  k_m I \frac{d\v{s}\times\v{u}_r}{r^2}
\end{equation}
per delle correnti che non scorrono in un filo ma sono localizzate possiamo sfruttare che
\begin{align}
  di\,d\v{ s } =  \v{ J } d\Sigma ds = \v{ J }d\uptau
\end{align}
e quindi riscrivere la \eqref{eq:dB} come
\begin{equation}
  d\v{ B } = k_m \frac{\v{ J }\times\v{ u_r }}{r^2}
\end{equation}
integrando si ottiene
\begin{equation}
  \v{ B } = k_m \int_\uptau \frac{\v{ J }\times\v{ u_r }}{r^2} d\uptau
\end{equation}


\subsection{Campo di un filo}
Consideriamo un filo di lunghezza $ l = 2a $ nel quale scorre una corrente stazionaria $ I $ (da un punto di vista della magnetostatica non è chiaro 
perché non sarebbe verificata l'equazione di continuità \eqref{eq:continuità}), vogliamo trovare il campo lungo il piano mediano del filo.
Siccome $ \v{ B } $ è uno pseudovettore e c'è simmetria dispari rispetto al piano $ xy $, il campo ha solo componente $ B_{\varphi} $.\\
Utilizziamo la \eqref{eq:dB} integrandola fra gli estremi del filo:
\begin{align}
  & d\v{ s } \times \v{ u }_r = dz \sin \theta \, \v{ u }_{\varphi}   \quad \sin\theta = \frac{R}{r} = \frac{R}{\sqrt{R^2 + z^2}}\\
  & d\theta \cos\theta = - R \frac{z}{\left( R^2 + z^2 \right)^{\frac{3}{2}}} dz \quad \text{ma}\quad \cos\theta =  \frac{z}{r} \\
  & \text{quindi} \quad dz = -\frac{r^2}{R} d\theta \implies \v{ B } 
  =  \frac{\mu_0 I}{4\pi} \int_{\theta_1}^{\pi-\theta_1} \frac{-\sin\theta}{R} \,d\theta \; \v{ u }_{\varphi} \\
  & = \frac{2\mu_0 I}{4\pi R} \int_{\frac{\pi}{2}}^{\pi-\theta_1} \left( -\sin\theta \right) \,d\theta 
  = \frac{\mu_0 I}{2\pi R} \cos \left( \pi -\theta_1 \right) 
  = \frac{\mu_0 I}{2\pi R} \frac{a}{\sqrt{R^2 + a^2}} \label{eq:campo_filo_lunghezza_a}
\end{align}


\subsection{Campo di un filo infinito} \label{section: B filo infinito}
Sia un filo rettilineo infinito dove scorre una corrente $ i $, il campo magnetico è dato dalla \eqref{eq:B_integrale_di_dB}, considerando il circuito 
chiuso come una semicirconferenza che ha come diametro il filo (infinito): siccome la circonferenza cresce $ \propto r $ mentre il campo 
decresce $ \propto \frac{1}{r^2} $, quindi sulla circonferenza è chiaro che si annulla per $ r \to \infty $.
L'integrale sul filo:
\begin{align}
  & d\v{ s }\times\v{ u }_r = dz\cos\theta, \quad z = R \tan\theta \implies dz = \frac{R}{\cos^2\theta}d\theta \\
  &\v{B} =  \frac{\mu_\text{0} i}{4\pi} \oint_C \frac{d\v{s}\times\v{u}_r}{r^2} \, 
  = \frac{\mu_\text{0} i}{4\pi} \int_{-\frac{\pi}{2}}^{+\frac{\pi}{2}} \frac{\cos\theta}{R}\,dx
  \, = \frac{\mu_0 i}{2\pi R}\quad \text{[Legge di Biot-Savart]} \label{eq:Biot-Savart}
\end{align}
Notiamo che lo stesso risultato si ottiene facendo il limite del filo di lunghezza $ l = 2a $ per $ a \to \infty $ nella \eqref{eq:campo_filo_lunghezza_a}:
\begin{align}
  &\frac{\mu_0 I}{2\pi R} \frac{a}{\sqrt{R^2 + a^2}} = \frac{\mu_0 I}{2\pi R} \frac{1}{\sqrt{1 + \left( \frac{R}{a} \right) ^2 }} 
  = \frac{\mu_0 I}{2\pi R} \left( 1 - \frac{1}{2} \left( \frac{R}{a} \right) ^2 + ... \right)
\end{align}

\subsection{Filo di raggio a}
Corrente uniforme J
\begin{equation}
  \oint_S \v{ B } d\v{ S } = \mu_0 I_{\text{conc}}
\end{equation}
ma la corrente concatenata dipende dal raggio della circonferenza chiusa che ho preso
\begin{equation}
  I_{\text{conc}} = J \left(  \frac{r}{a} \right)^2 \quad \text{se}\quad r<a
\end{equation}
Se il filo si muove con velocità $ \v{ v } $ parallela al suo asse.
$ \lambda = \frac{Q}{L} $, $ \v{ J } = \rho_Q \v{ v }$ ma $ \rho_Q = \frac{\lambda}{\Sigma} $
quindi $ I = \lambda \v{ v }\cdot \v{ n } $

\section{Equazioni di Maxwell per il campo magnetico} \label{section: eq Maxwell per B}
In generale per delle correnti descritte da una densità di corrente $ \v{ J(x) } $, il campo magnetico in un punto $ \v{ x } $ vale quella che 
è nota come \textit{legge di Ampère-Laplace}:
\begin{equation} 
  \v{ B(x) } = k_m \int \v{ J(x') } \times \frac{\v{ x-x' }}{\v{ |x-x'| }^3} \,d^3x' \tag*{[Ampère-Laplace]}
\end{equation}\label{B_Laplace}
che non è altro che la generalizzazione della legge di Biot-Savart a una densità di corrente. La forma della \eqref{B_Laplace} è analoga
all'espressione del campo elettrico mediante integrale di una densità di carica $ \rho $:
\begin{equation}
  \v{ E(x) } = k_e \int \rho(\v{ x }) \frac{\v{ x-x' }}{\v{ |x-x'| }^3}\,d^3x'
\end{equation}
Entrambe queste due espressioni sono così generali che risultano difficili da usare. Come abbiamo fatto per il campo elettrico, utilizziamo che
\begin{equation}
  \v{\nabla}\left(\frac{1}{|\v{ x-x'}|}\right) = - \frac{\v{ x-x' }}{|\v{ x-x' }|^3}
\end{equation}
e una proprietà del rotore di funzione scalare per vettore:
\begin{equation} \label{rot_proprietà}
  \v{\nabla}\times(a\v{ b }) = \nabla(a)\times\v{ b } + a\nabla\times\v{ b }
\end{equation}
così la \eqref{B_Laplace} si riscrive come
\begin{align}
  \v{ B } &= k_m \int \v{ J(x') }\times \v{\nabla}\left(\frac{-1}{|\v{ x-x'}|}\right)  \,d^3x'\\
  &= k_m \int \v{\nabla}\left(\frac{1}{|\v{ x-x'}|}\right) \times \v{ J(x') }   \,d^3x' \nonumber \\
  &=  \nabla \times \left( k_m \int \frac{\v{ J(x') } }{|\v{ x-x' }|} \,d^3x' \right) \equiv \nabla \times \v{ A(x) } \label{eq:rot_potenziale_vettore}
\end{align}
dove nell'ultima uguaglianza abbiamo sfruttato la \eqref{rot_proprietà} notando che il rotore è fatto rispetto 
alle coordinate $ \v{ x } $ e non rispetto alle $ \v{ x' } $, ma $ \v{ J(x') } $ è funzione delle coordinate primate 
(stiamo integrando sulle correnti nel volume), quindi il suo rotore è 0. Da qui segue la prima equazione di Maxwell:
\begin{align}
  \nabla\cdot\v{ B } &= 0 \quad \text{dato che} \quad \nabla\cdot \left( \v{ \v{\nabla}\times\v{ a } } \right) = 0 \quad \forall \v{a} \in C^1(\R^3,\R^3)
\end{align}
Per quanto riguarda la seconda, mostriamo che
\begin{equation}
  \v{\nabla}\times\v{ B } = 4\pi k_m \v{ J(x) } = \mu_0 \v{ J(x) }
\end{equation}
sfruttando la riscrittura di prima
\begin{align}
  \v{\nabla}\times\v{ B } = k_m \nabla \times \left( \nabla \times  \int \frac{\v{ J(x) } }{|\v{ x-x' }|} \,d^3x' \right)
\end{align}
per il rotore di rotore vale l'identità
\begin{equation}
  \nabla \times  \left(\nabla \times \v{ A } \right) = \nabla\left( \nabla\cdot \v{ A } \right) - \nabla^2 \v{ A } \quad \forall \v{ A } \in C^2
\end{equation}
dove l'ultimo Laplaciano è componente per componente cioè
\begin{align}
  \left( \nabla^2 \v{ A } \right) _i = \, \nabla^2 A_i  \tag*{[Laplaciano vettoriale]}
\end{align}
calcoliamo prima $ \nabla\cdot \v{ A(x) } $:
\begin{align}
  \nabla \cdot \v{ A(x) } = \nabla \cdot \left(k_m \int \frac{\v{ J(x') }}{\v{ |x-x'| }} \, d^3x' \right) 
  = k_m \int \v{ J(x') } \left( \nabla \frac{1}{\v{ |x-x'| }} \right) \,d^3x'
\end{align}
$ \v{ J } $ esce dal gradiente perché è rispetto alle coordinate non primate.
\begin{align}
  &= k_m \int \v{ J(x') } \left( -\nabla' \frac{1}{\v{ |x-x'| }} \right) \,d^3x' 
  = k_m \int \frac{\left( \nabla' \cdot \v{ J(x') } \right)}{\v{ |x-x'| }} - \nabla' \cdot \left( \frac{\v{ J(x') }}{\v{ |x-x'| }} \right) \,d^3x'\\
  &= 0 - k_m \oint_S \frac{\v{ J(x') }}{\v{ |x-x'| }} \,d\v{ S } = 0
\end{align}
per la localizzazione delle sorgenti mandando S a infinito e nel primo termine abbiamo anche usato l'equazione di continuità \eqref{eq:continuità}.
Quindi rimane solo il termine col laplaciano:
\begin{align}
  \left( \v{\nabla}\times\v{ B } \right)_i &= -\nabla^2\v{ A(x) }_i = k_m \int J_i(\v{ x' }) \left( -\nabla^2 \frac{1}{\v{ |x-x'| }} \right) \,d^3x'\\
  &= k_m \int J_i \v{ (x') } \; 4\pi \delta \left( \v{ x-x' } \right) \,d^3x' = 4\pi k_m J_i (\v{ x }) = \mu_0 J_i (\v{ x })
\end{align}
otteniamo quindi
\begin{align}
  \v{\nabla}\times\v{ B } = \mu_0 \v{ J(x) }
\end{align}
Per ottenere la forma globale possiamo integrare su una superficie $ S $ e applicare il teorema del rotore:
\begin{align}
  4\pi k_m \int_S \v{ J } \cdot \,d\v{ S } = \int_S \v{\nabla}\times\v{ B } \,d\v{ S } = \oint_{\partial S} \v{ B } \cdot \, d\v{ l }
\end{align}
Abbiamo ottenuto che la circuitazione del campo magnetico lungo una linea chiusa $ \partial S $ è il flusso di $ \v{ J } $ attraverso la 
superficie $ S $, cioè la corrente concatenata alla linea:
\begin{equation} \label{eq:circuitazione_B=i_concatenata}
  \oint_{\partial S} \v{ B } \cdot \, d\v{ l } = 4\pi k_m I_{\text{conc}} = \mu_0 I_{\text{conc}}
\end{equation}



\section{Potenziale vettore}
Proprio come il potenziale scalare (del campo elettrico e successivamente del campo magnetico) cui facendone il gradiente si ottiene il campo 
(cambiato di segno), definiamo come "potenziale vettore" un campo vettoriale sufficientemente differenziabile e regolare tale che il suo 
rotore è il campo magnetico
\begin{equation}
  \v{ B } = \v{\nabla}\times\v{ A(x) }
\end{equation}
Questo in generale è garantito che possiamo farlo perché la divergenza di $ \v{ B } $ è nulla.

Nella sezione \ref{section: eq Maxwell per B} avevamo già ottenuto $ \v{ B } $ come rotore di un altro campo; 
diciamo che la forma generale di $ \v{ A(x) } $ è la \eqref{eq:rot_potenziale_vettore} più un termine a rotore nullo, che possiamo sempre esprimere 
come il gradiente di un campo scalare $ \Psi $:
\begin{equation}
  \v{ A(x) } = k_m \int \frac{\v{ J(x') }}{\v{ |x-x'| }} \,d^3x' + \nabla \Psi (\v{ x })
\end{equation}
L'arbitrarietà di $ \Psi(\v{ x }) $ ci dice che:
\begin{itemize}
  \item $ \v{ A(x) } $ non è unico
  \item le equazioni di Maxwell sono invarianti per la trasformazione 
  \begin{align}
    \v{ A(x) } \mapsto \v{ A }' = \v{ A(x) } + \nabla \Psi \v{ (x) } \tag*{[trasformazione di gauge]}.
  \end{align}
\end{itemize}
Significa che siamo liberi di scegliere (questa scelta è detta \textit{scelta di gauge}) $ \Psi $ a seconda del problema che stiamo affrontando.
Applicando la seconda equazione di Maxwell per il campo magnetico si ottiene che
\begin{align}
  \v{\nabla}\times\v{ B } &= \v{\nabla}\times\v{ \left( \v{\nabla}\times\v{ A(x) } \right) } 
  = \nabla \left( \nabla \cdot \v{ A(x) } \right) - \nabla^2\v{ A(x) } \\
  &= - \nabla^2 \v{ A(x) } = \mu_0 \v{ J(x) }
\end{align}
e dall'elettrostatica sappiamo che la soluzione a questa equazione è
\begin{equation}
  \v{ A(x) } = k_m \int \frac{\v{ J(x') }}{\v{ |x-x'| }} \,d^3x' \qquad \Psi \; \text{scelto costante}
\end{equation}
che è ciò che abbiamo ottenuto prima partendo dalla legge di Ampère-Laplace, quindi è chiaro che
\begin{equation}
  \text{Ampère-Laplace} \iff \text{Equazioni di Maxwell}.
\end{equation}
\subsection{Possibili scelte di Gauge}
Ci sono varie scelte di gauge rilevanti in base all'ambito che si sta trattando, ad esempio in magnetostatica si utilizza spesso la cosiddetta 
\textit{gauge di Coulomb} che consiste nell'imporre
\begin{align}
  \nabla\!\cdot\!\v{ A } = 0,
\end{align}
un'altra scelta è la \textit{gauge simmetrica}:
\begin{align}
  \v{ A } = \frac{1}{2} \v{ B } \times \v{ r }.
\end{align}
La \textit{gauge di Landau}, che è simile alla gauge simmetrica ma impone una direzione preferenziale, ad esempio per $ \v{ B } = B \vers{z} $ 
uniforme si ha
\begin{align}
  \v{ A } = (0,Bx,0) \quad \text{oppure} \quad \v{ A } = (-By, 0,0).
\end{align}
Infine abbiamo la \textit{gauge di Lorenz}, che ha una particolare importanza in elettrodinamica e in relatività, dato che è si dimostra che è invariante 
per trasformazioni di Lorentz, e consiste in
\begin{align}
  \nabla\!\cdot\!\v{ A } = -\frac{1}{c^2} \der{V}{t}
\end{align}
dove $ c $ è la velocità della luce nel vuoto e $ V $ è il potenziale elettrico. Per la trattazione di questa gauge si rimanda alla 
sezione \ref{section:_soluzione_Maxwell_con_sorgenti}.


\section{Espansione in multipolo del campo magnetico} \label{section :B espansione multipolo}
In analogia col potenziale elettrostatico, voglia trovare il termine dominante nel calcolo del campo magnetico di una distribuzione di correnti nota.
Appurato che non esistono ancora (2025) evidenze sperimentali del monopolo magnetico, non potremo avere (come nel caso di una distribuzione di cariche) 
un termine dovuto alla "carica" magnetica netta. Ci aspettiamo quindi di avere come primo termine non nullo un termine che assomigli a quello di un dipolo
 elettrico, sempre per la somiglianza fra le espressioni dei potenziali elettrico e vettore. La quantità che vogliamo sviluppare è 
\begin{equation}
 \frac{1}{\v{ |x-x'| }} = \frac{1}{\v{ |x| }} + \frac{\v{ x \cdot x' }}{\v{ |x| }^3} + ...
\end{equation}
Prendiamo la componente i-esima del potenziale vettore

\begin{align}
  A_i(\v{ x }) &= k_m \int J_i(\v{ x' }) \left( \frac{1}{\v{ |x| }} + \frac{\v{ x' \cdot x }}{\v{ |x| }^3} + ... \right) \,d^3x' \\
  &= k_m \left[ \frac{1}{\v{ |x| }} \int J_i(\v{ x' }) \,d^3x' + \frac{\v{ x }}{\v{ x }^3} \cdot \int \v{ x' } J_i(\v{ x' }) \,d^3x' + ...\right]
\end{align}
Valutiamo il primo termine: vogliamo trovare una funzione tale che la sua divergenza è $ J_i(\v{ x' }) $ così da applicare il Teorema della divergenza
per ottenere un integrale sulla superficie (bordo) che racchiude le correnti. 
\begin{align}
  \nabla \cdot \left( x_i \v{ J(x) } \right) = \nabla x_i \cdot \v{ J(x) } + x_i\nabla \cdot \v{ J(x) } = J_i(\v{ x })
\end{align}
nell'ultima uguaglianza abbiamo sfruttato l'equazione di continuità \eqref{eq:continuità}. Applichiamo il Teorema della divergenza ($ \partial V = S $)
\begin{align}
  \int_V J_i(\v{ x' }) \,d^3x' = \int_V \nabla \cdot \left( x_i ' \, \v{ J(x') } \right)  \,dV = \int_S x_i ' \, \v{ J(x') } \,d\v{ S } = 0
\end{align}
è nullo perché il flusso di $ \v{ J } $ sulla superficie è $ 0 $ per l'ipotesi di correnti localizzate (le correnti non escono dalla superficie).\\
Valutiamo il secondo termine; senza perdita di generalità consideriamo la componente lungo $ \v{ u }_x $:
\begin{align}
  A_x(\v{ x }) = k_m \frac{1}{\v{ |x| }^3} \int J_x(\v{ x' }) \left( xx' + yy' + zz' \right) \,d^3x' \label{eq:A_x_intermedio}
\end{align}
e notiamo due cose:\\
(1)
\begin{align}
  &\nabla ' \left( (x')^2\v{ J(x') } \right) = 2x'J_x(\v{ x' }) + (x')^2 \nabla' \v{ J(x') } = 2x'J_x(\v{ x' })\\
  &\implies \int_V x'J_x \,d^3x' \overset{\text{T.divergenza}}{=} \int_S (x')^2\v{ J(x') } \,d\v{ S } = 0
\end{align}
(2)
\begin{align}
  & 0= \int_V \nabla'\left( y'x'\v{ J(x') } \right) \,d^3x' = \int_V y'J_x + x'J_y  \,d^3x'\\
  &\implies \int_V y'J_x \,d^3x' = \frac{1}{2} \left( \int_V y'J_x \,d^3x'- \int_V x'J_y \,d^3x' \right) \eqcolon -m_z
\end{align}
così la \eqref{eq:A_x_intermedio} si riscrive come:
\begin{align}
  A_x (\v{ x }) &= k_m \frac{1}{\v{ |x| }^3} \left[ y \left( \frac{1}{2} \int_V y'J_x - x'J_y \,d^3x' \right) 
  + z \left( \frac{1}{2} \int_V z'J_x - x'J_z \,d^3x' \right) \right]
  \\
  & = k_m \frac{1}{\v{ |x| }^3} \left[ -ym_z + zm_y \right] = k_m \frac{\left( \v{ m \times x } \right)_x}{\v{ |x| }^3} 
\end{align}
e quindi il secondo termine del potenziale vettore vale:
\begin{equation} \label{eq:A_termine_dipolo}
  \v{ A(x) } = k_m \frac{\left( \v{ m \times x } \right)}{\v{ |x| }^3} 
\end{equation}
e di conseguenza il campo magnetico vale
\begin{equation}
  \v{ B(x) } = k_m \nabla \times \left( \frac{ \v{ m \times x } }{\v{ |x| }^3}  \right) = ... 
  = k_m \frac{3 \v{ u }_r \left( \v{ u }_r \cdot  \v{ m } \right) - \v{ m }}{r^3}
\end{equation}
Quello che abbiamo appena definito con $ \v{m} $ prende il nome di \textit{dipolo magnetico}
\begin{align}
  m_i &= \frac{1}{2} \int \left( \v{ x' \times J(x') } \right)_i \,d^3x'
  \\
  &= \frac{1}{2} \epsilon_{ijk}\int x'_jJ_k(\v{ x' }) \,d^3x'
\end{align}
cioè 
\begin{align}
  \v{m} = \frac{1}{2} \int \left( \v{ x' \times J(x') } \right) \,d^3x'\\
\end{align}
\section{Potenziale scalare magnetico}
Dalla seconda equazione di Maxwell notiamo che in assenza di correnti il campo diventa conservativo, e questo ci permette di definire un potenziale 
magnetico scalare il cui gradiente cambiato di segno ci restituisce il campo:
\begin{equation}
  \v{\nabla}\times\v{ B } = 0 \implies \v{ B } = - \nabla W
\end{equation}
che calcolandone la divergenza si ottiene
\begin{equation}
  0 = \nabla\cdot\v{ B } = - \nabla^2 W
\end{equation}
In questo modo lo studio del campo magnetico si riconduce a risolvere l'equazione di Laplace con opportune condizioni al bordo, 
proprio come avevamo fatto per il potenziale elettrico.\\
Consideriamo il caso del filo infinito con corrente uniforme $ I $. Il campo magnetico a distanza R lo abbiamo calcolato nella sezione 
\eqref{section: B filo infinito}
e vale
\begin{equation}
  \v{ B } = \frac{\mu_0 I}{2\pi R} \v{ u }_{\varphi} 
\end{equation}
Se calcoliamo la circuitazione lungo un cammino chiuso che contiene l'origine, per la \eqref{eq:circuitazione_B=i_concatenata} si ha
\begin{equation}
  \mu_0 I_{conc} = \oint_{\gamma} \v{ B } d\v{ l } = \oint_{\gamma} \left( -\nabla W \right) d\v{ l } = W(A) - W(B)
\end{equation}
Possiamo prendere
\begin{equation}
  W = - \frac{\mu_0 I}{2\pi} \varphi \implies W(0) - W(2\pi) = \frac{\mu_0 I}{2\pi} 2\pi = \mu_0 I
\end{equation}
questo se il cammino chiuso comprende l'origine, altrimenti l'angolo iniziale e finale sono entrambi $ 0 $ e quindi $ W(A) - W(B) = 0 $ 
perché non c'è corrente concatenata. Questo potenziale funziona per questo caso.\\
È possibile dimostrare che un buon potenziale magnetico per una spira generica nello spazio è dato da
\begin{equation}
  W(\v{ r }) = - k_m I \Omega (\v{ r })
\end{equation}
dove $ \Omega (\v{ r }) $ è l'angolo solido che la spira sottende rispetto al punto $ \v{ r } $. Se il punto si trova all'interno della 
superficie (chiusa) allora l'angolo solido è $ 4\pi $, se è all'esterno è in generale difficile e spesso inutile. Nel caso della spira circolare, per un punto 
sull'asse l'angolo solido è quello di un disco e si può ottenere facilmente che
\begin{equation}
  \Omega = 2\pi (1 - \cos \theta) = 2\pi \left( 1 - \frac{z}{\sqrt{R^2 + z^2}} \right)
\end{equation}


\section{Energia del dipolo magnetico}
Trattare l'energia associata al campo magnetico è un compito che rientra nell'elettrodinamica, questo perché entrano in gioco fenomeni quali la legge 
di Faraday che fanno variare nel tempo le grandezze che trattiamo. Tuttavia possiamo definire l'energia di un dipolo magnetico $ \v{ m } $ immerso 
in un campo magnetico esterno $ \v{ B }_{\text{ext}} $ come
\begin{align}
  U = - \v{ m } \cdot \v{ B }_{\text{ext}} = - m { B }_{\text{ext}} \cos\theta
\end{align}
dove $ \theta $ è l'angolo compreso fa il dipolo e il campo esterno. Da questa espressione si riconosce il flusso nel caso la spira sia piana, infatti
\begin{align}
  \v{ m } = S I \vers{n} \implies U = - \Phi_S({\v{ B }}) I
\end{align}
dove $ S $ è l'area della spira. Questa definizione minima di energia magnetica vale nel caso in cui il dipolo non 
dipenda dalle coordinate, e di conseguenza è possibile ricavare la forza esercitata come
\begin{align}
  \v{ F } = \nabla\left( \v{ m } \cdot \v{ B }_{\text{ext}} \right).
\end{align}
Per una trattazione più completa dell'energia elettromagnetica si rimanda alla sezione \ref{section:_Onde_elettromagnetiche}.

\section{Forza magnetica}
\# BOZZA
\begin{equation} \label{eq:forza_magnetica}
  \v{ F } = \int_V \v{ J(x') } \times \v{ B(x') } \,d^3x'
\end{equation}
Di conseguenza il momento magnetico vale:
\begin{equation} \label{eq:momento_forza_magnetica}
  \v{ N } = \int_V \v{ x' \times \left( J(x') \times B(x') \right) } \,d^3x'
\end{equation}
Supponiamo di avere un campo magnetico esterno e una distribuzione di correnti, siamo interessati a conoscere la forza e il momento 
esercitati su tale distribuzione. Le formule scritte sopra sono di carattere generale ma spesso il campo 
risulta difficile o impossibile da integrare, possiamo quindi svilupparlo se varia di poco nella regione di nostro interesse.
Grazie alla formula di Taylor abbiamo
\begin{equation}
  B_i(\v{ x }) = B_i(0) + \v{ x } \cdot \nabla B_i(0) + ...
\end{equation}
La componente i-esima della forza è 
\begin{equation}
  F_i = \epsilon_{ijk} \left[ B_k(0)\int J_j(\v{ x' }) \,d^3x' + \int J_j(\v{ x' }) \v{ x' } \cdot \nabla B_k(0) \,d^3x' + ... \right]
\end{equation}
Per l'equazione di continuità il primo integrale si annulla, il secondo termine diventa
\begin{align}
  F_i = \epsilon_{ijk} \int J_j(\v{ x' }) \v{ x' } \cdot \nabla B_k(0) \,d^3x' = \epsilon_{ijk} \left( \v{ m } \times \nabla \right)_j B_k(0)
\end{align}
e quindi la forza è
\begin{align}
  \v{ F } = \left( \v{ m } \times \nabla \right) \times \v{ B } &= - \v{ B \times \left( m \times \nabla \right) } 
  = - \left( \v{ m \left( \nabla \cdot B \right) } - \nabla \v{ \left( B \cdot m \right) }\right) \\
  &= \nabla \left( \v{ B \cdot m } \right) - \v{ m \left( \nabla \cdot B \right)} \\
  & \implies \quad \v{ F } = \nabla \left( \v{ m \cdot B } \right)
\end{align}

Il momento magnetico per un dipolo magnetico vale
\begin{equation}
  \v{ \uptau } = \v{ m } \times \v{ B }
\end{equation}
ma per ordini successivi conviene sviluppare in multipoli l'energia 
magnetica data da
\begin{equation}
  U = -\frac{1}{2c} \int \v{ J }(\v{ r' }) \v{ A_{\text{ext}}(r') } \,d^3r'
\end{equation}
e quindi in generale il momento e la forza valgono
\begin{align}
  \v{ \uptau } = - \frac{\partial U }{\partial \theta } && \v{ F } = - \nabla (U)
\end{align}

\chapter{Magnetismo nella materia}
Fin ora abbiamo trattato i fenomeni magnetici nel vuoto, in questo capitolo studieremo come i materiali reagiscono e generano campi magnetici.

\section{Magnetizzazione}
Sperimentalmente si osserva che alcuni materiali subiscono una forza quando sottoposti a un campo magnetico, infatti attaccando un dinamometro a un 
materiale che vogliamo studiare e sottoponendolo a un campo magnetico generato da un solenoide (noto), questo respinto o attratto. Nel seguito utilizzeremo 
sempre l'approssimazione di dipolo magnetico. Abbiamo visto che l'energia associata al campo magnetico in questa approssimazione si può esprimere come
\begin{equation}
  U = - \v{ m } \cdot \v{ B } \implies F_z = - \left( \nabla U \right)_z = m_z \der{B_z}{z}
\end{equation}
dove abbiamo supposto che il dipolo e il campo siano allineati lungo l'asse $ z $. Possiamo definire la forza per unità di volume
\begin{equation}
  \v{ F }_{\uptau} = \frac{\v{ F }}{\uptau} = \frac{m'_z}{\uptau} \frac{\partial B'_z }{\partial z } \coloneqq M_z \frac{\partial B'_z }{\partial z }
\end{equation}
dove $ M_z $ rappresenta un coefficiente di una grandezza che chiameremo \textbf{magnetizzazione} definita come
\begin{equation}
  \v{ M } = \frac{\v{ m }}{\uptau}
\end{equation}
Per diversi valori di $ M $ possiamo classificare i materiali come:
\begin{itemize}
  \item diamagnetici,  dove la forza è repulsiva;
  \item paramagnetici, dove la forza è debolmente attrattiva;
  \item ferromagnetici, dove la forza è fortemente attrattiva.
\end{itemize}
Passando a un volume infinitesimo otteniamo
\begin{equation}
  \v{ M } = \frac{d \v{ m }}{d \uptau}
\end{equation}
che è una relazione locale la quale lega il momento di dipolo magnetico con $ \v{ M } $.

\section{Permeabilità e suscettibilità magnetica}
Prendiamo un solenoide di $ n $ spire di cui conosciamo il campo nel vuoto e lo riempiamo con il materiale di cui vogliamo conoscerne le
proprietà magnetiche. A livello sperimentale questo procedimento si effettua tramite una \textit{sonda Hall}.
Il rapporto tra il campo (il suo modulo) misurato con il materiale all'interno e il campo nel vuoto viene chiamato permeabilità magnetica relativa:
\begin{equation}
  \mu_r = \frac{B}{B_0}
\end{equation}
Conoscendo la corrente all'interno del solenoide allora il suo campo nel vuoto vale $ B_0 = \mu_0 \, n \, I $ e quindi
\begin{equation}
  B = \mu_r \mu_0 \, n \, I := \mu \, n \,I
\end{equation}
dove $ \mu $ è la permeabilità magnetica del mezzo. Possiamo quindi definire la suscettività magnetica $ \chi_m $ come
\begin{equation}
  \chi_m = \frac{B - B_0}{B_0} = \frac{\left( \mu_r - 1 \right) B_0}{B_0} = \mu_r -1
\end{equation}
Dalla relazione precedente si vede che il campo magnetico totale è dovuto al campo nel vuoto più un termine dovuto al materiale in cui è immerso:
\begin{align}
  \v{ B } &= \v{ B }_0 + \chi_m \v{ B }_0
  \\
  &= \mu_0 \,n \,I + \mu_0 \,\chi_m \, n \,I
\end{align}
dove il termine $ \chi_m \, n \, i $ è dovuto alle \textit{correnti Amperiane}, se $ \chi_m < 0 $ è il caso di diamagnetismo, mentre se
$ \chi_m > 0 $ è il caso del paramagnetismo e ferromagnetismo

\subsection{Leggi di Curie}
Nei mezzi paramagnetici vale la relazione, detta \textit{1° legge di Curie},
\begin{equation}
  \chi_m = \frac{C \rho}{T}
\end{equation}
dove $ T $ è la temperatura, $ \rho $ la resistività e $ C $ una costante. Il fatto che la suscettività magnetica sia inversamente proporzionale alla 
temperatura ci dice 
I mezzi ferromagnetici mantengono le loro proprietà fino a una certa temperatura (dipendente dal mezzo) oltre la quale diventano paramagnetici. Un modello 
che lo descrive è la \textit{2° legge di Curie}:
\begin{align}
  \chi_m = \frac{C \rho}{T - T_c}
\end{align}
dove $ T_C $ è detta temperatura di Curie.
\section{Correnti amperiane}
Nel modello di Bohr, ogni atomo genera una corrente grazie all'orbita degli elettroni attorno al nucleo. Modellizzare questo fenomeno è molto complicato 
e infatti ci limitiamo all'approssimazione di dipolo, che comunque porta a dei risultati più che sufficienti.
\begin{figure}
    \center
    \begin{tikzpicture}[ % FISSA IL RIFERIMENTO IN MODO CORRETTO, ALTRIMENTRI --> (X,Z,Y)
            x={(0.866cm,-0.5cm)},   % cos30, -sin30
            y={(0.866cm,0.5cm)},    % cos30,  sin30
            z={(0cm,0.8cm)},
            scale = 1.,
            name path global/.style={},
            ]
        % Parametri
        \def \xM {2}
        \def \yM {2}
        \def \R {1.5}
        \def \h {4}
        \coordinate (M) at ({\xM},{\yM},0); % centro della base inferiore
        \coordinate (N) at ({\xM},{\yM},\h); % centro della base superiore
        
        % Assi cartesiani
        \draw[->] (0,0,0) -- ({\xM + 1},0,0) node[below left] {$x$};
        \draw[->] (0,0,0) -- (0,{\yM + 1},0) node[above left] {$y$}; % below/above
        \draw[->] (0,0,0) -- (0,0,{\h + 1}) node[left] {$z$};

        % basi
        % base inferiore come path nominato
        \path[name path=BaseCircle] (M) circle [radius={\R}];
        \draw[name path=BaseCircleDraw] (M) circle [radius={\R}];
        \draw (N) circle [radius={\R}];
        % altezze

        % punti estremi sinistra/destra della base inferiore
        %\coordinate (LeftBase)  at ({\xM - \R},{\yM},0);
        %\coordinate (RightBase) at ({\xM + \R},{\yM},0);

        % generatrici verticali dai punti estremi
        %\draw (LeftBase)  -- ++(0,0,\h); % il ++ aggiunge il vettore al vettore precedente (coordinate relative)
        %\draw (RightBase) -- ++(0,0,\h);

        % retta nel piano xy (esempio)
        \path[name path=MyLine] (0,0,0) -- (4,4,0);
        %\draw[thick] (0,1,0) -- (5,3,0);

        % calcolo intersezioni

        \path[name intersections={of=BaseCircle and MyLine, by={I1,I2}}];
        % Salva I1 e I2 come coordinate vere
        \coordinate (C1) at (I1);
        \coordinate (C2) at (I2);

        \draw (I1) -- ++(0,0,\h);
        \draw (I2) -- ++(0,0,\h);
        % punti di intersezione evidenziati
        %\fill (I1) circle (2pt);
        %\fill (I2) circle (2pt);

        \draw[->] (N) -- ++(0,0,1) node[right] {$\mathbf{ M }$};


        % Calcola centro e raggio della semicirconferenza
        \coordinate (CenterArc) at ($(N)+(0,0,1)$);
        % Calcola i punti di partenza e arrivo
        \coordinate (StartArc) at ($(C2)+(0,0,1)$);
        \coordinate (EndArc) at ($(C1)+(0,0,1)$);

        % Calcola angoli rispetto al centro
        \pgfmathanglebetweenpoints{\pgfpointanchor{CenterArc}{center}}{\pgfpointanchor{StartArc}{center}}
        \let\StartAngle\pgfmathresult
        \pgfmathanglebetweenpoints{\pgfpointanchor{CenterArc}{center}}{\pgfpointanchor{EndArc}{center}}
        \let\EndAngle\pgfmathresult
        % Calcola raggio
        \pgfmathsetmacro{\ArcRadius}{veclen((\xM+\R)-(\xM), (\yM)-(\yM))}

        % Disegna la semicirconferenza
        \draw[thick,->] (StartArc) arc[start angle=\StartAngle, end angle=405, radius=\ArcRadius] node[above right] {$ \mathbf{ K }_m $};


    \end{tikzpicture}
    \caption{Cilindro con vettore magnetizzazione}
    \label{figure:_cilindro_con_magnetizzazione}
\end{figure} % Cilindro con momenti di dipolo -> magnetizzazione
Consideriamo un piccolo volume $ \uptau $, allora la magnetizzazione $ \v{ M } $ dovuta ai momenti di dipolo degli atomi del volume è
\begin{equation}
  \v{ M } = \frac{N_{\uptau} \, \langle \v{ m } \rangle}{\uptau} = n \, \langle \v{ m } \rangle
\end{equation}
con $ N_{\uptau} $ il numero di atomi nel volume $ \uptau $ e $ n $ la densità di atomi. Per facilità di calcolo consideriamo un cilindro 
(figura \ref{figure:_cilindro_con_magnetizzazione}), sezionandolo in modo perpendicolare all'asse troviamo che
\begin{equation}
  d\v{ m } = \v{ M } d\uptau = \v{ M } d\Sigma dz
\end{equation}
ma ricordando che il momento di dipolo magnetico è associato a una corrente dalla relazione $ d\v{ m } = dI_m d\Sigma \vers{z} $ si ottiene
\begin{align}
  dI_m = M dz
\end{align}
Integriamo lungo l'asse del cilindro
\begin{align}
  I_m = \int_0^h M \,dz = M h \implies M = \frac{I_m}{h} = K_m
\end{align}
e da questa relazione si nota che la magnetizzazione è dovuta alle correnti superficiali del cilindro
\begin{equation}
  \v{ K }_m = \v{ M } \times \vers{n}
\end{equation}
proprio come la densità di carica di polarizzazione $ \sigma_P = \v{ P \cdot n } $


Calcolando la circuitazione di $ \v{ M } $ lungo un circuito chiuso $ C $ che passa dall'asse del cilindro otteniamo:
\begin{equation}
  \oint_C \v{ M } \, d \v{ s } = M \,h = K_m \, h = I_m = \int_{\Sigma(C)} \v{ J }_m \,d\v{ \Sigma }
\end{equation}
posso applicare il teorema di Stokes alla circuitazione della magnetizzazione e si ha che
\begin{equation}
  \int_{\Sigma(C)} \v{\nabla}\times\v{ M } \,d\v{ \Sigma } = \int_{\Sigma(C)} \v{ J }_m \,d\v{ \Sigma }
  \implies \v{\nabla}\times\v{ M } = \v{ J }_m
\end{equation}
che questa volta è l'analogo di $ \displaystyle \nabla\!\cdot\!\v{ P } = - \rho_P $, e se $ \v{ M } $ è uniforme allora $ \v{ J }_m = 0 $. 
Riportiamo quindi i risultati ottenuti:
\begin{align}
  \v{ K }_m = \v{ M } \times \vers{n} \qquad \v{ J }_m = \nabla\!\times\!\v{ M }
\end{align}
che valgono in generale per qualsiasi volume e rappresentano le manifestazioni macroscopiche delle proprietà atomiche del materiale considerato.

Applichiamo le relazioni appena ottenute nel caso di un solenoide con all'interno un mezzo (cilindrico): calcoliamo la circuitazione di $ \v{ B } $
\begin{equation}
  \oint_C \v{ B }  \,d\v{ s } = \mu_0(I + I_m) = \mu_0 I + \mu_0 \oint_C \v{ M } \, d\v{ s }
\end{equation}
passando alla versione locale
\begin{align}
  &\v{\nabla}\times\v{ B } = \mu_0\left( \v{ J + J }_m \right) = \mu_0\v{ J } + \mu_0\v{\nabla}\times\v{ M }
  \\
  &\implies \nabla \times \left( \frac{\v{ B }}{\mu_0}-\v{ M } \right) = \v{ J } 
\end{align}
dove l'ultima $ \v{ J } $ è di conduzione. Possiamo definire il campo $ \v{ H } $ detto \textit{campo magnetizzante o campo ausiliario}\!
\footnote{In alcuni testi il vettore $ \v{ H } $ è indicato come campo magnetico, mentre $ \v{ B } $ è chiamato vettore di induzione magnetica.}: 
\begin{align}
  \v{ H } \coloneqq \frac{\v{ B }}{\mu_0} - \v{ M }
\end{align}
cioè è un campo che tiene conto sia del campo magnetico esterno sia degli effetti microscopici del materiale. In questo modo le equazioni di Maxwell 
statiche si riscrivono nella forma
\begin{align}
  &&\nabla\!\cdot\!\v{ D } &= \rho_L & \nabla\!\times\!\v{ E } &= 0&&
  \\
  &&\nabla\!\cdot\!\v{ B } &= 0 & \nabla\!\times\!\v{ H } &= \v{ J }&&
\end{align}
Grazie al vettore $ \v{ H } $ la legge di Ampère in forma globale diventa
\begin{align}
  \oint_\gamma \v{ H } \cdot d\v{ l } = I_{\text{conc}}
\end{align}

\subsection{Mezzi omogenei e isotropi}
In seguito tratteremo soltanto mezzi materiali detti omogenei e isotropi, cioè quelli per cui in ogni punto del materiale (isotropia) il campo 
$ \v{ H } $ risulta parallelo al campo magnetico esterno $ \v{ B } $ (omogeneità), cioè vale
\begin{align}
  \v{ B } = \mu \v{ H } \implies \v{ M } = \chi_m \v{ H } = \frac{\chi_m}{\mu} \v{ B }
\end{align}
in tutto il mezzo. Altri modelli per i mezzi magnetici prevedono risposte quadratiche oltre al termine lineare oppure orientazioni diverse in base alla 
geometria.

\section{Interfaccia fra materiali magnetici}
Avevamo già visto come la discontinuità di campo dovuta una corrente superficiale si poteva esprimere come
\begin{equation}
  \v{ B }_2 - \v{ B }_1 = \mu_0 \v{ K } \times \v{ n }
\end{equation}
Vogliamo vedere adesso come si comportano $ \v{ B } $ e $ \v{ H } $ fra due materiali diversi a contatto. Consideriamo il caso in cui i 
campi sono sullo stesso piano, prendiamo un cilindro e ne calcoliamo il flusso di $ \v{ B } $, e otteniamo
\begin{align}
  B_1 \cos \theta_1 &= B_2 \cos \theta_2 \implies B_{1|} = B_{2|}
  \\
  \mu_1 H_1 \cos \theta_1 &= \mu_2 H_2 \cos \theta_2 \implies \mu_1 H_{1|} = \mu_2 H_{2|}
\end{align}
Per la discontinuità della componente parallela prendiamo un circuito amperiano e calcoliamo la circuitazione ottenendo
\begin{align}
  B_{2||} - B_{1||} = K
\end{align}
\dots

\subsection{Elettromagneti}
Consideriamo il solito circuito di un ferromagnete con una forza magnetomotrice, interrotto per una lunghezza $ h $. Nell'intercapedine 
$ B = \mu H = \mu_0 H_0$, mentre la circuitazione di $ \v{ H } $ vale 
\begin{equation}
  \oint H \,ds = H (S-h) + H_0 h = Ni
\end{equation}
e sfruttando la relazione di prima fra $ H_0 $ e $ B $ si ottiene 
\begin{equation}
  H (S-h) + \frac{B h }{\mu_0} = Ni \implies B = \frac{\mu_0 N i }{h} - \frac{S-h}{h} H
\end{equation}
e tracciando la curva di isteresi osserviamo che

\chapter{Esercizi della magnetostatica}
\section{Sfera carica rotante}
Consideriamo un guscio sferico con densità di carica superficiale $ \sigma $ e raggio $ R $, che ruota con velocità angolare 
costante $ \overset{\to}{\omega} = \omega \v{ u }_z $. Si vuole trovare il campo magnetico in tutto lo spazio (dentro e fuori).\\
È possibile risolverlo con due approcci: 
\begin{itemize}
  \item il primo è il più formale e consiste nel sfruttare il potenziale magnetico scalare, risolvendo l'equazione di 
  Laplace e imponendo le condizioni al bordo dovute alla densità superficiale di carica
  \item il secondo è più fisico e consiste nell'osservare che la sfera ruotando è come se fossero infinite spire di spessore $ Rd\theta $ ognuna con 
  una corrente $ \Delta i $. Integriamo tutti i dipoli magnetici $ d\v{ m } = m \v{ u }_z $ e poi osserviamo che
\end{itemize}
Intanto calcoliamo la densità superficiale
\begin{equation}
  \v{ K } = \sigma \v{ v } = \sigma \v{ \omega } \times \v{ r } = \sigma \omega R \sin \theta := K_0 \sin \theta
\end{equation}
Primo approccio:
\begin{equation}
  \nabla^2 W = 0 \implies W = \sum_{l=0}^{\infty} \left( A_l \, r^l + \frac{B_l}{r^{l+1}} \right) P_l\left( \cos \theta \right) 
\end{equation}
\[
\begin{cases}
  W_> = \sum_{l=0}^{\infty}  \frac{B_l}{r^{l+1}} P_l\left( \cos \theta \right) \\
  \\
  W_< = \sum_{l=0}^{\infty}  A_l \, r^l  P_l\left( \cos \theta \right) 
\end{cases}
\]
Le condizioni che dobbiamo imporre sono date dalle leggi di Maxwell, cioè che la componente di $ \v{ B } $ ortogonale al guscio sferico sia continua 
(in modo che il flusso su una qualsiasi superficie chiusa sia $ 0 $) e che la componente tangenziale abbia una discontinuità pari a $ \mu_0 \v{ K } $.
In formule:
\begin{align}
  &(1) \quad \v{ n } \cdot \left( \v{ B_> - B_< } \right) = 0\\
  &(2) \quad \v{ n } \times \left( \v{ B_> - B_< } \right) = - \mu_0 \v{ K }
\end{align}
da notare il segno $ - $ nella seconda perché per il verso della circuitazione ($ \v{ -u }_{\varphi} $) è opposto a quello di $ \v{ K } $ 
($ \v{ u }_{\varphi} $). Sviluppando si ottiene
\begin{align}
  (1)\;\left. \frac{\partial W_> }{\partial r } \right|_{r=R} - \left. \frac{\partial W_< }{\partial r }\right|_{r=R} &= \left( -(l+1)\frac{B_l}{R^{l+2}} -lA_lR^{l-1} \right) P_l(\cos\theta) = 0 \\
   A_l l &= -(l+1) \frac{B_l}{R^{2l+1}}
\end{align}
\begin{align}
  &(2)\; \left. \left[ \frac{1}{r} \frac{\partial W_> }{\partial \theta } - \frac{1}{r} \frac{\partial W_- }{\partial \theta } \right]  \right|_{r=R} 
  = \left( \frac{B_l}{R^{l+2}} - A_l R^{l-1} \right)\frac{\partial P_l(\cos\theta) }{\partial \theta } = -\mu_0K_0\sin\theta\\
  &\text{ma} \quad -\sin\theta = \frac{\partial \cos\theta }{\partial \theta } = \frac{\partial P_1(\cos\theta) }{\partial \theta } \implies A_l = B_l = 0 \; \forall l\not= 1
\end{align}
e combinandole otteniamo
\begin{align}
    \begin{cases}
        A_1 +2 \frac{B_1}{R^{3}} = 0 \\
        \frac{B_1}{R^{3}} - A_1 = \mu_0K_0
    \end{cases}
    &\implies
    A_1 = - \frac{2\mu_0K_0}{3}  \, , \quad B_1 = \frac{\mu_0 K_0 R^3}{3} 
\end{align}
Il potenziale è quindi
\begin{align}
  W_< &= - \frac{2\mu_0K_0}{3} r \cos\theta = - \frac{2\mu_0K_0}{3} \,z\\
  W_> &= \frac{\mu_0 K_0 R^3}{3} \frac{\cos\theta }{r^2} = \frac{\mu_0}{4\pi} \frac{\v{ m } \cdot \v{ r }}{r^3}
\end{align}
dove nell'ultima uguaglianza abbiamo sostituito $ \frac{4\pi}{3}K_0R^3 = m $, che come calcoleremo nel secondo approccio, è il momento di dipolo del 
guscio sferico carico rotante. Questo ci permette di riconoscere un campo di dipolo all'esterno e un campo uniforme all'interno:
\begin{align}
  \v{ B }_< &=  \frac{2\mu_0K_0}{3} \left( \frac{\partial (r \cos\theta) }{\partial r } \v{ u }_r + \frac{1}{r} \frac{\partial (r\cos\theta) }{\partial \theta } \v{ u }_{\theta}\right) \\
  &= \frac{2\mu_0K_0}{3} \left( \cos\theta \v{ u }_r - \sin\theta \v{ u }_{\theta} \right) = \frac{2\mu_0K_0}{3} \v{ u }_z
\end{align}
\begin{align}
  \v{ B }_> = \frac{\mu_0}{4\pi} \frac{\left( \v{ r (r \cdot m) } - \v{ m } r^2 \right)}{r^5} = \frac{\mu_0}{4\pi} \frac{\left( \v{ u }_r (\v{ u }_r \cdot \v{ m }) - \v{ m } \right)}{r^3}
\end{align}

Secondo approccio:\\
Sia $ d $ un ipotetico spessore, la corrente in funzione di $ \v{ K } $ è
\begin{align}
  &\Delta i = \v{ J } \cdot \v{ S } = \frac{K}{d} \; d \; R d\theta = K_0 R \sin \theta \,d\theta\\
  &dm =  \pi (R \sin \theta )^2 \Delta i
\end{align}
integrando su $ \theta $ i momenti magnetici si ottiene
\begin{align}
  m &= \int \pi (R \sin \theta )^2 di = \int_{\pi}^{0} \pi (R \sin \theta )^2 K_0 R \sin \theta \,d\theta \\
  &= K_0 R^3 \pi \int_{\pi}^{0} \left( \sin\theta \right)^3 \,d\theta = K_0 R^3 \pi \int_{\pi}^{0} \left( 1 - \cos^2 \theta \right)  \,d\left( -\cos\theta \right) \\
  &= K_0 R^3 \pi \left[ \cos\theta - \frac{\cos^3 \theta}{3} \right]_{0}^{\pi} = \frac{4}{3} K_0 R^3 \pi = \frac{4}{3} \pi \sigma \omega R^4 = \frac{Q \omega R^2}{3}
\end{align}

SIMMETRIE

Concludiamo osservando che il sistema genera un campo uniforme all'interno e un campo di dipolo all'esterno ma senza alcuna approssimazione: i risultati 
ottenuti sono risultati esatti. Inoltre, a seguito delle considerazioni sulle simmetrie fatte nel secondo approccio, questo sistema è il sistema più 
elementare di dipolo magnetico esatto in una regione di illimitata. 

\section{Cilindri con correnti opposte}
Siano due cilindri coassiali di raggi $ a < b $, altezza $ h \gg b $ collegati a un estremo da una \textit{ddp} e all'altro da una resistenza $ R $. 
Trovare il campo magnetico e il coefficiente di autoinduttanza.

\textit{Soluzione} Il campo esterno è nullo perché la corrente concatena totale è $ i - i = 0 $, per lo stesso motivo è nullo anche il campo 
nella regione $ \rho < a $

\chapter{Elettrodinamica}
Continuiamo la trattazione dell'elettromagnetismo passando al caso dinamico. Nei precedenti capitoli abbiamo discusso i campi elettrico e magnetico 
e le relative forze nel caso in cui l'evoluzione temporale non influenzi l'evoluzione del sistema, e questo ci ha portato alle equazioni di Maxwell 
(nel vuoto) statiche:
\begin{align}
  &\nabla\cdot\v{ E } = \frac{\rho}{\varepsilon_0} &
  &\v{\nabla}\times\v{ E } = 0 \\
  &\nabla\cdot\v{ B } = 0 &
  &\v{\nabla}\times\v{ B } = \mu_0 \v{ J }
\end{align}
Le equazioni così come sono non vanno chiaramente bene dato che l'equazione di continuità non è rispettata nel caso non banale e da un punto di vista 
sperimentale si osserva che in alcuni casi il campo elettrico non è conservativo. Tuttavia le leggi di Gauss per campo elettrico e magnetico rimangono 
invariate.

\section{Rotore del campo elettrico}
Per passare al caso dinamico dobbiamo tenere presente che l'equazione di continuità \eqref{eq:continuità} non ci fornisce più delle correnti 
stazionarie, ma in generale la quantità di cariche all'interno di un volume può variare nel tempo. Il fatto che $ \nabla\cdot\v{ J } \neq 0 $
fa perdere di validità la legge di Ampère così come l'abbiamo formulata.\\
Analizziamo adesso la seconda legge di Maxwell per il 
campo elettrico: Faraday osserva che in un circuito sottoposto a un campo magnetico, effettuando delle variazioni al campo, si misurava una 
differenza di potenziale agli estremi del circuito. La legge di Faraday è espressa come
\begin{equation}
  \mathscr{E} = - \frac{d \Phi (\v{ B })}{dt}
\end{equation}
dove $ \mathscr{E} $ è la differenza di potenziale e il flusso è sulla superficie che ha come bordo il circuito $ \gamma $ (che abbiamo mostrato 
essere equivalente per qualsiasi superficie che abbia $ \gamma $ come bordo). Vogliamo mettere in relazione questa legge sperimentale con la 
definizione di differenza di potenziale ai capi di un circuito che abbiamo dato:
\begin{align}
  \mathscr{E} &= \oint_{\gamma} \frac{\v{ F }}{q} \, d\v{ l } = \oint \left( \v{ E } +  \v{ v } \times \v{ B } \right) \, d\v{ l } 
  = \oint_{\gamma} \v{ E } \, d\v{ l } + \oint_{\gamma} \v{ v } \times \v{ B } \, d\v{ l }
\end{align}
il secondo pezzo dell'integrale rappresenta fisica "vecchia", che andremo a sviluppare nel capo il circuito si deformi nel tempo, mentre il primo 
termine qualora sia non nullo rappresenta "nuova" fisica.\\
Per sviluppare il secondo termine osserviamo che la velocità di una carica possiamo scomporla come velocità di traslazione del circuito e 
velocità relativa al circuito
\begin{equation}
  \v{ v } = \v{ V }_{\text{T}} + \v{ v }_{\text{rel}} \quad \text{con} \quad \v{ v }_{\text{rel}} \parallel d\v{ l }
\end{equation}
quindi il fatto che $ \v{ v }_{\text{rel}} $ sia parallelo al tratto del circuito il suo integrale è nullo e rimane
\begin{align}
  \oint_{\gamma} \v{ v } \times \v{ B } \, d\v{ l } = \oint_{\gamma} \v{ V }_{\text{T}} \times \v{ B } \, d\v{ l }
\end{align}

\textbf{Caso $ V_{\text{T}} = 0$:} abbiamo solo il termine con la circuitazione di $ \v{ E } $ quindi
\begin{align}
  \mathscr{E} &= - \frac{d \Phi(\v{ B })}{dt} = -\frac{d}{dt} \left( \int_{S_{\gamma}} \v{ B }\! \cdot\! d\v{ S } \right) 
  = - \int_{S_{\gamma}} \frac{\partial \v{ B } }{\partial t } d\v{ S }
  \\ &= \oint_{\gamma} \v{ E } \, d\v{ l } \overset{\text{T. Rotore}}{=}
   \int_{S_{\gamma}} \left( \v{\nabla}\times\v{ E } \right) \cdot d\v{ S }
\end{align}
dove abbiamo sfruttato che nella derivata totale rispetto al tempo le coordinate sono costanti, quindi diventa una derivata parziale, ed eguagliando gli 
integrandi otteniamo
\begin{equation} \label{eq:rotore_E_der_B}
  \v{\nabla}\times\v{ E } = - \frac{\partial \v{ B } }{\partial t }
\end{equation}

\textbf{Caso generale:} vogliamo vedere se è possibile ottenere lo stesso risultato anche per un circuito che evolve in modo generico nel tempo.
Sia $ \gamma (t) $ il circuito dipendente dal tempo, sviluppiamo la variazione del flusso di $ \v{ B } $:
\begin{align}
  \Delta \Phi (\v{ B }) = \Phi ()
\end{align}
\dots

\section{Induzione elettromagnetica}
Per un sistema di conduttori è possibile definire i coefficienti di capacità, cioè dei valori numerici che dipendono dalla geometria del 
conduttore, utili a determinare il potenziale del conduttore $ i $-esimo in funzione della carica presente sugli altri conduttori. Siano i $ C_{ij} $ tali 
coefficienti, allora
\begin{align}
  V_i = \sum_{j} C_{ij} Q_j
\end{align}
dove $ V_i $ è il potenziale del conduttore $ i $-esimo. È possibile intenderli come il coefficiente che determina il potenziale indotto da un conduttore 
su un altro, ed è possibile dimostrare che sono simmetrici cioè $ C_{ij} = C_{ji} $.\\
Lo scopo di questa sezione non è trattare i coefficienti di capacità ma sfruttare lo stesso ragionamento per definire i coefficienti di \textit{autoinduzione} 
e di \textit{mutua induzione}. Consideriamo $ N $ circuiti distinti capaci di condurre, dall'espressione dell'energia del campo elettromagnetico 
(che viene fornita nella sezione \ref{section:_Onde_elettromagnetiche}) è possibile arrivare alla seguente forma:
\begin{align}
  U = \frac{1}{2} \sum_{i}^N L_i I_i^2 + \sum_i^N \sum_{j>i}^N M_{ij}I_i I_j
\end{align}
dove gli $ L_i $ sono i coefficienti di autoinduzione e gli $ M_{ij} $ sono i coefficienti di mutua induzione. In condizioni quasi stazionarie il flusso 
del campo magnetico è proporzionale al campo che a sua volta è proporzionale alla corrente, tale costante è detta \textit{induttanza} ed è definita come:
\begin{align}
  L \coloneq \frac{\Phi(B)}{I}
\end{align}
e in questo modo l'energia immagazzinata in un induttore si può esprimere come
\begin{align}
  U = \frac{1}{2} \Phi(\v{ B }) I = \frac{1}{2} L I^2
\end{align}
I coefficienti di mutua induzione sono particolarmente utili per determinare l'forza elettromotrice indotta su un circuito da un altro, infatti considerando 
noto due circuiti, una volta trovato $ M = M_{12} = M_{21} $ si ha che la \textit{fem} del primo circuito è data da
\begin{align}
  \mathcal{E}_1 = - \frac{d\Phi_{\v{ B }_2}}{dt} = - M \frac{dI_2}{dt}
\end{align}
e la stessa cosa vale per il secondo circuito. Determinare i coefficienti di mutua induzione è sempre possibile svolgendo l'integrale che esprime il flusso 
del campo generato da un circuito attraverso la superficie dell'altro, e come si può intuire i calcoli sono notevolmente più semplici quando si trattano 
solenoidi o spire circolari con particolari simmetrie.

\subsection{Autoinduttanza di un solenoide}
Consideriamo un solenoide di $ N $ spire, lungo $ l $, di raggio $ R\ll l $, densità di spire$ n = \frac{N}{l} $, in cui scorre una corrente $ I(t) $. 
Allora
\begin{align}
  B = \mu_0 n I &\implies \Phi(\v{ B }) = N \pi R^2 B = \mu_0 \pi n N R^2 I = \mu_0 \pi n^2 l R^2 I
  \\
  &\implies L = \mu_0 \pi n^2 l R^2
\end{align}
Se riempiamo il solenoide con un materiale di permeabilità magnetica $ \mu_0 $ abbiamo che
\begin{align}
  H = n I \qquad B = \mu H \implies L_\mu = \mu_r L = \mu \pi l R^2 n^2
\end{align}

\subsection{Circuito RL}
Sfruttiamo l'autoinduzione di un solenoide calcolata nel precedente paragrafo per risolvere (cioè trovare la corrente in funzione del tempo) 
un circuito formato da generatore, solenoide e resistenza come in figura \ref{figure:_circuito_RL}.\begin{figure}
    \center
    \begin{circuitikz}
        \def \x {3}

        \draw
        (0, 0) to[battery1 = $ V $, invert] (0,\x)
        to (\x,\x) 
        to[cute inductor = $ L $] (\x,0)
        to[american resistor= $ R $] (0,0);

    \end{circuitikz}
    \caption{Circuito RL}
    \label{figure:_circuito_RL}
\end{figure}
Scriviamo la legge delle maglie
\begin{align}
  \mathcal{E} + \left( - \frac{d\Phi_{\v{ B }}}{dt} \right) + \left(  -RI \right) &= 0
  \\
  \mathcal{E} - L \dot I - RI &= 0
\end{align}
che è una equazione differenziale lineare del prim'ordine. Le soluzioni generale e particolare sono
\begin{align}
  I_g = I_0 e^{-\frac{R}{L} t}  \qquad I_p = \frac{\mathcal{E}}{R}
\end{align}
la condizione da imporre è $ I(0) = 0 $ quindi la soluzione vale
\begin{align}
  I(t) = \frac{\mathcal{E}}{R} \left( 1-e^{-\frac{R}{L} t} \right)
\end{align}

\section{Rotore del campo magnetico}
Abbiamo visto che con la nuova espressione per il rotore di $ \v{ E } $ le equazioni sono accoppiate fra campo magnetico e campo elettrico. 
Esaminando le equazioni in questo stato (non in questa forma precisa) Maxwell si accorge di una inconsistenza fra queste: se si prova a calcolare la 
divergenza del rotore di $ \v{ B } $ questa dovrebbe essere nulla eppure
\begin{align}
  \nabla\! \cdot\! \left( \nabla\! \times\!\v{ B } \right) = 4\pi k_m \nabla \cdot \v{ J } 
  \overset{\text{eq. continuità}}{=} 4\pi k_m \left(- \der{\rho}{t} \right)
\end{align}
che in generale è non nulla. Notando questa inconsistenza, Maxwell nel 1864 corregge le equazioni mosso soltanto da tali motivazioni teoriche  
e non perché avesse fatto esperimenti che evidenziassero una problema nel modello, anche perché la conferma sperimentale delle equazioni dopo la 
modifica di Maxwell arrivò soltanto nel 1884 (ben 20 anni dopo) dal laboratorio di Hertz.

La modifica consiste nell'imporre che la quarta equazione abbia divergenza nulla, aggiungendo un termine apposito:
\begin{align}
  \v{\nabla}\times\v{ B } = \mu_0 \v{ J } + \v{ \Lambda }
\end{align}
calcoliamone la divergenza:
\begin{align}
  0 = \nabla\! \cdot\! \left( \nabla\! \times\!\v{ B } \right) &= \mu_0 \left( - \der{\rho}{t} \right) + \nabla \cdot \v{ \Lambda } 
  = \mu_0 \left( - \varepsilon_0 \der{\left( \nabla \! \cdot \!\v{ E } \right)}{t} \right) + \nabla\!\cdot\!\v{ \Lambda }
  \\
  &= \nabla \! \cdot \! \left( - \mu_0\varepsilon_0 \der{\v{ E }}{t} + \v{ \Lambda } \right)
\end{align}
e una scelta possibile è
\begin{equation}
  \v{ \Lambda } = \mu_0 \varepsilon_0 \der{\v{ E }}{t}
\end{equation}
È una scelta perché potremmo aggiungere un qualsiasi campo vettoriale a divergenza nulla e otterremmo delle equazioni consistenti dal punto di 
vista matematico, tuttavia questo non è possibile dal punto di vista fisico per delle proprietà che le equazioni devono rispettare che tratteremo 
più avanti. La forma definitiva delle equazioni di Maxwell è
\begin{align} \label{equazioni di Maxwell dinamiche}
  &\nabla\cdot\v{ E } = \frac{\rho}{\varepsilon_0} &
  &\v{\nabla}\times\v{ E } = - \der{\v{ B }}{t} \\
  &\nabla\cdot\v{ B } = 0 &
  &\v{\nabla}\times\v{ B } = \mu_0 \v{ J } + \mu_0 \varepsilon_0 \der{\v{ E }}{t}
\end{align}
In questo modo l'equazione di continuità diventa un corollario della prima e della quarta equazione di Maxwell.
\subsection{Corrente di spostamento nel condensatore piano}
\begin{figure}[ht]
    \centering
    \begin{tikzpicture}[>=Stealth, scale=1.1,
        plate/.style={line width=1pt},
        wire/.style={line width=1pt},
        efield/.style={gray!70, line width=0.6pt, ->}]

        \def\r{1.5}   % raggio verticale delle armature
        \def\h{2}     % distanza tra le armature
        \def\e{0.4}   % schiacciamento orizzontale delle ellissi

        % Campo E nel gap: flusso che attraversa l'estremita' chiusa di S
        \foreach \y in {-0.7, 0, 0.7}{ \draw[efield] (0.30,\y) -- (1.70,\y); }
        \node[gray!70] at (1.5,1.25) {$\vec{E}$};

        % Armatura destra
        \draw[plate] (\h,\r) arc[start angle=90, end angle=270, x radius=\e, y radius=\r];
        \draw[plate, dashed] (\h,\r) arc[start angle=90, end angle=-90, x radius=\e, y radius=\r];

        % Fili
        \draw[wire, dashed] (0,0) -- (-\e,0);          % tratto dietro l'armatura
        \draw[wire] (-\e,0) -- (-4.6,0);               % filo che porta la corrente I
        \draw[wire] (\h,0) -- (\h+1.2,0);              % terminale destro
        \draw[wire, -{Stealth[length=3mm]}] (-4.1,0) -- (-3.3,0);
        \node at (-3.7,0.30) {$I$};

        % Superficie S, profilo ad anfora:
        %  - tangente orizzontale in partenza da C
        %  - si gonfia sopra l'armatura sinistra (non la interseca)
        %  - si richiude a tangente verticale fra le armature
        \path[fill=blue!50, fill opacity=0.25, draw=blue!55!black, line width=0.6pt]
            (-2.5,1.0)
            .. controls (-1.9, 1.00) and (-1.2, 1.95) .. (-0.1, 1.95)
            .. controls ( 0.7, 1.95) and ( 1.1, 1.05) .. ( 1.1, 0.00)
            .. controls ( 1.1,-1.05) and ( 0.7,-1.95) .. (-0.1,-1.95)
            .. controls (-1.2,-1.95) and (-1.9,-1.00) .. (-2.5,-1.0)
            arc[start angle=-90, end angle=90, x radius=0.45, y radius=1.0]
            -- cycle;

        % Armatura sinistra (davanti al sacco)
        \draw[plate] (0,\r) arc[start angle=90, end angle=270, x radius=\e, y radius=\r];
        \draw[plate, dashed] (0,\r) arc[start angle=90, end angle=-90, x radius=\e, y radius=\r];

        \node at (-0.62,0.45) {$+$};
        \node at (\h+0.62,0.45) {$-$};

        % Bordo C (loop attorno al filo) e verso di circuitazione
        \draw[fill=red, fill opacity = 0.25, red!75!black, line width=1pt] (-2.5,0) ellipse[x radius=0.45, y radius=1.0];
        \draw[red!75!black, -{Stealth[length=2.4mm]}, line width=1pt]
            (-2.105,-0.5) arc[start angle=-30, end angle=35, x radius=0.45, y radius=1.0];

        % Etichette
        \node[red!75!black] at (-2.5,1.35) {$C$};
        \node[blue!55!black] at (-0.30,1.55) {$S$};
        \node[red!75!black] at (-2.5,0.35) {$\Sigma$};
        \draw[gray!60,line width=0.5pt] (1.13,-0.12) -- (1.7,-1.65);
        \node at (2.15,-1.9) {\footnotesize "corrente di spostamento"};

    \end{tikzpicture}
    \caption{Condensatore piano con una curva $ C $ concatenata al filo e due superfici $ \Sigma $ e $ S $ con bordo la stessa $ C $.}
    \label{figure: condensatore corrente spostamento}
\end{figure}
Un esempio della consistenza delle equazioni col termine della corrente di spostamento può essere visto analizzando il caso del 
condensatore piano a facce parallele in figura \ref{figure: condensatore corrente spostamento}. Consideriamo il condensatore isolato e le facce collegate 
entrambe a un filo conduttore nel quale scorre la corrente 
$ I $: se prendiamo la superficie circolare $ \Sigma $ attraversata dal filo con bordo $ C $ e calcoliamo la circuitazione del campo magnetico otteniamo
\begin{equation}
  \oint_{C} \v{ B } \, d\v{ l } = 2\pi r B_{\varphi} = 4\pi k_m I
\end{equation}
Ma se prendiamo la superficie $ S $ che ha sempre $ C $ come bordo ma passa in mezzo alle armature in modo che non ci sia nessuna 
corrente concatenata, otteniamo che la stessa circuitazione del campo magnetico è nulla, che è assurdo. Considerando anche il 
termine di Maxwell otteniamo che
\begin{equation}
  \oint_{C} \v{ B } \, d\v{ l } = \mu_0 \Phi_S(\v{ J }) + \mu_0 \varepsilon_0 \Phi_S(\v{ E }) 
  = 0 + \mu_0 \varepsilon_0 \der{}{t} \left( \frac{Q}{\varepsilon_0} \right) = \mu_0 I
\end{equation}
dove per calcolare il flusso di $ \v{ E } $ abbiamo sfruttato che il campo è non nullo solo fra le armature quindi possiamo 
applicare il teorema di Gauss alla superficie chiusa $ S \cup \Sigma $. 
Da questo esempio deriva il termine \textit{corrente di spostamento} proprio in analogia con la legge di Biot-Savart, tuttavia ribadiamo che il termine di 
Maxwell non rappresenta in alcun modo una corrente.

\section{Campi lentamente variabili} \label{section: campi lentamente variabili}
Studiare i regimi in cui i campi sono considerati lentamente variabili è utile perché il termine della corrente di spostamento risulta trascurabile, 
inoltre da un punto di vista storico ci sono voluti molti anni per riuscire a realizzare dei campi rapidamente variabili che richiedono tecnologie 
diverse. In questa sezione vogliamo capire quali sono le condizioni per stabilire se un campo è lentamente variabile.

\begin{figure}[h]
    \centering
    \begin{tikzpicture}[scale=1]

        % Parametri
        \def\r{1.5}      % raggio orizzontale ellisse
        \def\h{2}        % distanza tra le armature
        \def\e{0.4}      % schiacciamento verticale ellisse

        % Armatura superiore
        \draw (-\r,\h) arc[start angle=180, end angle=360, x radius=\r, y radius=\e];
        \draw[dashed] (-\r,\h) arc[start angle=180, end angle=0, x radius=\r, y radius=\e];

        % Armatura inferiore
        \draw (-\r,0) arc[start angle=180, end angle=360, x radius=\r, y radius=\e];
        \draw[dashed] (-\r,0) arc[start angle=180, end angle=0, x radius=\r, y radius=\e];

        % Collegamenti laterali
        %\draw (-\r,0) -- (-\r,\h);
        %\draw (\r,0) -- (\r,\h);

        % Fili verticali
        \draw (0,\h) -- (0,\h+1);
        \draw [dashed] (0,0) -- (0,-\e);
        \draw (0,-\e) -- (0,-1);

        % Campo elettrico (frecce interne)
        %\foreach \x in {-1, -0.5, 0, 0.5, 1}{  \draw[->, thick] (\x,0.4) -- (\x,\h-0.4);}

        % Etichette
        \node at (-0.5,\h+0.7) {$-$};
        \node at (-0.5,-0.7) {$+$};

    \end{tikzpicture}
    \caption{Condensatore piano visto leggermente dall'alto}
    \label{figure: condensatore verticale}
\end{figure}

Consideriamo un condensatore a facce piane parallele circolari come in figura \ref{figure: condensatore verticale} dove la carica su una armatura è 
data da $ Q(t) = Q_0 \sin(\omega t) $, e quindi il campo è uniforme e vale  
\begin{equation}
  \v{ E } = \frac{Q}{\varepsilon A} \vers{z} = E_0 \sin(\omega t) \vers{z}.
\end{equation}
Vedremo come questo termine è il primo di 
uno sviluppo in potenze di $ \omega $. Ne calcoliamo la derivata rispetto al tempo
\begin{equation}
  \der{\v{ E }}{t} = E_0 \omega \cos(\omega t) \vers{z}
\end{equation}
Calcoliamo ora la circuitazione del campo magnetico lungo una circonferenza $ \gamma $ di raggio $ r $ concentrica al condensatore
\begin{align}
  \oint_{\gamma} \v{ B }  \,d\v{ l } = 2\pi r B_{\varphi} = \mu \varepsilon \int_{S_{\gamma}} \der{\v{ E }}{t} \,d\v{ S } 
  = \mu \varepsilon E_0 \omega \cos(\omega t) \pi r^2
\end{align}
quindi il campo vale
\begin{align}
  B_{\varphi}(r) = \mu \varepsilon  r \frac{\omega}{2} E_0\cos(\omega t)
\end{align}
quindi,
\begin{equation}
  \v{ B }_1(r) = \frac{\mu \varepsilon \omega r}{2} E_0 \cos(\omega t) \vers{\varphi}
\end{equation}
Siccome $ \v{ B }_1 $ dipende dal tempo possiamo usare nuovamente la 3° equazione. 
Prendiamo l'asse del condensatore e calcoliamo la circuitazione del campo elettrico lungo un piccolo circuito $ \gamma $:
\begin{align}
  \oint_{\gamma} \v{ E } \, d\v{ l } &= h E_z(r= 0) - h E_z (r \neq 0)
  = -\der{}{t} \int_{S_\gamma} \v{ B } \, d\v{ S } 
  \\
  &= - \der{}{t} \int_0^h dz \int_0^r \frac{\mu \varepsilon \omega r}{2} E_0 \cos(\omega t) dr 
  = \mu \varepsilon h E_0 \left( \frac{\omega r}{2} \right)^2 \sin(\omega t)
\end{align}
Chiamiamo $ E_z(r = 0) = E_1 $ e semplificando $ h $ otteniamo una espressione per $ E_z(r \neq 0) $ 
\begin{align}
  \v{ E }_2 \coloneqq \v{ E }_z(r \neq 0) = \v{ E }_1 - \mu \varepsilon E_0 \left( \frac{\omega r}{2} \right)^2 \sin(\omega t) \vers{z}
  = E_0 \left( 1 - \mu \varepsilon \left( \frac{\omega r}{2} \right)^2  \right) \sin(\omega t) \vers{z}
\end{align}

% condensatore con profilo di parabola
\begin{figure}[h]
  \center
  \begin{tikzpicture}
    % parametri
    \def \r{1.5}
    \def \h{2}
    \def \e{0.4}

    % ellissi
    % Armatura superiore
    \draw (-\r,\h) arc[start angle=180, end angle=360, x radius=\r, y radius=\e];
    \draw[dashed] (-\r,\h) arc[start angle=180, end angle=0, x radius=\r, y radius=\e];

    % Armatura inferiore
    \draw (-\r,0) arc[start angle=180, end angle=360, x radius=\r, y radius=\e];
    \draw[dashed] (-\r,0) arc[start angle=180, end angle=0, x radius=\r, y radius=\e];
    
    % Fili verticali
    \draw (0,\h) -- (0,\h+1);
    \draw [dashed] (0,0) -- (0,-\e);
    \draw (0,-\e) -- (0,-1);

    % profilo del campo
    \pgfmathsetmacro{\a}{\h/(2*\r*\r)}

    %\draw[domain=-\r:\r, samples=200, thick, black] plot (\x, {\h - \a*\x*\x});

    % funzione
    \def \f {\h - 3/4 * \a* \x *\x}
    \draw [domain = - \r : \r , samples = 200, dashed, black] plot (\x, {\f});
    
    % frecce campo elettrico
    \foreach \x in {-1.3 , -1, -0.5, 0, 0.5, 1, 1.3}{  \draw[->] (\x,0) -- (\x,{\f});}


  \end{tikzpicture}
  \caption{Profilo del campo elettrico nel condensatore}
  \label{figure: condensatore piano con profilo}
\end{figure}



Notiamo che $ \v{ E }_2 $ differisce da $ \v{ E }_1 $ per un coefficiente 
$ 1 - \alpha^2 $ con $\displaystyle \alpha = \frac{\omega r}{2} \sqrt{\mu \varepsilon} = \frac{\omega r}{2 c} n $ dove $ n \coloneq \sqrt{\eps_r \mu_r}$ 
è l'indice di rifrazione del mezzo che per ora ha il solo ruolo di una costante reale; per la sua definizione rimandiamo al capitolo 
\ref{chap: elettromagnetismo nella materia}. Grazie a ciò siamo in grado di dare una condizione di 
\textit{\enquote{lentamente variabili}}: quando $ \alpha $ è piccolo possiamo considerare solo il termine $ \v{ E }_1 $. La correzione $ \alpha $ è una 
quantità adimensionale dunque la consideriamo piccola rispetto ad $ 1 $, tuttavia dopo ci soffermeremo cosa davvero stiamo richiedendo che sia piccolo 
in questa quantità. Possiamo stimare l'ordine di grandezza limite della frequenza (quando non è più trascurabile) in funzione del raggio $ R $ del 
condensatore: approssimando $ \omega = 2 \pi f \approx 6 f $ si ha
\begin{align}
  \alpha^2 = \left( \frac{ 6 f \cdot r }{2 \cdot 3 \cdot 10^8} \right)^2 \leq 10^{-2} \implies f \leq \frac{10^7}{R} \si{\Hz}
\end{align}
quindi per un condensatore di raggio $ R = \SI{1}{\cm} $ otteniamo che per frequenze $ f \leq \SI{e9}{\Hz} = \SI{1}{\GHz} $ la correzione al campo elettrico 
è al massimo dell'$ 1\% $. Se ci soffermiamo su quello che abbiamo richiesto, notiamo che scrivendo $ \alpha $ in termini della lunghezza d'onda 
$ \lambda = \frac{v}{f} = \frac{2\pi c}{n\omega} $ si ha che 
\begin{align}
  \alpha = \frac{\pi r}{\lambda} 
\end{align}
e dunque la richiesta che abbiamo fatto equivale a 
\begin{align}
  \alpha \ll 1 \implies R \ll \lambda \quad \text{oppure} \quad \omega R \ll c
\end{align}
cioè che il condensatore sia "invisibile" rispetto alla lunghezza d'onda, o che la scala di velocità $ \omega R $ non sia comparabile a quella della luce. 
Questa seconda condizione acquista maggiormente significa riscrivendo in termini del periodo $ T = \frac{2\pi}{\omega} $ ottenendo
\begin{align}
  \frac{2R}{c} \ll \frac{T}{\pi}
\end{align}
dove $ 2R/c $ è il tempo che la luce impiega ad attraversa il condensatore, che richiediamo essere molto più piccolo rispetto al tempo di variazione 
dei campi.

\section{Onde elettromagnetiche piane} \label{section:_Onde_elettromagnetiche}
In questa sezione risolveremo le equazioni di Maxwell nel vuoto e in assenza di sorgenti, e questo ci porterà a delle soluzioni che possono sembrare 
inutili dato che se non ci sono sorgenti è naturale chiedersi come si sono generate. Tuttavia le soluzioni che troveremo, chiamate \textit{onde piane}, sono 
il limite per quando siamo molto lontani da correnti e cariche libere, in cui le soluzioni sono localmente onde piane.

Dalla matematica sappiamo che il sistema $ \{1, \cos\left( \frac{n\pi}{a} x \right), \sin\left( \frac{n\pi}{a} x \right)\} $ è completo nello spazio di 
Hilbert $ L^2(-a, a) $, e quindi tramite serie di Fourier è possibile ricostruire una qualsiasi funzione di tale spazio. Allora una volta verificato che 
una funzione del tipo $ h = h_0\cos(kx - \omega t) $ (o il seno) è soluzione, dato che le equazioni sono lineari, possiamo comporle linearmente e scrivere 
ogni soluzione come serie di termini del tipo $ h $. Questo approccio è classico ma vogliamo provare ad affrontare il problema 
con un metodo più formale e che successivamente risulterà più comodo per i calcoli, il quale richiede una raffinatezza a livello teorico. 
Per snellire i calcoli indicheremo con $ g = g (x,t) = kx - \omega t $. Sfruttiamo la notazione complessa:
\begin{align}
  h = h_0 e^{i \left( kx - \omega t \right)} 
\end{align} \label{eq: onda piana monodimensionale}
con
\begin{equation}
  h_0 \in \mathbb{C},\; h_0 = \alpha + i \beta \quad \alpha, \beta \in \R
\end{equation}
In questo modo prendendo la parte reale di $ h $ (cerchiamo soluzioni nello spazio $ \mathbb{R} $ dato che i campi fisici sono reali):
\begin{align}
  \text{Re}\left( h \right) = \text{Re} \left[ (\alpha + i \beta)\left( \cos g + i \sin g \right) \right] = \alpha \cos g - \beta \sin g
\end{align}
Scegliendo il numero complesso $ h_0 $ otteniamo una combinazione lineare di seno e coseno di $ g(x,t) $, quindi la notazione complessa non ci fa 
effettivamente perdere generalità. Generalizzando nello spazio $ \R^3 $ l'oggetto in equazione \eqref{eq: onda piana monodimensionale} otteniamo 
($ h:\R^3 \to \mathbb{C} $)
\begin{align}
  h = h_0 e^{i (\v{k\cdot r} - \omega t)} \tag*{[Onda piana]}
\end{align} \label{eq: onda piana}
che è l'onda piana di cui parlavamo sopra, chiamata così perché il fronte d'onda, cioè il luogo dei punti con stessa immagine, è un piano. Sarà proprio 
questo oggetto l'elemento fondamentale con cui descriveremo le soluzioni alle equazioni di Maxwell, tuttavia non è banale il perché possiamo farlo dato che 
le onde piane non sono funzioni $ L^2 $. Dal punto di vista teorico stiamo cambiando lo spazio su cui ci appoggiamo, infatti mandando $ a \to \infty $ la 
nostra (discreta) serie di Fourier diventa una (continua) trasformata di Fourier. Per ora ci basta sapere che in questo nuovo spazio dove vivono le onde 
piane, la trasformata di una onda piana rimane ancora in questo spazio e dunque possiamo comporle per formare la soluzione generale, e questo chiude le 
giustificazioni della notazione complessa. 
Vediamo adesso come Maxwell ha manipolato le sue equazioni in modo da trovare l'equazione differenziale delle onde. Calcoliamo allora il rotore del 
rotore di $ \v{ E } $:
\begin{align}
  \v{\nabla}\!\times\! \left( \v{\nabla}\!\times\!\v{ E } \right) &= \v{\nabla}\!\times\! \left( - \der{\v{ B }}{t} \right)
  = - \mu_0 \varepsilon_0 \der[2]{\v{ E }}{t}
  \\
  \text{ma} \quad \v{\nabla}\!\times\! \left( \v{\nabla}\!\times\!\v{ E } \right) 
  &= \nabla \left( \nabla \cdot \v{ E } \right) - \nabla^2 \v{ E } = 0 - \nabla ^2 \v{ E }
\end{align}
e unendo questi due risultati si ottiene:
\begin{equation}
  \nabla^2\v{ E } -\mu_0 \varepsilon_0 \der[2]{\v{ E }}{t} = 0
\end{equation}
quindi ogni componente $ E_i $ del campo elettrico soddisfa l'equazione tridimensionale delle onde. Facendo lo stesso procedimento per il campo 
magnetico si ha che:
\begin{align}
  -\nabla^2 \v{ B } = \nabla\!\times\!\left( \nabla\!\times\!\v{ B } \right) 
  = \nabla \! \times \! \left( \mu_0 \varepsilon_0 \der{\v{ E }}{t} \right) = - \mu_0 \varepsilon_0 \der[2]{\v{ B }}{t}
\end{align}
e quindi
\begin{equation}
  \quad \nabla^2\v{ B } - \mu_0 \varepsilon_0 \der[2]{\v{ B }}{t} = 0
\end{equation}
Con queste due nuove equazioni, $ \v{ B } $ ed $ \v{ E } $ non sono accoppiati, quindi nel cercare le soluzioni dovremo poi verificare che 
risolvano effettivamente anche le equazioni di Maxwell (questo perché non sono effettivamente equivalenti). Per chiarire, se una funzione è soluzione 
delle equazioni di Maxwell, allora risolve anche le due equazioni delle onde, ma non è vero il viceversa.\\
Proviamo con le onde piane
\begin{align} \label{onde piane E,B}
  \v{ E }(\v{ r },t) &= \v{ E }_0 e^{i\left( \v{ k \cdot r } - \omega t \right)}
  \\
  \v{ B }(\v{ r },t) &= \v{ B }_0 e^{i\left( \v{ k \cdot r } - \omega t \right)}
\end{align}
ricordando che dobbiamo sempre prenderne la parte reale prima o poi. È facile verificare per quali condizioni sono soluzioni dell'equazione delle onde 
che abbiamo trovato ma per quanto detto sopra noi verificheremo che risolvono le equazioni di Maxwell:
\begin{align}
  \nabla\!\cdot\!\v{ E } &= i \left( E_{0x}k_x + E_{0y}k_y + E_{0z}k_z \right) e^{i\left( \v{ k \cdot r } - \omega t \right)} = 0 
  \iff \v{ k\cdot E }_0 = 0
  \\
  \nabla\!\cdot\!\v{ B } &= i \left( B_{0x}k_x + B_{0y}k_y + B_{0z}k_z \right) e^{i\left( \v{ k \cdot r } - \omega t \right)} = 0 
  \iff \v{ k \cdot B }_0 = 0
\end{align}
Da questo abbiamo capito che entrambi i campi sono trasversali rispetto a $ \v{ k } $. Per calcolare i rotori sfruttiamo la proprietà di 
essere autofunzioni dell'operatore derivata e ci convinciamo che
\begin{align}
  \nabla\!\times\!\v{ E } = det 
  \begin{pmatrix}
    \v{ u }_x & \v{ u }_y & \v{ u }_z \\
    \partial_x & \partial_y & \partial_z \\
    E_x & E_y & E_z
  \end{pmatrix}
  &= det
    \begin{pmatrix}
    \v{ u }_x & \v{ u }_y & \v{ u }_z \\
    ik_x& ik_y & ik_z \\
    E_x & E_y & E_z
  \end{pmatrix}
  \\= i \v{ k \times E }_0 e^{i\left( \v{ k \cdot r } - \omega t \right)} &= - \der{ \v{ B }}{t} 
  = i \omega \v{ B }_0 e^{i\left( \v{ k \cdot r } - \omega t \right)}
\end{align}
e otteniamo così una relazione sulla direzione e sul modulo:
\begin{equation}
  \v{ k \times E }_0 = \omega \v{ B }_0
\end{equation}
Con gli identici passaggi sulla quarta equazione si ottiene la relazione:
\begin{align}
  \nabla\!\times\!\v{ B } = i \v{ k\times B }_0 e^{i\left( \v{ k \cdot r } - \omega t \right)} 
  &= \mu_0 \varepsilon_0 \der{\v{ E }}{t}   = \mu_0 \varepsilon_0 (- i \omega ) \v{ E }_0 e^{i\left( \v{ k \cdot r } - \omega t \right)}
  \\
  \implies \v{ k\times B }_0 &= - \mu_0 \varepsilon_0 \omega \v{ E }_0
\end{align}
Non è un caso che otteniamo gli stessi risultati per $ \v{ B } $ ed $ \v{ E } $ ma è una conseguenza della simmetria delle equazioni di 
Maxwell \eqref{equazioni di Maxwell dinamiche}, il che significa che campo magnetico e campo elettrico sono due facce della stessa medaglia, l'uno 
la conseguenza (e causa) dell'altro.\\
Dalle onde meccaniche sapevamo già che $ |\v{ k }| = \frac{\omega}{v} $, con $ v $ la velocità di propagazione dell'onda, per trovarla in questo 
caso prendiamo il coefficiente davanti alla derivata temporale nell'equazione delle onde e otteniamo
\begin{equation} \label{definizione di c}
  \mu_0 \varepsilon_0 = \frac{1}{v^2} \implies v = \frac{1}{\sqrt{\mu_0 \varepsilon_0}} =: c
\end{equation}
e quindi le relazioni trovate prima si riscrivono come
\begin{align}
  \vers{k} \times \v{ E }_0 = c \v{ B }_0 \implies E = c B
\end{align}
dove $ \vers{k} $ è il versore di propagazione dell'onda, cioè $ \vers{k} = \frac{\v{ k }}{|\v{ k }|} $. A oggi è scontato che la luce sia una un'onda 
elettromagnetica ma al tempo di Maxwell non erano ancora stati fatti gli esperimenti che evidenziassero il comportamento da onda elettromagnetica della 
luce. Tuttavia la velocità della luce era nota e con un incertezza dell'ordine del $ 1\% $, dunque Maxwell disse che 
\textit{"non è possibile escludere che la luce stessa sia una perturbazione elettromagnetica"}.

\subsection{Gradi di libertà}
Soffermiamoci a vedere quante variabili sono necessarie per fissare in modo univoco una soluzione. La scelta della fase 
$ \v{ k } = (k_x,k_y,k_z) $ comporta 3 parametri ma $ \omega = c|\v{ k }| $ quindi nessun grado di libertà. Per quanto riguarda le due 
ampiezze sfruttiamo le relazioni trovate e prime e abbiamo che $ \v{ E }_0 \cdot \v{ k } = 0 $, quindi solo due 
gradi di libertà. Infine per il campo magnetico abbiamo già modulo e direzione fissati quindi non ci sono effettivamente dei parametri liberi, 
dunque in totale ne abbiamo $ 5 $.

\subsection{Polarizzazione}
Consideriamo il caso in cui $ \v{ k } = k_z \vers{z} $ dell'onda che si propaga lungo $ z $, $ \v{ E }_0 = (\mathcal{E}_x, 0 , 0)$ 
con $ \mathcal{E}_x \in \mathbb{R} $ (notiamo che $ \mathcal{E}_z $ è necessariamente nulla per la condizione di onda trasversale rispetto alla 
direzione di propagazione), allora il campo vale
\begin{equation}
  \v{ E } = Re \left( \v{ E }_0 e^{i\left( \v{ k\cdot x} - \omega t \right)} \right) 
  = \mathcal{E}_x \cos \left( k_z z - \omega t \right) \vers{x}
\end{equation}
che è detta \textbf{polarizzazione lineare}, in questo caso lungo $ x $. Con $ \v{ E }_0 =  \mathcal{E} \left( 1, i, 0 \right) $, 
$ \mathcal{E}\in \mathbb{R} $ il campo vale:
\begin{equation}
  \v{ E } = \mathcal{E} \cos \left( k_z z - \omega t \right) \vers{x} - \mathcal{E} \sin \left( k_z z - \omega t \right) \vers{y}
\end{equation}
che se la visualizziamo in un piano $ xy $ è un vettore di modulo $ \mathcal{E} $ che sta ruotando con velocità angolare $ \omega $. Questa 
è detta \textbf{polarizzazione circolare}, in generale per coefficienti diversi per seno e coseno si ha la \textbf{polarizzazione ellittica}.

\textbf{Non polarizzata}: luce che varia con una successione che varia molto lentamente rispetto alla frequenza angolare dell'onda ma molto 
rapidamente rispetto alla sensibilità dei nostri strumenti. Ad esempio possiamo aggiungere una fase stocastica 
\begin{align}
  \v{ E } = \mathcal{E}_0 \left( \vers{x} + \vers{y} e^{i\delta(t)}\right) e^{i (kz - \omega t)}
\end{align}
in questo modo la luce risultante è una sovrapposizione di un fascio lungo $ \vers{x} $ e un treno d'onda lungo $ \vers{y} $. Facendone la media 
utilizzando la media per i campi complessi sopravvive solo il termine di $ x $ con se stesso, il termine di $ y $ con se stesso, ma vanno a zero i 
termini incrociati:
\begin{align}
  \langle \v{ E }^2 \rangle = \frac{1}{2} \v{ E } \cdot \v{ E }^*
\end{align}

Fase relativa costante $ \implies $ coerenti, altrimenti non sono coerenti. Quando sono incoerenti posso sommare le intensità perché sto supponendo 
che i fenomeni di interferenza non esistono, ma nella realtà questo è vero nel caso in cui il mio strumento faccia la misura (la media) su una scala 
maggiore a quella a cui avvengono i fenomeni di interferenza. 
\subsection{Somma di onde progressive e regressive}
Nella nostra trattazione supponiamo sempre che oltre a essere reale, $ \omega $ sia positiva, che è una usuale convenzione dato che per ottenere la 
frequenza negativa basta mandare $ t \to -t $. Per quanto riguarda $ \v{ k } $, in base al segno otteniamo l'onda progressiva o l'onda regressiva, 
riferito alla direzione di propagazione. Vogliamo vedere cosa succede se combiniamo due soluzioni, una progressiva e l'altra regressiva.
Consideriamo le soluzioni progressive ($ \v{ k }=k \vers{z} $)
\begin{equation}
  \v{ E }_+(z,t) = A_0  e^{i\left( kz-\omega t \right)}\vers{x} \implies \v{ B }_+ = A_0  e^{i\left( kz-\omega t \right)}\vers{y}
\end{equation}
mentre le corrispondenti soluzioni regressive si ottengono mandando $ k \to -k $ e stando attenti che $ B_0\to -B_0 = A_0 $ dato che 
$ \omega \v{ B }_0 = \v{ k\times E }_0 $, cioè
\begin{equation}
  \v{ E }_-(z,t) = A_0  e^{i\left( -kz-\omega t \right)}\vers{x} \implies \v{ B }_- = -A_0  e^{i\left( -kz-\omega t \right)}\vers{y}
\end{equation}
Sommando, i campi valgono
\begin{align}
  \v{ E } &= \v{ E }_+ + \v{ E }_- = A_0 \vers{x} e^{-i\omega t}\left( e^{ikz} + e^{-ikz} \right)= A_0\vers{x}e^{-i\omega t}2 \cos (kz)
  \\
  \v{ B } &= \v{ B }_+ + \v{ B }_- = A_0 \vers{x} e^{-i\omega t}\left( e^{ikz} - e^{-ikz} \right)= A_0\vers{x}e^{-i\omega t}2\, i \sin(kz)
\end{align}
e prendendone la parte reale otteniamo:
\begin{align}
  \v{ E } &= 2A_0 \cos(\omega t) \cos (kz) \vers{x}
  \\
  \v{ B } &= 2A_0 \sin(\omega t) \sin (kz) \vers{y}
\end{align}
Possiamo visualizzare questi campi in un piano cartesiano mettendo in ascissa la $ z $, in ordinata i campi e li disegniamo in vari tempi. La cosa 
apparentemente strana è che ci sono dei punti per ogni campo in cui questi sono sempre nulli, ad esempio per $ z = n\pi, n\in \mathbb{N} $ il campo 
magnetico è sempre nullo mentre quello elettrico varia nel tempo. Sembra contraddire l'accoppiamento dei due campi dettato dalle leggi di Maxwell, 
che a un campo elettrico prevedono sempre associato un campo magnetico e viceversa. Tuttavia la soluzione al (presunto) paradosso è che sì è vero 
il fatto dell'accoppiamento dei campi ma dobbiamo sempre ricordare che le equazioni di Maxwell sono delle relazioni \textbf{locali} fra i campi, 
non puntuali e nemmeno istantanee, ma vanno sempre considerati degli intorni.

\section{Teorema di Poynting}
In questa sezione otterremo un risultato fondamentale dell'elettrodinamica, che riguarda l'energia associata al campo elettromagnetico, noto 
come \textit{Teorema di Poynting}. Tale teorema è una conseguenza delle equazioni di Maxwell, non ci fornisce relazioni ulteriori fra i campi 
ma bensì mette in chiaro il fatto che i campi sono un'unica identità e come sono legati alla trasmissione di energia.

Consideriamo un volume $ V $ con superficie regolare $ \partial V = A $ orientata uscente, dove sono presenti delle sorgenti. 
Svolgiamo i seguenti prodotti scalari: $ \v{ E } $ con il rotore di $ \v{ B } $ e $ \v{ B } $ con il rotore di $ \v{ E } $.
\begin{align}
  \v{ E } \cdot \left( \nabla\!\times\!\v{ B } \right) &= 
  \v{ E } \cdot \v{ J } + \mu_0 \varepsilon_0 \v{ E } \cdot \der{\v{ E }}{t}
  \\
  \v{ B } \cdot \left(\nabla\!\times\!\v{ E } \right) &= - \v{ B } \cdot \der{\v{ B }}{t}
\end{align}
Sottraiamo la prima alla seconda applichiamo l'identità vettoriale \eqref{identity: div(A times B)}:
\begin{align}
  &\v{ B } \cdot \left(\nabla\!\times\!\v{ E } \right) - \v{ E } \cdot \left( \nabla\!\times\!\v{ B } \right) 
  = \nabla \! \cdot \! \left( \v{ E\times B } \right) 
  \\
  &= - \mu_0 \v{ E \cdot J } - \v{ B } \cdot \der{\v{ B }}{t} - \mu_0 \eps_0 \v{ E } \cdot \der{\v{ E }}{t}
  \\
  &= - \mu_0 \v{ E \cdot J} - \mu_0 \der{}{t} \left( \frac{\v{ B }^2}{2\mu_0} +  \frac{\eps_0\v{ E }^2}{2} \right)
\end{align}
In questa espressione riconosciamo la densità di energia elettrica $ u_e $ e allo stesso modo definiamo la densità di energia magnetica $ u_m $ 
(entrambe nel vuoto):
\begin{align}
  u_e \coloneqq  \frac{\eps_0\v{ E }^2}{2}
  \\
  u_m \coloneqq \frac{\v{ B }^2}{2\mu_0}
\end{align}
riscrivendo in termini di $ u $ otteniamo la relazione
\begin{align} \label{eq:teorema_di_Poynting_locale}
  \nabla\! \cdot \! \left( \v{ E } \times \frac{\v{ B }}{\mu_0} \right) + \v{ E \cdot J } = - \der{u}{t}
\end{align}
che è la forma locale del Teorema di Poynting. Applicando il teorema della divergenza passiamo alla sua formulazione globale su un volume $ V $:
\begin{align}
  \int_V \nabla \! \cdot \!\left( \v{ E } \times \frac{\v{ B }}{\mu_0}  \right) \, dV + \int_V \v{ E \cdot J } \,dV &=  -\der{}{t} \int_V u \,dV
  \\
  \implies \oint \left( \v{ E \times } \frac{\v{ B }}{\mu_0} \right) + W &= - \der{U}{t}
\end{align}
dove $ W $ rappresenta la potenza ceduta alle cariche e $ U $ è l'energia elettromagnetica immagazzinata nel volume (nei campi). Analizzando i termini 
del teorema di Poynting osserviamo che esprime il principio di conservazione dell'energia in termini dell'energia trasferita alle cariche, l'energia dei 
campi e un certo flusso uscente dal volume. Questa quantità viene definita come \textit{vettore di Poynting} $ \v{ S } $:
\begin{align}
  \v{ S } \coloneqq \frac{1}{\mu_0} \left(  \v{ E \times B } \right) = \v{ E \times H }
\end{align}
Sotto una certa interpretazione si può pensare il vettore di Poynting come un flusso di energia, ma comparendo nel teorema solo la sua divergenza potremmo 
aggiungere un qualsiasi rotore di un campo vettoriale (che ha quindi divergenza nulla) e il teorema rimarrebbe invariato. Benché questo non abbia conseguenze 
fisiche, dal solo teorema di Poynting il campo $ \v{ S } $ non è unico, infatti per dirlo servono delle considerazioni relativistiche.

Abbiamo detto che il teorema di Poynting rappresenta quindi una sorta di equazione di continuità per l'energia, ma proprio perché i campi scambiano energia 
con le cariche vogliamo trovare le implicazioni dinamiche di questo fatto. 
Consideriamo la forza di Lorentz che agisce su una singola carica $ q $:
\begin{align}
  \v{ F } = q(\v{ E } + \v{ v \times B}),
\end{align}
se la scriviamo in termini delle distribuzioni di carica e corrente $ \rho, \v{ J } $ otteniamo
\begin{align}
  \v{ F } = \int_V \v{ f } \,dV = \int_V \rho\v{ E } + \v{ J \times B } \,dV
\end{align}
Sostituiamo le due densità applicando le equazioni di Maxwell
\begin{align}
  f &= \eps_0 \v{ E } \left( \nabla\!\cdot\!\v{ E } \right) + \left( \nabla\!\times\!\frac{\v{ B }}{\mu_0} - \eps_0 \der{\v{ E }}{t} \right) \times \v{ B }
  \\
  &= \eps_0 \left[ \v{ E } \left( \nabla\!\cdot\!\v{ E } \right) + c^2 \left( \nabla\!\times\!\v{ B } \right) \times \v{ B } 
  -  \der{\v{ E }}{t} \times \v{ B } \right]
  \\
  &= \eps_0 \left[ \v{ E } \left( \nabla\!\cdot\!\v{ E } \right) + c^2 \left( \nabla\!\times\!\v{ B } \right) \times \v{ B } 
  - \der{\v{ \left( E \times B \right)}}{t} + \v{ E } \times \der{\v{ B }}{t} \right]
  \\
  &= \eps_0 \left[ \v{ E } \left( \nabla\!\cdot\!\v{ E } \right) + c^2 \left( \nabla\!\times\!\v{ B } \right) \times \v{ B } 
  - \der{\v{ \left( E \times B \right)}}{t} - \v{ E } \times \left( \nabla\!\times\!\v{ E } \right) \right]
\end{align}
Per il secondo principio della dinamica abbiamo che
\begin{align}
  \frac{d\v{p}}{dt} = \v{F} = \int_V \v{f} \,dV 
\end{align}
dove $ \v{p} $ è il momento totale delle cariche del sistema considerato, dunque è il \textit{momento meccanico} 
(spesso indicato con $ \v{p}_{\text{mecc}} $). A questo punto è naturale riconoscere nei termini che compongono $ \v{f} $ una densità di momento,  
ovvero il termine con la derivata temporale, che indicheremo con $ \v{g} $, mentre il suo integrale di volume lo indicheremo con $ \v{G} $. 
Tali quantità rappresentano rispettivamente la densità di momento e il momento del campo elettromagnetico, e cioè $ \v{p}_{\text{em}} \coloneq \v{G}  $: 
\begin{align}
  \v{g} \coloneq \eps_0 \left( \v{E \times B} \right) = \frac{\v{S}}{c^2} \qquad \v{G} = \int_V \v{g} \,dV
\end{align}
dove abbiamo evidenziato la proporzionalità col vettore di Poynting. In questi termini il secondo principio si riscrive come
\begin{align}
  \frac{d}{dt} \left( \v{p}_{\text{mecc}} + \v{G} \right) 
  = \int_V \eps_0 \left[ \v{ E } \left( \nabla\!\cdot\!\v{ E } \right) + c^2 \left( \nabla\!\times\!\v{ B } \right) \times \v{ B } 
  - \v{ E } \times \left( \nabla\!\times\!\v{ E } \right) \right] \,dV
\end{align}
Adesso manipoliamo il membro di destra aggiungendo la quantità nulla $ 0 = c^2 \v{ B } \left( \nabla\!\cdot\!\v{ B } \right) $ perché ci permette di 
ottenere una espressione che è simmetrica per i due campi (a meno del fattore $ c^2 $):
\begin{align}
  \eps_0 \left[ \v{ E } \left( \nabla\!\cdot\!\v{ E } \right) + c^2 \v{ B } \left( \nabla\!\cdot\!\v{ B } \right) 
  -  c^2 \, \v{ B } \times \left( \nabla\!\times\!\v{ B } \right)  
  - \v{ E } \times \left( \nabla\!\times\!\v{ E } \right) \right].
\end{align}
Da qui il nostro obiettivo è quello di compattare questa espressione e per farlo sfruttiamo un risultato noto per i campi vettoriali, il quale ci permette 
di costruire un tensore a due indici:
\begin{align}
  \left[ \v{ E } \left( \nabla \cdot \v{ E } \right) - \v{ E } \times \left( \nabla \times \v{ E } \right) \right]_i 
  = \sum_{j} \der{}{x_j} \left( E_i E_j - \frac{1}{2} \v{ E }^2 \delta_{i j} \right)
  \\
  \implies T_{i j} \coloneqq \eps_0 \left[ E_i E_j + c^2 \, B_i B_j - \frac{1}{2} \left( \v{ E }^2 + c^2 \, \v{ B }^2 \right) \delta_{i j} \right]
\end{align}
e il tensore $ T_{i j} $ è detto \textit{tensore di Maxwell} $ T_{i j} $ (o \textit{tensore degli sforzi}). La componente $ i $-esima $ f_i $ della densità 
di forza si riscrive come
\begin{align}
  f_i = \sum_j \der{T_{i j}}{x_j} - \der{g_i}{t}
\end{align}
e riportandola in termini della formulazione iniziale si ha
\begin{align}
  \der{}{t} \left( \frac{\v{S}}{c^2} \right) + \rho \v{ E } + \v{ J } \times \v{ B } = \sum_{i,j} \der{T_{i j}}{x_j} \vers{x}_i
\end{align}
A questo punto per ottenere il secondo principio in forma compatta manca solo da integrare sul volume la densità di forza $ \v{f} $:
\begin{align}
  \frac{d}{dt} \left( \v{ p }_{\text{mecc}} + \v{ p }_{\text{em}} \right) 
  = \sum_{i,j} \int_V \der{T_{i j}}{x_j} \vers{x}_i \,d^3r = \sum_{i,j} \oint_{\partial V} T_{i j} n_j \vers{x}_i \,d^2r
\end{align}
dove abbiamo usato il teorema della divergenza sul tensore di Maxwell, che risulta più chiaro nella seguente forma per la componente $ i $-esima
\begin{align}
  \frac{d}{dt} \left( p_i^{\text{mecc}} + p_i^{\text{em}} \right) 
  = \sum_{j} \int_V \der{T_{i j}}{x_j} \,d^3r = \sum_{j} \oint_{\partial V} T_{i j} n_j \,d^2r.
\end{align}
Il risultato che abbiamo ottenuto ci dice che la quantità di moto (derivata rispetto al tempo) contenuta nel volume è uguale tensore di Maxwell 
valutato sulla superficie, questo perché l'informazione che è contenuta sul bordo è quella dei campi iniziali (che evolvono nel tempo) e quella della 
materia che reagisce ai campi modificandoli.

\subsection{Vettore di Poynting per il filo}
Nel caso del filo percorso da densità di corrente $ \v{ J } $ uniforme, $ \v{ E } $ è lungo il filo ($ \vers{z} $) mentre $ \v{ B } $ è 
lungo $ \vers{\varphi} $, dunque il vettore di Poynting è 
\begin{align}
  \v{ S } = \frac{1}{\mu_0} \v{ E } \times \v{ B } = S\; \vers{z} \times \vers{\varphi} = - S \, \vers{\rho}
\end{align}
quindi è diretto verso il centro, il che è abbastanza controintuitivo.

Il vettore $ \v{ S } $ non è additivo, infatti consideriamo due sorgenti che generano due coppie di campi
\begin{align}
  \v{ E }_1 , \v{ B }_1 \qquad \v{ E }_2 , \v{ B }_2
\end{align}
il vettore di Poynting risultante non è banalmente la somma dei due vettori $ \v{ S }_1, \v{ S }_2 $ ma è
\begin{align}
  \v{ S } = \frac{1}{\mu_0} \left[ (\v{ E }_1 + \v{ E }_2) \times (\v{ B }_1 + \v{ B }_2 ) \right]
\end{align}


\section{Cavità elettromagnetiche o radiofrequenza (RF)}
Nella sezione \ref{section: campi lentamente variabili} abbiamo visto che in un condensatore a facce circolari il profilo del campo elettrico è quello di 
una parabola rovesciata. Abbiamo inoltre trovato una relazione che ci dice se i campi possono essere considerati lentamente variabili, ovvero
\begin{align}
  \alpha  \ll 1 \qquad \text{con} \qquad \alpha = \frac{\omega R}{2 c}
\end{align}
dove $ \omega $ è la frequenza angolare con la quale oscillano i campi e $ R $ è il raggio del condensatore.
Vogliamo capire cosa succede nel limite opposto, cioè quando i campi sono rapidamente variabili all'interno del condensatore. Dato che stiamo considerando 
frequenze molto alte possiamo supporre senza perdere rigorosità che il campo sul bordo del condensatore sia nullo, questo perché la lunghezza d'onda è 
inversamente proporzionale a $ \omega $:
\begin{align}
  \lambda = \frac{2\pi v}{\omega} \qquad k = \frac{2\pi}{\lambda},
\end{align}
quindi basta scegliere una frequenza tale che il diametro del condensatore sia un multiplo di $\lambda/2$. In formule
\begin{align}
  R = n \, \frac{\lambda}{2} = \frac{n \pi c}{\omega} \implies \omega_n = \frac{n \pi c}{R}
\end{align}
quindi prendendo una qualsiasi $ \omega_n $ avremo che il campo elettrico sarà nullo sul bordo, come mostrato in figura 
\ref{figure:_condensatore_piano_cavità} dove abbiamo preso 
$ n = 1 $ ottenendo quindi $ \lambda = 2 R $. Il ragionamento che abbiamo appena fatto parte dall'assunzione che le soluzioni siano effettivamente delle 
sinusoidi, il che non è vero per il condensatore cilindrico. Infatti le soluzioni esatte di questo modello sono le funzioni di Bessel e per vederlo si 
rimanda all'esercizio \ref{section:_esercizio_cavità_cilindrica}. Benché esistano alcuni limiti in cui le funzioni di Bessel diventano delle funzioni 
trigonometriche, il ragionamento appena fatto è utile solo ai fini motivazionali di chiudere la cavità e sapere che per alcune condizioni il campo 
elettrico tangenziale è nullo sul bordo.

\input{build/figure/condensatore_piano_cavità.tex}

Adesso che abbiamo dato le condizioni per avere il campo nullo al bordo possiamo chiudere il condensatore con un metallo di spessore $ \delta $ e un 
sistema così fatto è detto \textit{cavità}, rappresentata in figura \ref{figure:_cavità:_A}. Vogliamo studiare questo tipo di sistema considerando due casi: 
\begin{itemize}
  \item quando la cavità è considerata ideale e quindi all'interno delle pareti metalliche il campo elettrico è nullo;
  \item quando non può più essere considerata ideale per via del campo elettrico e quindi si aggiungerà un effetto di dissipazione;
\end{itemize}

Calcoliamo la circuitazione dei campi su un circuito amperiano $ \gamma $ al bordo della cavità come in figura \ref{figure:_cavità:_B}.
\begin{figure}[h]
  % parametri
  \def \r{1.5}
  \def \h{2}
  \def \e{0.4}
  \def \d {1}
  % sotto-figura 1
  \begin{subfigure}{0.48\textwidth}
    \centering
    \begin{tikzpicture}
      % ellissi
      % Armatura superiore
      \draw (-\r,\h) arc[start angle=180, end angle=360, x radius=\r, y radius=\e];
      \draw[dashed] (-\r,\h) arc[start angle=180, end angle=0, x radius=\r, y radius=\e];

      % Armatura inferiore
      \draw (-\r,0) arc[start angle=180, end angle=360, x radius=\r, y radius=\e];
      \draw[dashed] (-\r,0) arc[start angle=180, end angle=0, x radius=\r, y radius=\e];
      
      % Fili verticali
      \draw (0,\h) -- (0,\h+1);
      \draw [dashed] (0,0) -- (0,-\e);
      \draw (0,-\e) -- (0,-1);

      % bordi cavità
      \draw (\r, 0) -- (\r , \h);
      \draw (-\r, 0) -- (-\r , \h);
    \end{tikzpicture}
    \caption{Cavità cilindrica}
    \label{figure:_cavità:_A}

  \end{subfigure} % lasciare una riga vuota fa andare a capo la seconda figura (e \hfill)
  \hfill
  % sotto-figura 2
  \begin{subfigure}{0.48\textwidth}
    \centering
    \begin{tikzpicture}

      \def \hb {1.5*\h}
      \def \l {0.3}
      \begin{scope}[shift = {(0,1)}]
        \draw (0,0) -- (0,\hb) -- (\d,\hb) -- (\d,0) -- (0,0);
        % circuito amperiano tratteggiato con frecce
        \draw [dashed, decorate, decoration={mark=at position 0.5 with {\arrow[triangle 45]{>}}}] (-\l, \hb/4) -- (-\l,\hb* 3/4);
        \draw [dashed, decorate, decoration={mark=at position 0.5 with {\arrow[triangle 45]{>}}}] (-\l,\hb* 3/4) -- (\l, \hb * 3/4);
        \draw [dashed, decorate, decoration={mark=at position 0.5 with {\arrow[triangle 45]{>}}}] (\l, \hb * 3/4) -- (\l, \hb/4 );
        \draw [dashed, decorate, decoration={mark=at position 0.5 with {\arrow[triangle 45]{>}}}] (\l, \hb/4 ) -- (-\l , \hb/4);
      \end{scope}
      % Aggiungo una linea invisibile per "riempire" lo spazio sotto
      \path[draw=none] (0,0) -- (0,2);
      % Etichette
      \node at (-\l - 0.3 ,\hb/4 + 1.2) {$\gamma$};
      \node at (\d/2,0.7) {$\delta$};
    \end{tikzpicture}
    \caption{Sezione della cavità e circuito amperiano}
    \label{figure:_cavità:_B}
  \end{subfigure}
\end{figure}

\begin{align}
  E_z(\text{metallo}) - E_z(\text{cavità}) = 0 &\qquad E_\rho(\text{metallo}) - E_\rho(\text{cavità}) = \frac{\sigma}{\eps_0}
  \\
  B_z(\text{metallo}) - B_z(\text{cavità}) = \mu_0 K &\qquad B_\rho(\text{metallo}) - B_\rho(\text{cavità}) = 0
\end{align}
Queste sono le condizioni al bordo che dovremo imporre in seguito. Supponiamo di avere una onda monocromatica 
$ f(t) \sim \exp \left\{ -i \omega t \right\} $, l'equazione da risolvere la troviamo facendo il rotore del rotore del campo elettrico:
\begin{align}
  \nabla\!\times\!\v{ (\nabla\!\times\!\v{ E }) } &= \nabla \! \times \! \left( i \omega \v{ B } \right) 
  = i \omega \left( -i \omega \v{ E } \right)
  \\
  \nabla\!\times\!\v{ (\nabla\!\times\!\v{ E }) } &= \nabla (\nabla \cdot \v{ E }) - \nabla ^2 \v{ E } = - \nabla ^2 \v{ E }
  \\
  \implies &\nabla^2 \v{ E } = - \omega^2 \v{ E }
\end{align}
che si tratta di una equazione agli autovalori per l'operatore differenziale laplaciano. Per quanto riguarda il campo magnetico vale la stessa relazione 
trovata sopra: 
\begin{align}
  \v{ B } = - \frac{i}{\omega} \nabla\!\times\!\v{ E } \implies \nabla\!\cdot\!\v{ B } = 0 \implies B_\rho = 0
\end{align}
che sono proprio le condizioni che dovevamo imporre, dunque ci possiamo concentrare solo sul campo elettrico. Dalla matematica sappiamo che c'è un 
insieme numerabile di modi propri $ \omega_n $ che costituiscono le soluzioni dell'equazione, le quali costituiscono una base (numerabile) ortogonale e 
completa. Nel seguito forniremo le soluzioni a questa equazione in coordinate cartesiane.

Per quanto riguarda l'energia all'interno della cavità, sappiamo dalla trattazione sulle onde elettromagnetiche che
\begin{align}
  U = \int_V u \,dV = \int \frac{\eps_0 \v{ E }^2}{2} + \frac{\v{ B }^2}{2\mu_0} \,dV \implies \langle U \rangle 
  = \int \frac{\eps_0 \langle \v{ E }^2 \rangle}{2} + \frac{\langle \v{ B }^2 \rangle}{2\mu_0} \,dV
\end{align}
e vogliamo mostrare che i due termini elettrico e magnetico forniscono esattamente lo stesso contributo, cioè
\begin{align}
  \int_V u_E \,dV = \int_V u_B \,dV.
\end{align}
\begin{align}
  \int_V \langle \v{ B }^2 \rangle \,dV &= \int_V \frac{1}{2} \v{ B } \v{ B }^* \,dV 
  = \int_V \frac{1}{2} \left( \frac{i}{\omega} \nabla\!\times\!\v{ E } \right) \left( \frac{i}{\omega} \nabla\!\times\!\v{ E } \right)^* \,dV 
  \\
  &= \frac{1}{2 \omega^2} \int_V \frac{1}{2} \left(  \nabla\!\times\!\v{ E } \right) \left( \nabla\!\times\!\v{ E } \right)^* \,dV \dots
\end{align}

\subsection{Cavità in coordinate cartesiane}
\begin{comment}
\begin{tikzpicture}[scale=1.2,
    x={(1cm,0cm)},
    y={(0.6cm,0.4cm)},
    z={(0cm,1cm)}
]

    % Parametri
    \def\a{4.1}
    \def\b{4}
    \def\d{2}

    % Assi cartesiani
    \draw[->] (0,0,0) -- (5,0,0) node[below left] {$x$};
    \draw[->] (0,0,0) -- (0,3,0) node[below right] {$z$};
    \draw[->] (0,0,0) -- (0,0,-4.5) node[right] {$y$};

    % Vertici del parallelepipedo
    \coordinate (O) at (0,0,0);
    \coordinate (A) at (\a,0,0);
    \coordinate (B) at (\a,\d,0);
    \coordinate (C) at (0,\d,0);

    \coordinate (O1) at (0,0,-\b);
    \coordinate (A1) at (\a,0,-\b);
    \coordinate (B1) at (\a,\d,-\b);
    \coordinate (C1) at (0,\d,-\b);

    % Spigoli base
    \draw (O) -- (A) -- (B) -- (C) -- cycle;

    % Spigoli verticali
    \draw (O) -- (O1);
    \draw (A) -- (A1);
    \draw (B) -- (B1);
    \draw (C) -- (C1);

    % Spigoli superiori
    \draw (O1) -- (A1) -- (B1) -- (C1) -- cycle;

    % Etichette dimensioni
    \node at (\a/2, -0.3, 0) {$a$};
    \node at (\a+0.3,0,-\b/2) {$b$};
    \node at (-0.3,\d/2,0) {$d$};


    \end{tikzpicture}

\end{comment}
\begin{figure}
    \center
    \begin{tikzpicture}[ % FISSA IL RIFERIMENTO IN MODO CORRETTO, ALTRIMENTRI --> (X,Z,Y)
        x={(1cm,0cm)},
        y={(0.6cm,0.4cm)},
        z={(0cm,1cm)}
        ]
            % Parametri
        \def\a{4}
        \def\b{3.5}
        \def\d{2}

        % Assi cartesiani
        \draw[->] (0,0,0) -- ({\a + 1},0,0) node[below left] {$x$};
        \draw[->] (0,0,0) -- (0,{\b + 1},0) node[below right] {$y$}; % below/above
        \draw[->] (0,0,0) -- (0,0,{\d + 1}) node[left] {$z$};

        % Vertici del parallelepipedo
        \coordinate (O) at (0,0,0);
        \coordinate (A) at (\a,0,0);
        \coordinate (B) at (\a,\b,0);
        \coordinate (C) at (0,\b,0);

        \coordinate (O1) at (0,0,\d);
        \coordinate (A1) at (\a,0,\d);
        \coordinate (B1) at (\a,\b,\d);
        \coordinate (C1) at (0,\b,\d);

        % Spigoli base
        \draw (O) -- (A) -- (B) -- (C) -- cycle;

        % Spigoli verticali
        \draw (O) -- (O1);
        \draw (A) -- (A1);
        \draw (B) -- (B1);
        \draw (C) -- (C1);

        % Spigoli superiori
        \draw (O1) -- (A1) -- (B1) -- (C1) -- cycle;

        % Etichette dimensioni
        \node at (\a/2, -0.5, 0) {$a$};
        \node at (\a+0.5, \b/2, 0) {$b$};
        \node at (-0.3, 0, \d/2) {$d$};

    \end{tikzpicture}
    \caption{Cavità a forma di parallelepipedo}
    \label{figure:_cavità_parallelepipedo}
\end{figure}
Vogliamo risolvere l'equazione trovata nella sezione precedente per una cavità a forma di parallelepipedo di lati $ a > b > d $ allineati (in ordine) lungo 
$ xyz $ (figura \ref{figure:_cavità_parallelepipedo}). Proviamo una soluzione del tipo 
\begin{align}
  \v{ E }(\v{ x }, t) = e^{-i\omega t}& \left[  \mathcal{E}_x \cos\left( \frac{n_1\pi}{a}x \right) \sin\left( \frac{n_2\pi}{a}y \right) 
  \sin\left( \frac{n_3\pi}{a}z \right)\vers{x} \right.
  \\
  &+  \mathcal{E}_y \sin\left( \frac{n_1\pi}{a}x \right) \cos\left( \frac{n_2\pi}{a}y \right) 
  \sin\left( \frac{n_3\pi}{a}z \right)\vers{y}
  \\
  &+ \left. \mathcal{E}_z \sin\left( \frac{n_1\pi}{a}x \right) \sin\left( \frac{n_2\pi}{a}y \right)
  \cos\left( \frac{n_3\pi}{a}z \right)\vers{z} \right]
\end{align}
vogliamo verificare se e per quali condizioni può essere una soluzione del nostro problema. Calcolandone il laplaciano vettoriale otteniamo 
(componente $ x $)
\begin{align}
  - \left( \nabla^2 \v{ E } \right)_x = - \frac{n_1\pi}{a} E_x \implies - \left( \nabla^2 \v{ E } \right) 
  = -\left[ \left( \frac{n_1\pi}{a} \right)^2 + \left( \frac{n_2\pi}{b} \right)^2 + \left( \frac{n_3\pi}{c} \right)^2 \right] \v{ E }
\end{align}
dunque una condizione affinché sia soluzione è che
\begin{align}
  \left[ \left( \frac{n_1\pi}{a} \right)^2 + \left( \frac{n_2\pi}{b} \right)^2 + \left( \frac{n_3\pi}{c} \right)^2 \right] = \frac{\omega^2}{c^2} \label{condizione_1}
\end{align}
Applichiamo adesso la condizione di divergenza nulla:
\begin{align}
  0 = \nabla\!\cdot\!\v{ E } &= e^{-i \omega t} C_{pqs} \left( -\frac{n_p \pi}{l_p} \right) \mathcal{E}_p \sin(\frac{n_p \pi}{l_p}x_p) 
  \sin(\frac{n_q \pi}{l_q}x_q) \sin(\frac{n_s \pi}{l_s}x_s)
  \\
  &= - \left( \mathcal{E}_x \frac{n_1 \pi}{a} + \mathcal{E}_y \frac{n_2 \pi}{b} + \mathcal{E}_z \frac{n_3 \pi}{d} \right) e^{-i\omega t} 
  \sin(\frac{n_1 \pi}{a}x) \sin(\frac{n_2 \pi}{b}y) \sin(\frac{n_3 \pi}{d}z) 
  \\
  &\implies \left( \mathcal{E}_x \frac{n_1 \pi}{a} + \mathcal{E}_y \frac{n_2 \pi}{b} + \mathcal{E}_z \frac{n_3 \pi}{d} \right) = 0 \label{condizione_2}
\end{align}
dove $ l_i \in \{ a, b, d\} $ e $ C_{pqs} $ è il tensore ciclico definito come $ C_{pqs} = 1 $ se $ (p,q,s) $ è ciclica, $ 0 $ altrimenti. Vediamo quali 
configurazioni di $ (n_1,n_2,n_3) $ sono possibili in base ai vincoli che abbiamo trovato:
\begin{itemize}
  \item con due o tre parametri nulli ho la soluzione banale
  \item con un solo parametro nullo la soluzione è completamente fissata (quindi anche la polarizzazione)
  \item con tutti i parametri non nulli ho due diverse polarizzazioni possibili
\end{itemize}
Ponendo a $ 0 $ un solo parametro, affinché la frequenza sia minima, in base alle dimensioni che abbiamo dato ($ a>b>d $), occorre annullare $ n_3 $ e in 
questo modo avremo che
\begin{align}
  \v{ E } &= E_z \vers{z} = e^{-i\omega t} \mathcal{E}_z \sin\left( \frac{n_1 \pi}{a} x \right) \sin\left( \frac{n_2 \pi}{b}y \right) \vers{z}
  \\
  \v{ B } &= -\frac{i}{\omega} \nabla\!\times\!\v{ E } 
  \\
  &= - \frac{i}{\omega} e^{-i\omega t} \mathcal{E}_z \left[ \frac{n_1 \pi}{a} \sin\left( \frac{n_1 \pi}{a}x 
  \right) \cos\left( \frac{n_2 \pi}{b}y \right)\vers{x} + \frac{n_2 \pi}{b} \cos\left( \frac{n_1 \pi}{a}x 
  \right) \sin\left( \frac{n_2 \pi}{b}y \right)\vers{y}\right]
\end{align}
di cui una rappresentazione dei profili si può osservare in figura \ref{figure:_cavità_parallelepipedo_campi}.
\begin{figure}[h]
    \center
    \begin{tikzpicture}[ % FISSA IL RIFERIMENTO IN MODO CORRETTO, ALTRIMENTRI --> (X,Z,Y)
        x={(1cm,0cm)},
        y={(0.6cm,0.4cm)},
        z={(0cm,1cm)}
        ]
            % Parametri
        \def\a{4}
        \def\b{3.5}
        \def\d{2}

        % Assi cartesiani
        \draw[->] (0,0,0) -- ({\a + 1},0,0) node[below left] {$x$};
        \draw[->] (0,0,0) -- (0,{\b + 1},0) node[below right] {$y$}; % below/above
        \draw[->] (0,0,0) -- (0,0,{\d + 1}) node[left] {$z$};

        % Vertici del parallelepipedo
        \coordinate (O) at (0,0,0);
        \coordinate (A) at (\a,0,0);
        \coordinate (B) at (\a,\b,0);
        \coordinate (C) at (0,\b,0);

        \coordinate (O1) at (0,0,\d);
        \coordinate (A1) at (\a,0,\d);
        \coordinate (B1) at (\a,\b,\d);
        \coordinate (C1) at (0,\b,\d);

        % Spigoli base
        \draw (O) -- (A) -- (B) -- (C) -- cycle;

        % Spigoli verticali
        \draw (O) -- (O1);
        \draw (A) -- (A1);
        \draw (B) -- (B1);
        \draw (C) -- (C1);

        % Spigoli superiori
        \draw (O1) -- (A1) -- (B1) -- (C1) -- cycle;

        % Etichette dimensioni
        \node at (\a/2, -0.5, 0) {$a$};
        \node at (\a+0.5, \b/2, 0) {$b$};
        \node at (-0.3, 0, \d/2) {$d$};

        % Centro della base
        \coordinate (M) at ({\a/2},{\b/2},0);

        % Ellisse sul piano xy centrata in M
        \draw[dashed] (M) ellipse [x radius={\a * 7/16 }, y radius={\b/4}] node[right] {$B$};

        % Funzione seno sul piano xy: z = 0, y = b/2
        \draw[domain=0:\a, smooth, variable=\x]
            plot ({\x},{\b/2},{ \d * 7/8 *sin(180/\a*\x)}) node[above] {$E$};
        
        \draw[domain=0:\b, smooth, variable=\y]
            plot ({\a/2},\y,{ \d * 7/8 *sin(180/\b *\y)});

    \end{tikzpicture}
    \caption{Campi all'interno della cavità}
    \label{figure:_cavità_parallelepipedo_campi}
\end{figure}

\subsection{Effetto pelle}
Fin ora abbiamo supposto che la cavità sia ideale e che quindi il campo elettrico ai bordi sia nullo, tuttavia quello che si verifica nella realtà è che il 
campo ai bordi è molto piccolo ed entrando nel metallo decresce esponenzialmente. Questo succede perché il campo diverso da zero nel bordo di metallo della 
cavità induce una densità di corrente secondo la legge di Ohm microscopica:
\begin{align}
  \v{ J } = \sigma \v{ E }
\end{align}
con $ \sigma $ la conducibilità del metallo. In questo modo l'equazione da risolvere diventa
\begin{align}
  \nabla\!\times\! (\nabla \! \times \! \v{ E }) = - \nabla^2 \v{ E } = i \omega \nabla\!\times\!\v{ B } 
  = i \omega \left( \mu_0 \v{ J } + \frac{1}{c^2} \der{\v{ E }}{t} \right) 
  = \left( i \omega \mu_0 \sigma + \frac{\omega^2}{c^2} \right) \v{ E }
\end{align}
e definiamo $ \chi $ come
\begin{align}
  \chi^2 \coloneqq \frac{\omega^2}{c^2} + i \omega \mu_0 \sigma 
\end{align}
ma dato che siamo ancora nell'approssimazione di campi lentamente variabili abbiamo $ \omega \ll c $, quindi $ \chi \approx i \omega \mu_0 \sigma $, 
e l'equazione diventa
\begin{align}
  -\nabla^2 \v{ E } = \chi^2 \v{ E }
\end{align}
È facile vedere che una decadenza esponenziale del tipo $ E_{\parallel} \sim \exp(i\chi z) $ risolve una tale equazione, dal momento che il campo può 
essere solo parallelo alla cavità. Per ottenere il campo reale troviamo prima $ \chi $:
\begin{align}
  \chi^2 = i \omega \mu_0 \sigma \implies \chi = \sqrt{\frac{\omega \mu_0 \sigma}{2}} (1+i)
\end{align}
in questo modo possiamo definire una lunghezza caratteristica $ \delta $ chiamata \textit{spessore pelle}
\begin{align}
  \text{SI} \qquad \delta &= \frac{1}{\text{Im}(\chi)} = \sqrt{\frac{2}{\omega \mu_0 \sigma}}
  \\
  \text{cgs} \qquad \delta &= \frac{c}{\sqrt{2\pi \omega \sigma}}
\end{align}
e il campo risulta così
\begin{align}
  E_\parallel = E_0 e^{-z/\delta} \text{Re}(e^{i(z/\delta - \omega t)}) 
  = E_0 e^{-z/\delta} \cos \left( z/\delta -\omega t\right)
\end{align}
Possiamo definire un fattore che ci permette di confrontare vari materiali in base a quanto velocemente decadono i campi in funzione dello spessore del 
metallo:
\begin{align}
  Q \coloneqq \frac{l}{\delta} 
\end{align}
ed è detto \textit{fattore di qualità}. Per gli scopi di confinamento di campi elettromagnetici si impiegano materiali con $ Q \gg 1 $. Il concetto di 
fattore di qualità è molto più generale che per le sole onde elettromagnetiche, infatti fa parte dei parametri per descrivere gli smorzamenti di segnali, 
che siano onde meccaniche o di altra natura.
%% AGGIUNGERE CALCOLO DI LUNGHEZZA PELLE PER ALCUNI MATERIALI

\section{Pacchetti d'onda}
In natura le onde piane sono un fenomeno irrealistico perché si estendono in tutto lo spazio e dovrebbero esistere da sempre. Si può parlare di onde piane 
quando si cerca di approssimare un segnale in certe condizioni, ad esempio le onde elettromagnetiche emesse dal sole possono essere trattate come onde 
piane quando si infrangono sulla Terra. Un modello più realistico consiste nel considerare i \textit{pacchetti d'onda}:
\begin{align}
  f(x) = e^{i k_0 x - \frac{x^2}{L^2}} = e^{i k_0 x} \cdot e^{- \frac{x^2}{L^2}}
\end{align}
che rappresenta una onda piana moltiplicata per una gaussiana. La condizione $ k_0 L \gg 1 $ ci dice che ci sono molte oscillazioni all'interno della 
lunghezza ($ L $) della gaussiana. Questo tipo di onda è molto diverso da una onda piana classica ma si può ottenere da questa tramite il teorema di Fourier 
cioè per sovrapposizione lineare.
% PACCHETTO D'ONDA
\begin{figure}[h]
  \center
  \begin{tikzpicture}
    % parametri
    \def\asse{15}
    \def\A{2}        % ampiezza
    \def\mu{0}       % centro
    \def\sigma{2}    % larghezza

    \def \gaussiana {\A*exp(-((\x-\mu)^2)/(2*\sigma^2))}
    \def \f {sin(500*\x)}
    \draw (-\asse/2, 0) -- (\asse/2, 0);

    \draw [domain = -\asse/2: \asse/2, samples=600] plot (\x , {\gaussiana});
    %\draw [domain = -\asse/2: \asse/2, samples=200] plot (\x , {-\gaussiana});

    \draw [domain = -\asse/2: \asse/2, samples=600] plot (\x , {\f*\gaussiana});

  \end{tikzpicture}
  \caption{Pacchetto d'onda gaussiano}
  \label{figure:_pacchetto d'onda}
\end{figure}


\subsection{Trasformata di Fourier}
Una trasformata è un operatore (in genere lineare) che va da uno spazio di funzioni in un altro spazio di funzioni; la trasformata di Fourier è una 
trasformata integrale e opera con funzioni che appartengono allo spazio $ \mathbb{L}^2 $ (spazio delle funzioni a quadrato sommabile) cioè:
\begin{equation}
  L^2 = \left\{f:\mathbb{R} \to \mathbb{C} \quad \text{tale che} \quad \int_{-\infty}^{+\infty} |f(x)|^2 \,dx < +\infty \right\}
\end{equation}
Definiamo la trasformata di Fourier come segue:
\begin{align}
  f(x) = \frac{1}{\sqrt{2\pi}} \int_{-\infty}^{+\infty} \hat{f}(k) e^{ikx} \,dk
\end{align}
dove $ \mathcal{F}(f) = \hat{f} $ è la trasformata di Fourier della funzione $ f $. Per ottenerla facciamo l'antitrasformata:
\begin{equation}
  \hat{f}(k) = \frac{1}{\sqrt{2\pi}} \int_{-\infty}^{+\infty} f(x) e^{-ikx} \,dx
\end{equation}
È come se fosse una somma (infinita) di funzioni sinusoidali pesate con la $f(k)$.
Qualunque funzione (che sia in $ L^2 $) si può scrivere come sovrapposizione di onde piane monocromatiche. $ \hat{f}(k) $ è l'ampiezza di ciascuna 
componente spaziale, mentre $ \exp(ikx) $ è la componente sinusoidale della frequenza spaziale $ k $. Quando c'è anche dipendenza dal tempo $ f(x,t) $ si ha
\begin{align}
  f(x,t ) = \frac{1}{2\pi} \int \!\!\!\int_{-\infty}^{+\infty} \hat{f}(k,\omega) e^{ikx-i\omega t} \,dkd\omega
\end{align}
e quindi la funzione trasformata vale
\begin{align}
  \hat{f}(k,\omega) = \frac{1}{2\pi}\int\!\!\! \int_{-\infty}^{+\infty} {f}(x,t) e^{-ikx+i\omega t} \,dxdt
\end{align}
La trasformata di Fourier è uno strumento che permette di trasformare equazioni differenziali in equazioni algebriche: applichiamo la trasformata a una 
equazione, troviamo un espressione per la variabile trasformata e infine applichiamo l'anti-trasformata.

\subsection{Trasformata di un oscillatore forzato}
Sfruttiamo questa tecnica nel caso del oscillatore armonico con forzante esterna $ F(t) $:
\begin{align}
  F(t) = \ddot{x} + \omega_0^2 x.
\end{align}
Applico la trasformata di Fourier
\begin{align}
  \hat{F}(\omega) = \frac{1}{\sqrt{2\pi}} \int_{-\infty}^{+\infty} \left( \ddot{x} + \omega_0^2 x \right) e^{-i\omega t} \,dt 
  = (i\omega)^2 \hat{x}(\omega) + \omega_0^2 \hat{x}(\omega)
\end{align}
dove abbiamo sfruttato l'integrazione per parti della $ \ddot{x} $ e che è effettivamente una funzione integrabile. La trasformata della $ x $ quindi vale
\begin{equation}
  \hat{x}(\omega) = \frac{\hat{F}}{\omega_0^2 - \omega^2}
\end{equation}
e infine per trovare la legge oraria ne facciamo l'anti-trasformata
\begin{equation}
  x(t) = \frac{1}{\sqrt{2\pi}} \int_{-\infty}^{+\infty} \frac{\hat{F}}{\omega_0^2 - \omega^2} e^{i\omega t} \,dt
\end{equation}

\subsection{Trasformata dell'equazione delle onde}
La trasformata di Fourier si comporta bene rispetto agli operatori differenziali, infatti abbiamo
\begin{align}
  \mathcal{F}(\partial_t f) = -i \omega \mathcal{F}(f)
  \\
  \mathcal{F}(\partial_x f) = i k \mathcal{F}(f)
\end{align}
le quali sono una diretta conseguenza dell'integrazione per parti e del fatto che le funzioni che andiamo a considerare sono integrabili 
(il loro integrale è finito). Applicando all'equazione delle onde si ottiene:
\begin{align}
  0 = \partial_t^2 \xi - v^2 \partial_x^2 \xi &= \left(- i \omega \right) ^2 \hat{\xi}- v^2 \left( i k  \right)^2 \hat{\xi} 
  = \left(-\omega^2 +v^2k^2 \right) \hat{\xi}
  \\
  \implies \omega &= v\,k \tag{relazione di dispersione}
\end{align}
abbiamo quindi ottenuto in pochi passaggi la relazione di dispersione per le onde nel vuoto. 
\subsection{Trasformata di un pacchetto d'onda}
Consideriamo adesso il pacchetto d'onda gaussiano di prima cioè 
\begin{align}
  f(x,t) = e^{i k_0 x - i\omega t} \cdot e^{- \frac{x^2}{L^2}}
\end{align}
e ne vogliamo fare la trasformata di Fourier, stavolta però stando attenti in che dominio la facciamo. $ f(x,t) $ è funzione dello spazio e del tempo, ma la 
parte temporale è quella di un onda piana, dunque si porta dietro gli stessi aspetti critici cioè che esiste in tutti i tempi. Se facessimo la trasformata 
considerandola funzione di due variabili allora questa non sarebbe ben definita dato che non appartiene a $ \mathbb{L}^2 $ proprio perché il suo integrale non 
converge. Possiamo aggirare questo problema in due modi: annullando completamente la parte temporale, o portarcela dietro e considerarla come una modifica 
alla distribuzione di $ k $, che è ciò che faremo.
\begin{align}
  \hat{f}(k) = \frac{1}{\sqrt{2\pi}} \int_{-\infty}^{+\infty} f(x) e^{-ikx} \,dx =
  \frac{e^{-i\omega t}}{\sqrt{2\pi}} \int_{-\infty}^{+\infty} e^{i(k_0 - k)x - \frac{x^2}{L^2}} \,dx
\end{align}
e per risolvere questo integrale completiamo il quadrato sommando e sottraendo $ \frac{L^2}{4} (k_0 - k)^2 $
\begin{align}
  \hat{f}(k) = \frac{e^{-i\omega t}\cdot e^{-\frac{L^2}{4}(k_0 - k)^2}}{\sqrt{2\pi}} \int_{-\infty}^{+\infty} 
  e^{-\left[ \frac{x}{L} - \frac{iL}{2}(k_0 - k) \right]^2} \,dx = \frac{L\;e^{-i\omega t}\cdot e^{-\frac{L^2}{4}(k_0 - k)^2}}{\sqrt{2}}
\end{align}
Questa funzione rappresenta nel dominio di $ k $ una gaussiana di centro $ \mu = k_0 $ e varianza $ \sigma^2 = \frac{2}{L^2} $, quindi all'aumentare di 
$ L $ la distribuzione si restringe attorno $ k_0 $. Con questi risultati siamo pronti per trattare cosa succede a un pacchetto d'onda che attraversa un 
mezzo materiale: la relazione di dispersione in generale non è lineare, ha dato che la trasformata del pacchetto è centrata nel valore $ k_0 $, è 
ragionevole cercare una nuova legge di dispersione sviluppando in Taylor la frequenza in funzione di del passo d'onda nel punto $ k_0 $. Dallo sviluppo
\begin{align}
  \omega(k) \approx \omega(k_0) + \partial_k \omega (k-k_0) + \frac{1}{2} \partial_k^2 \omega (k-k_0)^2
\end{align}
si definisce la \textit{velocità di gruppo} $ v_g $ come
\begin{align}
  v_g \coloneqq \der{\omega}{k}
\end{align}
a differenza della velocità di fase che è $ v = \frac{\omega}{k} $. La velocità di gruppo è la velocità con cui si muove il centro del pacchetto e
in un mezzo in cui la relazione di dispersione non è lineare, man mano che il pacchetto si sposta allo stesso tempo si disperde in base alle diverse 
velocità di fase delle onde che lo compongono. Tali mezzi sono detti \textit{dispersivi}.


\section{Guide d'onda}
I casi visti precedentemente riguardavano volumi limitati con alcune simmetrie (cilindrica, rettangolare), passiamo ora a considerare dei volumi che si 
estendono in modo (quasi) illimitato lungo una direzione. Tali oggetti si dicono \textit{guide d'onda} poiché vedremo che sono in grado di trasmettere 
segnali elettromagnetici.

Consideriamo un tubo di sezione rettangolare illimitato su $ z $ come in figura \ref{figure:_guida_rettangolare}, dato che non ci sono sorgenti 
dobbiamo risolvere le stesse equazioni delle cavità, ma stavolta con condizioni diverse su $ z $. Supponiamo i campi siano a variabili separabili 
e che abbiano una dipendenza del tipo
\begin{align}
  f(x,y,z,t) = e^{ikz - i\omega t} F(x,y)
\end{align} \begin{figure}
    \center
    \begin{tikzpicture}[ % FISSA IL RIFERIMENTO IN MODO CORRETTO, ALTRIMENTRI --> (X,Z,Y)
        z={(1cm,0cm)},
        x={(0.6cm,0.4cm)},
        y={(0cm,1cm)}
        ]
            % Parametri
        \def\a{2}
        \def\b{3.5}
        \def\d{9}

        % Assi cartesiani
        \draw[->] (0,0,0) -- ({\a + 1},0,0) node[above] {$x$};
        \draw[->] (0,0,0) -- (0,{\b + 1},0) node[below right] {$y$}; % below/above
        \draw[->] (0,0,0) -- (0,0,{\d + 1}) node[below left] {$z$};

        % Vertici del parallelepipedo
        \coordinate (O) at (0,0,0);
        \coordinate (A) at (\a,0,0);
        \coordinate (B) at (\a,\b,0);
        \coordinate (C) at (0,\b,0);

        \coordinate (O1) at (0,0,\d);
        \coordinate (A1) at (\a,0,\d);
        \coordinate (B1) at (\a,\b,\d);
        \coordinate (C1) at (0,\b,\d);

        % Spigoli base
        \draw[dashed] (O) -- (A) -- (B) -- (C) -- cycle;

        % Spigoli verticali
        \draw (O) -- (O1) -- (0,0,{\d -1});
        \draw (A) -- (A1);
        \draw (B) -- (B1);
        \draw (C) -- (C1);

        % Spigoli superiori
        \draw[dashed] (O1) -- (A1) -- (B1) -- (C1) -- cycle;

        % Etichette dimensioni
        % \node at (\a/2, -0.5, 0) {$a$};
        % \node at (\a+0.5, \b/2, 0) {$b$};
        % \node at (-0.3, 0, \d/2) {$d$};

    \end{tikzpicture}
    \caption{Guida rettangolare}
    \label{figure:_guida_rettangolare}
\end{figure}
Riscriviamo il rotore di rotore di $ \v{ E } $:
\begin{align}
  -\nabla^2 \v{ E } = \nabla\!\times\!\left( \nabla\!\times\!\v{ E } \right) = \nabla \! \times \! \left( i \omega \v{ B } \right) 
  =\frac{\omega^2}{c^2} \v{ E } 
\end{align}
e iniziamo a calcolare il laplaciano della $ f(x,y,z,t) $:
\begin{align}
  -\left( \der[2]{}{x} + \der[2]{}{y} + \der[2]{}{z} \right)f &= \left( k^2 - \der[2]{}{x} - \der[2]{}{y} \right) f
  \\
  &= \left( k^2 - \nabla_{xy}^2 \right) e^{ikz - i\omega t} F(x,y)
\end{align}
dove $ \nabla_{xy}^2 $ è il laplaciano $ 2 $-dimensionale rispetto a $ x $ e $ y $. Quindi l'equazione che dobbiamo risolvere è
\begin{align}
  \left( k^2 - \nabla_{xy}^2 \right) F(x,y) = \frac{\omega^2}{c^2} F(x,y)
\end{align}
che è una equazione agli autovalori per l'operatore differenziale $ \left( k^2 - \nabla_{xy}^2 \right) $, che introducendo 
$ \chi^2 = \frac{\omega^2}{c^2} - k^2 $ si ha
\begin{align}
  - \nabla_{xy}^2 F(x,y) = \chi^2 F(x,y)
\end{align}
Con la definizione che abbiamo dato di $ \chi^2 $ abbiamo che la relazione di dispersione rappresenta una iperbole nel piano $ (\omega/c , k) $ dato che
\begin{align}
  \frac{\omega}{c} = \sqrt{k^2 + \chi^2}
\end{align}
con asintoto la cosiddetta \textit{linea luce} $ \frac{\omega}{c} = k $, come possiamo vedere nel grafico \ref{figure:_linea_luce}. È chiamata linea luce 
perché corrisponde alla situazione di onda piana nel vuoto, quindi senza vincoli traversali ed è interessante notare che per $ k \to \infty $ si ha che 
ogni relazione di dispersione tende a quella del vuoto, questo è come se la guida "non si accorgesse" dei bordi, dato che con un passo d'onda molto grande 
la sua evoluzione spaziale è "rallentata". \begin{figure}
    \center
    \begin{tikzpicture}
        
        \draw[->] (0,0) -- (0,4) node[left] {$ \frac{\omega}{c} $};
        \draw[->] (0,0) -- (5,0) node[below] {$ k $};

        \draw[dashed, domain = 0:5, variable = \x] plot ({\x}, {\x}) node[below right, scale = 0.9] {$ \frac{\omega}{c} = k $};
        % parametro

        \def\a{1}
        \def\b{2}
        \def\c{0.5}        

        % iperbole y = sqrt(x^2 + a^2)
        \draw[domain=0:5,smooth,variable=\x,orange,thick]
            plot ({\x},{sqrt(\x*\x+\a*\a)});
        
        \node[ left,scale = 0.7, black] at (0,1) {$ \chi^2 = 1 $};

        \draw[domain=0:5,smooth,variable=\x,blue,thick]
            plot ({\x},{sqrt(\x*\x+\b*\b)});
        
        \node[ left,scale = 0.7, black] at (0,2) {$ \chi^2 = 4 $};
        
        \draw[domain=0:5,smooth,variable=\x,red,thick]
            plot ({\x},{sqrt(\x*\x+\c*\c)});
        
        \node[ left,scale = 0.7, black] at (0,0.5) {$ \chi^2 = 0.25 $};

    \end{tikzpicture}
    \caption{Grafico delle relazioni di dispersione per vari valori di $ \chi^2 $.}
    \label{figure:_linea_luce}
\end{figure}

Scriviamo le equazioni di Maxwell applicate al caso della guida senza sorgenti:
\begin{align}
  &\nabla \times \v{E} = -\der{\v{B}}{t} = i\omega \v{B} \qquad &\nabla \cdot \v{E} = 0
  \\
  &\nabla \times \v{B} = \frac{1}{c^2} \der{\v{E}}{t} = -\frac{i \omega}{c^2} \v{E} \qquad &\nabla \cdot \v{B} = 0
\end{align}
che esplicitandole direttamente diventano
\begin{align}
  i\omega B_x &= \partial_y E_z - \partial_z E_y  &{-\frac{i \omega}{c^2} E_x = \partial_y B_z - \partial_z B_y}
  \\
  i\omega B_y &= \partial_z E_x - \partial_x E_z  &{-\frac{i \omega}{c^2} E_y = \partial_z B_x - \partial_x B_z}
  \\
  i\omega B_z &= \partial_x E_y - \partial_y E_x  &{-\frac{i \omega}{c^2} E_z = \partial_x B_y - \partial_y B_x}
\end{align}
Adesso che abbiamo esplicitato i due rotori possiamo classificare le varie soluzioni secondo i modi TEM, TE, TM, che stanno per 
\textit{trasversal electro-magnetic}. Ognuno di questi modi consiste nell'imporre il campo associato trasversale rispetto alla direzione di propagazione, 
che in questo caso è $ z $. Il modo TEM impone che 
\begin{align}
  E_z = B_z = 0
\end{align}
quindi le equazioni che otteniamo sono
\begin{align}
  i\omega B_x &= - ik E_y & -\frac{i \omega}{c^2} E_x &= - ik B_y
  \\
  i\omega B_y &= ik E_x &  -\frac{i \omega}{c^2} E_y &= ik B_x 
  \\
  0 &= \partial_x E_y - \partial_y E_x & 0 &= \partial_x B_y - \partial_y B_x
\end{align}
dove abbiamo svolto le derivate lungo $ z $ delle varie componenti. Dalle equazioni notiamo che i campi rispettano le relazioni
\begin{align}
  \left( \nabla \times \v{E} \right)_z = 0 \qquad \v{B} = \frac{1}{c} \vers{z} \times \v{E}
\end{align}
cioè i campi rispettano le condizioni dell'elettrostatica, e con il fatto che il campo elettrico è conservativo abbiamo da risolvere l'equazione di 
Laplace per il potenziale $ \nabla^2 V = 0 $, con le condizioni al bordo del conduttore cioè $ V(x,y)|_{\text{bordo}} = V_0 $, ma una soluzione è proprio 
il potenziale costante in ogni punto, che per il teorema di unicità è unica. Ma se il potenziale è costante i campi sono di conseguenza banali, dunque 
concludiamo che in queste condizioni il modo TEM è solo il modo banale. Il ragionamento appena fatto vale pero nel caso in cui il dominio sia 
semplicemente connesso, infatti nel caso del cavo coassiale il dominio non è connesso e infatti esiste il modo TEM come soluzione.

% POTENZA
{Per quanto riguarda la potenza stiamo passando da una equazione differenziale nel tempo a una che è nello spazio: 
scriviamo il bilancio energetico come potenza entrante meno la potenza uscente: potenza entrante (il vettore di 
Poynting mediato sul tempo e integrato sulla superficie) che sarà proporzionale a $ \Delta z $ portiamo sotto e 
passiamo a $ dz $.}


\section{Soluzione equazioni di Maxwell in presenza di sorgenti} \label{section:_soluzione_Maxwell_con_sorgenti}
Fin ora abbiamo trattato i casi in cui le sorgenti non sono presenti, dunque le soluzioni che abbiamo trovato sono il limite in cui i campi sono molto 
lontani da ciò che li ha generati. Vogliamo ora rielaborare le equazioni di Maxwell in termini dei potenziali elettrico $ V $ e vettore $ \v{ A } $ 
in modo da ottenere una forma integrale per i campi.

Dato che $ \nabla \cdot \v{ B } = 0 $ è valida sia nel caso statico sia nel caso dinamico, la definizione che avevamo dato del potenziale vettore è 
ancora valida cioè
\begin{align}
  \nabla\!\times\!\v{ A } = \v{ B }.
\end{align}
Ciò non vale per il potenziale scalare $ V $ perché il campo elettrico non è più conservativo, serve rielaborare l'espressione del suo rotore:
\begin{align}
  \nabla\!\times\!\v{ E } &= -\der{\v{ B }}{t} = - \der{}{t} \left( \nabla\!\times\!\v{ A } \right) 
  = - \nabla \times \left( \der{\v{ A }}{t} \right) 
  \\
  &\implies \nabla \times \left( \v{ E } + \der{\v{ A }}{t} \right) = 0
\end{align}
dunque definiamo un nuovo potenziale $ V $ tale che il suo gradiente (cambiato di segno) sia $ \v{ E } + \der{\v{ A }}{t} $:
\begin{align}
  -\nabla V = \v{ E } + \der{\v{ A }}{t} \implies \v{ E } = - \nabla V - \der{\v{ A }}{t}.
\end{align}
Adesso che abbiamo sistemato i due potenziali passiamo a riscrivere le due equazioni non omogenee in termini di $ V $ e $ \v{ A } $:
\begin{align}
  \left.1\right)\qquad &\nabla\!\cdot\!\v{ E } = - \nabla^2 V - \der{}{t} \left( \nabla\!\cdot\!\v{ A } \right) = \frac{\rho}{\eps_0}
  \\
  \left.2\right)\qquad &\nabla\!\times\!\v{ B } = \nabla\!\times\!\left( \nabla\!\times\!\v{ A } \right) 
  = \nabla \left( \nabla\!\cdot\!\v{ A } \right) - \nabla^2 \v{ A } = \mu_0 \v{ J } + \frac{1}{c^2} \der{\v{ E }}{t} 
  \\
  &= \mu_0 \v{ J } - \frac{1}{c^2}\left[ \nabla \left( \der{V}{t} \right) + \der[2]{\v{ A }}{t} \right].
\end{align}
Le equazioni che abbiamo ottenuto sono complicate da risolvere così come sono, perciò per semplificarle sfruttiamo il fatto che il potenziale vettore 
non è unico: l'idea è di applicare una trasformazione di gauge per ottenere un nuovo potenziale vettore che ci semplifichi le equazioni. Ricordiamo 
che una trasformazione di gauge è della forma
\begin{align}
  \nabla \cdot \v{ A }' = 0, \qquad \v{ A }' = \v{ A } + \nabla \Psi,
\end{align}
ma trasformare il potenziale vettore adesso influenza anche quello scalare, dato che il campo elettrico deve rimanere lo stesso prima e dopo la 
trasformazione:
\begin{align}
  \v{ E }' = -\nabla V' - \der{\v{ A }'}{t} &= -\nabla V - \der{\v{ A }}{t} - \der{}{t} \left( \nabla \Psi \right) \overset{!}{=} \v{ E }
  \\
  \implies V' &= V - \der{\Psi}{t} 
\end{align}
In elettrostatica si preferisce lavorare nella cosiddetta \textit{gauge di Coulomb}, ovvero con $ \nabla\!\cdot\!\v{ A' } = 0 $ e $ V $ che rispetta 
l'equazione di Poisson 
\begin{align}
  \nabla^2 V = -\frac{\rho}{\eps_0} \implies V = \frac{1}{4\pi\eps_0} \int \frac{\rho(\v{ x }', t)}{|\v{ x }-\v{ x }'} \,d^3x'
\end{align}
mentre adesso sfrutteremo la \textit{gauge di Lorenz} (fisico danese, da non confondere con il fisico olandese Lorentz a cui si devono le omonime trasformazioni e 
la forza magnetica) dove il nuovo potenziale vettore $ \v{ A }' $ è tale che
\begin{align}
  \nabla\!\cdot\!\v{ A }' = -\frac{1}{c^2} \der{V'}{t}
\end{align}
che si ottiene imponendo che $ \Psi $ risolva l'equazione
\begin{align}
  \nabla\!\cdot\!\v{ A } + \nabla^2 \Psi &= -\frac{1}{c^2} \left( \der{V}{t} + \der[2]{\Psi}{t} \right)
  \\
  \implies \Box\Psi &= - \left( \nabla\!\cdot\!\v{ A } + \frac{1}{c^2} \der{V}{t} \right)
\end{align}
dove $ \Box \Psi $ è l'operatore di d'Alembert applicato alla funzione scalare $ \Psi $. Questa particolare gauge sarà molto utile per scrivere le 
equazioni di Maxwell in forma covariante poiché è invariante sotto trasformazioni fra sistemi di riferimento inerziali. Dopo questa trasformazione di Gauge 
possiamo riscrivere le equazioni (per semplicità di notazione indicheremo i nuovi potenziali semplicemente con $ \v{ A } $ e $ V $ poiché non 
compariranno termini contenenti i vecchi potenziali), partiamo dalla prima:
\begin{align}
 \left.1\right)\qquad &-\nabla^2 V - \der{}{t}\left( \nabla\!\cdot\!\v{ A } \right) = -\nabla^2 V + \frac{1}{c^2} \der[2]{V}{t} = \frac{\rho}{\eps_0}.
\end{align}
La seconda:
\begin{align}
  \left.2\right)\qquad &\nabla\left( \nabla\!\cdot\!\v{ A } \right) -\nabla^2 \v{ A } 
  = \mu_0 \v{ J } + \nabla \left( \nabla\!\cdot\!\v{ A } \right) -\frac{1}{c^2}\der[2]{\v{ A }}{t}
\end{align}
semplificando i termini otteniamo
\begin{align}
  \left.1\right)\qquad &-\nabla^2V + \frac{1}{c^2} \der[2]{V}{t} =-\Box V = \frac{\rho}{\eps_0}
  \\
  \left.2\right)\qquad &-\nabla^2\v{ A } + \frac{1}{c^2} \der[2]{\v{ A }}{t} =-\Box \v{ A } = \mu_0\v{ J }
  \\
  \left.3\right)\qquad &\nabla\!\cdot\!\v{ A } + \frac{1}{c^2} \der{V}{t} = 0 \label{eq:_condizione_di_Lorenz}
\end{align}
che sono delle equazioni più semplici da risolvere. Dunque abbiamo trovato che $ V $ e ogni componente $ A_i $ del potenziale vettore soddisfano 
l'equazione non omogenea di d'Alembert, ci basta quindi trovare una soluzione per una generica sorgente $ S(\v{ x }, t) $:
\begin{align}
  \nabla^2 f(\v{ x }, t) - \frac{1}{c^2} \der[2]{}{t} f(\v{ x },t) = -S(\v{ x },t).
\end{align}
Dato che è lineare, applichiamo a tutta l'equazione la trasformata di Fourier rispetto al tempo, in modo che dopo ricostruiremo la funzione originale 
$ f(\v{ x }, t) $ facendo l'anti-trasformata:
\begin{align}
  &f(\v{ x },t) = \int_{-\infty}^{+\infty} e^{-i\omega t} f_\omega(\v{ x }) \,d\omega
  \\
  &\implies \nabla^2 f_\omega + \frac{\omega^2}{c^2} f_\omega = -S_\omega
\end{align}
dove abbiamo sfruttato le proprietà della trasformata di Fourier con le derivate. Definendo $ k^2 \coloneqq \omega^2 /c^2 $ otteniamo quella che nota 
come \textit{equazione di Helmholtz}:
\begin{align}
  \left( \nabla^2 + k^2 \right)f_\omega = -S_\omega
\end{align}
Per trovarne una soluzione sfruttiamo la tecnica di Green, cioè scriviamo la funzione $ f_\omega(\v{ x }) $ come
\begin{align}
  f_\omega(\v{ x }) = \frac{1}{4\pi} \int S_\omega(\v{ x }') G(\v{ x }, \v{ x }') \,d^3x'
\end{align}
e adesso troviamo quale condizione deve rispettare la funzione di Green $ G(\v{ x },\v{ x }') $ affinché $ f_\omega $ sia soluzione:
\begin{align}
  \left( \nabla^2 + k^2 \right)f_\omega 
  = \frac{1}{4\pi} \int S_\omega(\v{ x }') \left[ \nabla^2 + k^2 \right]\!G(\v{ x }, \v{ x }') \,d^3x' \overset{!}{=} -S_\omega(\v{ x })
\end{align}
da cui la condizione da imporre è
\begin{align}
  \left( \nabla^2 + k^2 \right)G(\v{ x },\v{ x }') = -4\pi \delta(\v{ x }-\v{ x }')
\end{align}
È noto che la funzione di Green che risolve l'equazione di Helmholtz è 
\begin{align}
  G(r) = \frac{e^{ikr}}{r} \quad \text{con} \quad r = |\v{ x }-\v{ x }'|
\end{align}
espressa in coordinate sferiche (per verificare che risolve effettivamente l'equazione si usi il laplaciano in coordinate sferiche). Affinché sia proprio 
la soluzione che cercavamo dobbiamo verificare che risolva anche la \eqref{eq:_condizione_di_Lorenz}: 
\begin{align}
  \nabla\!\cdot\!\v{ A }_\omega &= k_m \int \v{ J }_\omega(\v{ x' }) \cdot \nabla G(\v{ x },\v{ x }') \,d^3x' 
  = - k_m \int \v{ J }_\omega(\v{ x' }) \cdot \nabla' G(\v{ x },\v{ x }') \,d^3x' 
  \\
  &= - k_m \int \left[ \nabla' \cdot \left( \v{ J }_\omega(\v{ x' }) G(\v{ x },\v{ x }') \right) 
  - G(\v{ x },\v{ x }') \nabla' \cdot \v{ J }_\omega(\v{ x' }) \right] \,d^3x' 
  \\
  &= 0 + k_m \int i\omega \rho_\omega(\v{ x }') G(\v{ x }, \v{ x }') \,d^3x' 
  = \frac{i\omega}{4\pi c^2} \int  \frac{\rho_\omega(\v{ x }')}{\eps_0} G(\v{ x }, \v{ x }') \,d^3x' = \frac{-1}{c^2} \der{V_\omega}{t}
\end{align}
dove abbiamo usato (in ordine): la differenziazione di un prodotto; il teorema della divergenza per fare diventare un termine un integrale di superficie 
e annullarlo grazie all'ipotesi di sorgenti localizzate; l'equazione di continuità con la proprietà di derivazione della trasformata di Fourier. 
Finalmente abbiamo la nostra soluzione, possiamo quindi applicare l'anti-trasformata per ottenere i potenziali:
\begin{align} 
  V(\v{ x },t) &= \int_{-\infty}^{+\infty} e^{-i\omega t}V_\omega(\v{ x }) \,d\omega 
  = \frac{1}{4\pi} \int_{-\infty}^{+\infty} e^{-i\omega t} \int \frac{\rho_\omega(\v{ x }')}{\eps_0} 
  \frac{e^{i\frac{\omega}{c} |\v{ x-x' }|}}{|\v{ x - x' }|} \,d^3x'\,d\omega
  \\
  &= \frac{1}{4\pi \eps_0} \int \frac{1}{|\v{ x-x' }|} \int_{-\infty}^{+\infty} \rho_\omega(\v{ x' }) 
  \exp\left\{ -i\omega \left( t - \frac{|\v{ x-x' }|}{c} \right) \right\} \,d\omega \,d^3x' 
  \\
  &= k_e \int \frac{\rho\left( \v{ x }', t - \frac{|\v{ x-x' }|}{c} \right)}{|\v{ x-x' }|} \,d^3x'
  \label{eq:_potenziale_scalare_ritardato}
\end{align}
gli stessi identici passaggi valgono per il potenziale vettore
\begin{align} \label{eq:_potenziale_vettore_ritardato}
  \v{ A }(\v{ x },t) = k_m  \int\frac{\v{J}\left(\v{x}',t-\frac{|\v{x-x'}|}{c} \right)}{|\v{x-x'}|} \,d^3x'
\end{align}
Quelli che abbiamo ottenuto sono i cosiddetti \textit{potenziali ritardati}, per evidenziare il fatto che il potenziale in un punto dello spazio-tempo
non dipenda dalla sorgente in quell'istante, ma in un intervallo di tempo precedente detto tempo ritardato $ t_r $, che corrisponde al tempo che la luce 
impiega a raggiungere il punto. Detto in termini relativistici, il potenziale in un evento dipende dalla sorgente valutata sul cono di luce passato di 
quell'evento, con il tempo ritardato che vale
\begin{align}
  t_r = t - \frac{|\v{ x-x' }|}{c}
\end{align}


\begin{comment}

\section{Elettrodinamica raw del 23 febbraio}
Fattore di qualità $ Q \sim \frac{l}{\delta} \approx 10^4 \gg 1 $(nei materiali) con $ l $ spessore della cavità e $ \delta $ spessore pelle.
\begin{align}
  \delta \sim \frac{c}{\sqrt{2\pi\omega \sigma}}.
\end{align}
Velocità di gruppo: $ v_g \coloneqq = \der{\omega}{k} $ che si ottiene dal primo termine dello sviluppo in Taylor.
L'effetto degli ordini superiori è quello di allargare il pacchetto mentre si propaga.

Consideriamo un tubo di sezione rettangolare illimitato su $ z $. Questo oggetto è detto guida d'onda, e abbiamo le stesse condizioni al bordo nel caso 
della cavità rettangolare. Supponiamo che i nostri campi abbiano una dipendenza da del tipo
\begin{align}
  f(x,y,z,t) = e^{ikz} e^{-i\omega t} F(x,y)
\end{align}
Applichiamo il laplaciano 
\begin{align}
  (\der[2]{}{z} + \dots ) f = \left( k^2 - \nabla^2_{2D} \right) F(x,y) e^{ikz - i\omega t}
\end{align}
che è un'altra equazione agli autovalori per un operatore differenziale che stavolta è $ k^2 - \nabla^2_{2D} $. 
\begin{align}
  -\nabla^2_{2D} \sim \frac{1}{l^2} \eqcolon \chi^2 \implies \frac{\omega^2}{c^2} = k^2 + \chi^2
\end{align}
\% GRAFICO CI SI ALLINEA ALLA LINEA LUCE $ (k, \frac{\omega}{c}) $. Per $ \frac{k^2}{\chi^2} $ piccoli si ha una dipendenza parabolica:
\begin{align}
  \alpha + \beta K^2 = \frac{\omega}{c} = \chi \sqrt{1- \frac{k^2}{\chi^2}}
\end{align}
Praticamente $ \chi^2 $ è il risultato che troveremo in aggiunta a $ k^2 $ nell'autovalore dell'operatore.

Risolviamo le stesse due equazioni:
\begin{align}
  \nabla\!\times\!\v{ E } = i \frac{\omega}{c} \v{ B } \qquad \nabla\!\times\!\v{ B } = -i \frac{\omega}{c} \v{ E }
  \\
  i \frac{\omega}{c} B_x = \der{E_z}{y} - ik \der{E_y}{z}
  \\
  \dots
\end{align}
stessa cosa per $ E $. Iniziamo a cercare delle soluzioni tali che $ E_z = B_z = 0 $ che si chiamano TED (trasverse electrical and magnetic) e otterremo 
che $ B = \vers{z} \times \v{ E } $ (come nelle onde piane).
\begin{align}
  i \frac{\omega}{c} B_x = -ik \der{B_y}{z} = i k ( k \frac{c}{\omega}) B_x 
\end{align}
($ \frac{\omega}{c} = k $)
\begin{align}
  B_z = 0 \implies \left( \nabla\!\times\!\v{ E } \right)_z = 0
  \\
  E_z = 0 \implies 0 = \der{E_x}{x} + \der{E_y}{y} = \nabla\!\cdot\!\v{ E } = 0
\end{align}
\begin{align}
  \v{ E } = -\nabla V(x,y) \qquad \nabla^2_{2D} V = 0
\end{align}
dobbiamo imporre che $ V(x,y) $ al bordo sia costante così da avere campo nullo agli estremi, ma quindi questo significa che $ V $ deve essere costante 
ovunque quindi campo nullo, a meno che il bordo non sia disconnesso: cavo coassiale. In una guida d'onda di sezione arbitraria questi modi (TEM) non 
esistono. Proviamo altri modi meno forti: TE $ E_z = 0, B_z \neq 0 $ e TM $ E_z \neq 0, B_z = 0 $. Iniziamo dal TE. (tutto in funzione di $ E_z $)
\begin{align}
  B_x = -i \frac{\omega}{c} \der{E_z}{y} \frac{1}{\frac{\omega^2}{c^2}- k^2}
  \\
  B_y = i \frac{\omega}{c}  \der{E_z}{x} \frac{1}{\frac{\omega^2}{c^2}- k^2}
  \\
  E_x = \frac{c}{\omega} k B_y = ik \der{E_z}{x} \frac{1}{\frac{\omega^2}{c^2}- k^2}
  \\
  E_y = - \frac{c}{\omega} k B_x = i k \der{E_z}{y} \frac{1}{\frac{\omega^2}{c^2}- k^2}
\end{align}
$ (E_x, E_y) \sim \nabla_{2D}E_z $ gradiente bidimensionale di $ E_z $, e $ \v{ B } \sim \vers{z} \times \nabla_{2D} E_z $
\begin{align}
  \der{}{x} \der{E_z}{y} - \der{}{y} \der{E_z}{x} = 0
  \\
  -i \frac{\omega}{c} E_z = \left( \der[2]{E_z}{x} + \der[2]{E_z}{y} \right) \frac{i\omega c}{\frac{\omega^2}{c^2}- k^2}
  \\
  \implies \nabla^2_{2D} E_z (x,y) = \left( {k^2 -\frac{\omega^2}{c^2}} \right) E_z (x,y)
\end{align}
dobbiamo imporre che $ E_z $ sul bordo è nullo e si ottiene un insieme di modi. Sul bordo abbiamo che il gradiente è ortogonale, $ E $ può essere solo 
ortogonale, ma quindi $ B $ può essere solo tangente alla superficie dal fatto che $ \v{ B } \sim \vers{z} \times \nabla_{2D} E_z $.
Per il modo TE ce li facciamo noi i conti: scrivere tutto in funzione di $ B_z $ e troviamo che $ (B_x, B_y) \sim \nabla_{2D}B_z $ e che
$ \v{ E } \sim \vers{z} \times \nabla_{2D} B_z $, otteniamo $ \nabla^2_{2D} E_z (x,y) = \left( {k^2 -\frac{\omega^2}{c^2}} \right) E_z (x,y )$
imponiamo
\begin{align}
  \der{B_z}{n}|_{\text{bordo}} = \nabla_{2D} B_z \cdot \vers{n}
\end{align}
Per sovrapposizione fra modi (TM e TE) si può ottenere campi che hanno tutte le componenti non nulle.

\subsection{Guida rettangolare}
\begin{align}
  E_z (x,y) =  \sin(\frac{n_1 \pi}{a}x) \sin(\frac{n_2 \pi}{b}y) e^{i(kz-\omega t)}
\end{align}
otteniamo la condizione
\begin{align}
  \frac{\omega^2}{c^2} = k^2 + \pi^2 \left( \frac{n^2_1}{a^2} + \frac{n^2_2}{b^2} \right)
\end{align}
e devono essere entrambi non nulli senno il termine esponenziale muore. Per TE: 
\begin{align}
  B_z (x,y) =  \cos(\frac{n_1 \pi}{a}x) \cos(\frac{n_2 \pi}{b}y) e^{i(kz-\omega t)}
\end{align}
stessa relazione (stessa equazione ma diverse C.B.). Cambia la frequenza minima (se $ a>b $ mettiamo a $ n_1 $)
\begin{align}
  TE_{10}: \quad \frac{\omega^2}{c^2} = k^2 + \frac{\pi^2}{b^2}
\end{align}
Una scorciatoia è usare la relazione $ E = \vers{z} \times $
\begin{align}
  B_z = \cos \left( \frac{\pi}{a}x \right) e^{i(kz-\omega t)}
  \\
  \nabla^2_{2D} B_z \propto \vers{x} \sin\left( \frac{\pi}{a} x \right)e^{i(kz-\omega t)}
\end{align}
\% GRAFICO DELL'ANDAMENTO DELLE SOLUZIONI
\begin{align}
  \v{ E } (x,y) = \vers{y} a \sin(\frac{\pi}{a}x) e^{i(kz-\omega t)}
\end{align}
con Maxwell ci troviamo il campo magnetico
\begin{align}
  \nabla\!\times\!\v{ E } = i\frac{\omega}{c} \v{ B }
\end{align}
Ora abbiamo tutto. Media temporale dell'energia:
\begin{align}
  \langle U \rangle = \frac{1}{8\pi} \left( \langle E^2 \rangle + \langle B^2 \rangle \right) 
  = \frac{a^2}{16\pi} \left( \sin^2(\frac{\pi}{a}x) + \left( \frac{ck}{\omega} \right)\sin^2(\pi/a x) + (\frac{c\pi}{\omega a})\cos^2(\pi/a x) \right)
\end{align}
è opportuno fare una ulteriore media stavolta sulla sezione, così i cos2 e i sen2 mediano $ 1/2 $
\begin{align}
  \langle \langle U \rangle \rangle = \frac{1}{8\pi} \left( \langle \langle E^2 \rangle \rangle + \langle \langle B^2 \rangle \rangle \right) = \frac{a^2}{16\pi}
\end{align}
\begin{align}
  \langle S_z \rangle = \frac{c}{4\pi} \langle- E_y B_x \rangle = \frac{c}{8\pi} a^2 \sin^2 (\frac{\pi}{a}x)ck/\omega
\end{align}
notiamo che per fortuna sono in fase (sennò la media era nulla) 
\begin{align}
  \langle \langle U \rangle \rangle = \frac{a^2 c^2 k}{16\pi \omega} \overset{!}{=} v_g \langle \langle u \rangle \rangle
\end{align}
ottengo 
\begin{align}
  v_g = \frac{c^2k}{\omega} \overset{?}{=} \der{\omega}{k}
  \\
  \der{}{k}\left( disp \right) = 
\end{align}
la relazione della doppia media del S e u non è locale, vale solo una volta integrato.
Quanto vale la pressione sulle pareti della guida? Fisica o tensore di Maxwell.
Sfera pulsante, E costante fuori dal raggio massimo, B sempre nullo a causa della simmetria sferica, quindi niente onde E-M, e niente irraggiamento di 
monopolo. $ \der{B}{t} $ associata a una j quando il piano passa, delta


\end{comment}
\section{Irraggiamento}
\begin{figure}
    \centering
    \begin{tikzpicture}
        \def \xa {4}
        \def \ya {1}
        \def \p {0.15}
        \pgfmathsetmacro{\angolo}{atan2(\ya,\xa) + 90}
        \pgfmathsetmacro{\m}{atan2(\ya,\xa)}
        \pgfmathsetmacro{\l}{7}

        \draw[thick, fill=gray!10]
            plot[smooth cycle, tension=0.9] coordinates {
            (1.2,0.1)
            (0.6,1.1)
            (-0.7,1)
            (-1.1,0.2)
            (-0.6,-0.9)
            (0.9,-0.7)
            };

        
        \draw (0,0) -- (\xa,\ya);


        \node[rotate=\angolo] at ({\xa + \p*cos(\m)},{\ya + \p*sin(\m)}) {$\sim$};
        \node[rotate=\angolo] at ({\xa + 2*\p*cos(\m)},{\ya + 2*\p*sin(\m)}) {$\sim$};
        \node[rotate=\angolo] at ({\xa + 3*\p*cos(\m)},{\ya + 3*\p*sin(\m)}) {$\sim$};

        \draw[->] ({\xa + 4*\p*cos(\m)},{\ya + 4*\p*sin(\m)}) -- ({\l*cos(\m)}, {\l*sin(\m)});
        \draw[->] (0,0) -- (0.3,0.8);
        % \draw (0.3,0.8) -- ({\xa}, {\ya + 0.3});
        % \draw[->] ({\xa + 4*\p*cos(\m)},{\ya + 0.3 + 4*\p*sin(\m)}) -- ({\l*cos(\m)}, {\l*sin(\m)});
    \end{tikzpicture}
    \caption{Sorgenti}
    \label{figure:_irraggiamento}
\end{figure}
Con le espressioni che abbiamo ottenuto per i potenziali ritardati vogliamo studiare il problema dell'irraggiamento: valutare in un punto lontano gli 
effetti elettromagnetici di sorgenti localizzate ed elettricamente neutre (altrimenti il primo termine sarebbe quello di monopolo). Quanto lontano è 
stabilito dalla condizione di \textit{far fields}: preso $ \v{ r } = R \vers{u} $ diciamo che siamo in far fields quando $ R \gg a $ e 
$ R \gg \lambda $ dove $ a $ è la dimensione (maggiore) delle sorgenti e $ \lambda $ è la lunghezza d'onda delle onde elettromagnetiche che otterremo, 
ovviamente da verificare a posteriori.

Consideriamo delle sorgenti $ \v{ J }(\v{ r },t) , \rho(\v{ r },t) $ con dipendenza temporale $ e^{-i\omega t} $, il punto in cui vogliamo 
valutare i campi è $ \v{ r } $ mentre la distanza fra il punto e un elemento della sorgente è 
\begin{align}
  |\v{ r-r' }| \simeq R - \v{ r' } \cdot \vers{u} = R - \frac{\v{ r'\cdot r }}{R}
\end{align}
passiamo ai potenziali ritardati nel dominio delle frequenze facendo una trasformata di Fourier
\begin{align}
  \v{ A }_\omega (\v{ r }) &= \mathcal{F}(\v{ A }(\v{ r },t)) 
  = k_m \int d^3r' \int_{-\infty}^{+\infty} dt \frac{\v{ J }(\v{ r' }) e^{-i\omega t_r}}{|\v{r-r'}|} e^{i\omega t} \,d^3r' 
  \\
  &= k_m \int  \frac{e^{i\frac\omega{c} |\v{r - r'}|}}{|r - r'} \left( \int \v{J(\v{r'}, t)} e^{i\omega t} \,dt \right) \, d^3 r'
  = k_m \int  \frac{\v{ J }_\omega(\v{ r' })}{|\v{r-r'}|} e^{i\frac{\omega}{c} |\v{r - r'}| } \,d^3r' 
  \\
  &\simeq k_m \int  \frac{\v{ J }_\omega(\v{ r' })}{R - \frac{\v{r\cdot r'}}{R}}\; e^{i k (R - \frac{\v{r\cdot r'}}{R})} \,d^3r' 
  \simeq k_m\; \frac{e^{ik R}}{R} \int \v{J}_\omega(\v{r'}) \; e^{-ik \frac{ \v{r\cdot r'} }{ R }} \,d^3r'
\end{align}
e da questo risultato vediamo che il potenziale vettore rappresenta un onda sferica che decresce come $ R^{-1} $, con l'aggiunta di un qualche 
coefficiente dovuto all'angolo a cui si trova il punto di interesse $ \v{r} $. A questo punto è chiaro che nell'espressione ottenuta possiamo sviluppare 
in serie di Taylor l'esponenziale e ottenere una serie di multipoli, ma affinché questa approssimazione sia valida occorre aggiungere un ulteriore ipotesi 
che ci permette di dire che l'esponenziale è abbastanza piccolo da essere sviluppato. L'esponente è $ k (\vers{u} \cdot \v{r'}) $ ma $ r' $ è al massimo $ a $
quindi dobbiamo imporre $ k\, a = 2\pi a/\lambda \ll 1 $ che è la condizione di \textit{sorgenti non relativiste} (anche chiamata approssimazione di 
dipolo come vedremo tra poco), e il motivo è che
\begin{align}
  ka = 2\pi \frac{a}{\lambda} = \frac{2\pi a}{T \,c} = 2\pi \frac{v}{c} \ll 1 \implies v \ll c.
\end{align}
Sotto questa ipotesi aggiuntiva scriviamo l'esponenziale in serie:
\begin{align}
  \v{A}_\omega(\v{r}) = k_m \frac{e^{ikR}}{R} 
  \sum_{n = 0}^\infty \frac{\left( -ik \right)^n}{n!} \int \v{J}_\omega(\v{r'}) \left( \vers{u}\cdot \v{r'} \right)^n \,d^3r'
\end{align}
e noi siamo interessati al termine che da il maggior contributo quindi per $ n = 0 $. L'idea per svolgere questo integrale è quella di ricondurci alla 
divergenza della densità di corrente e quindi usare l'equazione di continuità per trasformarlo in integrale sulla densità di carica $ \rho $, così da 
ottenere il termine di dipolo elettrico $ \v{p} $. Sfruttiamo l'identità vettoriale \eqref{identity:_divergenza(scalare_per_vettore)} applicata al caso
\begin{align}
  \nabla' \cdot \left( x' \, \v{J}_\omega(\v{r'}) \right) &= \vers{x}' \cdot \v{J}_\omega(\v{r'}) + x' \nabla'\!\cdot\!\v{ J }_\omega(\v{r'}) 
  \\
  &= J_{x'}^\omega - x' \der{\rho_\omega(\v{r'})}{t} = J_{x}^\omega + x' i\omega \rho_\omega(\v{r'})
\end{align}
dove nell'ultimo passaggio la $ J_{x'}^\omega $ è diventata $ J_{x}^\omega $, questo perché nell'integrale è valutata nel punto $ \v{r'} $ quindi la 
direzione è sempre la stessa. A questo possiamo subito notare che il termine $ \nabla' \cdot \left( x' \, \v{J}_\omega(\v{r'}) \right) $ può essere 
trasformato in un integrale di superficie per il teorema della divergenza ed è quindi nullo per l'ipotesi di sorgenti localizzate dato che la superficie 
va all'infinito. Rimane soltanto il termine con $ \rho(\v{r'}) $
\begin{align}
  A_{x}^\omega(\v{r}) &= -k_m \frac{e^{ikR}}{R} \int i\omega x'\rho_\omega(\v{r'}) \,d^3r' 
  = -k_m \, i \omega \frac{e^{ikR}}{R} p_{x}^\omega
  \\
  \implies \v{A}_{\omega}(\v{r}) &= -k_m \,i \omega \frac{e^{ikR}}{R} \v{p}_{\omega}
\end{align}
Per trovare i campi dobbiamo farne il rotore, ma il dipolo $ \v{p}_\omega $ non dipende dalle coordinate e a grandi distanze il termine dominante è la 
derivata rispetto a $ R $ della fase cioè $ \der{}{R} \left( e^{ikR} \right) $ perché l'altro termine va come $ 1/R^2 $. Quindi applicando il rotore 
otteniamo
\begin{align}
  \v{B}_\omega(\v{r}) &= \nabla\!\times\!\v{ A }_\omega(\v{r}) \simeq i \v{k} \times \v{A}_\omega(\v{r})
  \\
  &= - k_m \,i\omega \frac{e^{ikR}}{R} i \v{k} \times \v{p}_\omega
  = k_m\,\omega \, \frac{e^{ikR}}{R}\, \v{k} \times \v{p}_\omega
\end{align}
dove abbiamo riscritto $ \v{k} = k \vers{u} $ come la direzione di propagazione dell'onda, così da tenere conto direttamente nel prodotto vettore la 
dipendenza dall'angolo compreso ($ \sin\theta $). Otteniamo ora il campo elettrico
\begin{align}
  \v{E}_\omega(\v{r}) &= \frac{i\,c^2}{\omega} \nabla\!\times\!\v{ B }_\omega \simeq \frac{i\,c^2}{\omega} \, i\v{k} \times \v{B}_\omega 
  = -k_e \, k^2\, \frac{e^{ikR}}{R} \vers{u} \times (\vers{u} \times \v{p}_\omega)
  \\
  &= k_e \, k^2\, \frac{e^{ikR}}{R} (\vers{u} \times \v{p}_\omega) \times \vers{u}
\end{align}
I campi che abbiamo ottenuto sono trovati nel dominio delle frequenze, per tornare al dominio del tempo occorre fare l'anti-trasformata del campo nel 
dominio delle frequenze $ \v{B}_\omega(\v{r}) $, cioè
\begin{align}
  \v{B}(\v{r},t) &= \mathcal{F}^{-1} \left[ \v{B}_\omega(\v{r}) \right] 
  = k_m \frac1{cR} \mathcal{F}^{-1} \left[ {\omega^2} e^{i\frac{\omega}{c} R} \vers{u} \times \v{p}_\omega \right]
  \\
  &= k_m \, \frac{\ddot{\v{p}}_{(t - \frac{R}{c})} \times \vers{u} }{c\,R} 
\end{align}
dove abbiamo invertito il prodotto vettore per cambiare segno e abbiamo usato il teorema del ritardo causale (\textit{time shifting}) relativo alla 
trasformata di Fourier che ci ha permetto di trasformare l'esponenziale complesso nel dominio delle frequenze in una tralazione temporale nel dominio 
del tempo. Facciamo la stessa cosa per il campo elettrico
\begin{align}
  \v{E}(\v{r},t) &= \mathcal{F}^{-1} \left[ \v{E}_\omega(\v{r}) \right]
  = \frac{k_e}{c^2\,R} \mathcal{F}^{-1} \left[ {\omega^2} e^{i\frac{\omega}{c} R} (\vers{u} \times \v{p}_\omega) \times \vers{u} \right]
  \\
  &= k_e \,\frac{( \ddot{\v{p}}_{(t - \frac{R}{c})} \times \vers{u} ) \times \vers{u}}{c^2 \,R} 
\end{align}
Adesso che abbiamo ottenuto entrambi i campi possiamo trovare il vettore di Poynting $ \v{S} $ e la media temporale calcolata sul periodo del suo modulo che 
costituisce una grandezza di ampio interesse chiamata \textit{intensità}:
\begin{align} 
  I = \langle |\v{S}| \rangle_T.
\end{align}
Troviamo quindi $ \v{S} $, per svolgere questo prodotto vettore di cinque termini osserviamo che se chiamiamo 
$ \v{a} = \ddot{\v{p}}_{(t_r)} \times \vers{u} $ si può applicare la regola $"{bac-cab}"$:
\begin{align}
  \v{S} (\v{r}, t) &= \frac{1}{\mu_0} \v{E} \times \v{B} = \frac{k_e  k_m }{\mu_0 c^3\, R^2}\, (\v{a}\times \vers{u}) \times \v{a} 
  = \frac{1}{16 \pi^2 \eps_0 \, c^3\,R^2}\,a^2 \vers{u}
  \\
  &= \frac{\eps_0 \,k_e^2 }{c^3 R^2} \sin^2\theta\lvert \ddot{\v{p}}(t_r) \rvert^2 \vers{u}.
\end{align}
Notiamo che il vettore di Poynting è diretto verso il punto 
di osservazione e la sua distribuzione angolare varia con $ \sin^2\theta $ con $ \theta $ angolo polare, questo ci dice che il suo modulo è minimo per 
$ \theta = 0,\pi $ quindi sull'asse $ z $, mentre è massimo per $ \theta = \pi/2 $. Per trovare l'intensità ricordiamo che la sorgente è varia nel tempo 
come $ e^{-i\omega t} $ che è quindi un $ \cos\!\left( \omega(t - R/c) \right) $ calcolato al tempo ritardato, la derivata seconda elevata al quadrato 
ci porta un fattore $ \omega^4 $, rimane l'integrale sul periodo del $ \cos^2(\omega (t - R/c))$ che ci porta un fattore $ 1/2 $, dunque abbiamo che
\begin{align}
  I = \frac{\eps_0 k_e^2 }{2c^3 R^2} \omega^4 \sin^2 \theta\lvert \v{p}_0 \rvert^2.
\end{align}
Un dettaglio interessante è che il vettore di Poynting, e di conseguenza l'intensità, decrescono come $ 1/R^2 $, il loro integrale su una superficie è 
quindi costante ed è proprio ciò che ci aspettavamo di ottenere, che significa che osservare l'intensità emessa non dipende dalla distanza a cui lo facciamo, 
a patto di osservare su tutto l'angolo solido. Infatti la dipendenza dell'intensità dall'angolo solido è costante:
\begin{align}
  {dI} = R^2 \langle \v{S} \cdot \vers{n} \rangle {d\Omega}
  = \frac{\eps_0 k_e^2 }{2c^3} \omega^4 \sin^2 \theta\lvert \v{p}_0 \rvert^2 \,{d\Omega}
\end{align}
A questo punto è facile ottenere la potenza totale irraggiata che è il flusso del vettore di Poynting
\begin{align}
  W &= \Phi_\Sigma(\v{S}) = \int |\v{S}| R^2 \,d\Omega 
  = \frac{2\pi\eps_0 k_e^2}{c^3} \, \lvert \ddot{\v{p}}_0 \rvert^2 \int_0^\pi sin^2\theta \cdot \sin \theta \,d\theta 
  \\
  &= \frac{8\pi\eps_0 k_e^2}{3\,c^3}\, \lvert \ddot{\v{p}}_0 \rvert^2 = \frac{1}{6\pi \eps_0 \,c^3}\,\lvert \ddot{\v{p}}_0 \rvert^2 
  = \frac{2 k_e}{3c^3} \lvert \ddot{\v{p}}_0 \rvert^2 
\end{align}
che mediandola su un periodo si ha
\begin{align}
  \langle W \rangle =  \frac{2 k_e\, \langle \lvert \ddot{\v{p}}_0 \rvert^2 \rangle}{3\,c^3} 
  = \frac{k_0}{3\,c^3}  \omega^4 \lvert {\v{p}}_0 \rvert^2
\end{align}
e questa potenza è detta potenza di \textit{Larmor}. La formula di Larmor per la potenza riguarda più in generale la potenza emessa da una carica in 
moto accelerato, mentre noi l'abbiamo ricavata nel caso specifico di un dipolo che oscilla in moto armonico. Più avanti vedremo in generale le formule 
per ottenere la potenza irraggiata da una carica in un moto generico.

\subsection{Irraggiamento nell'atomo di idrogeno}
Applichiamo quello che abbiamo appena ricavato sulla teoria dell'irraggiamento all'atomo di idrogeno nel modello planetario (modello di Rutherford), 
cioè con il protone carico positivo al centro e l'elettrone che lo orbita, e supponiamo che ci sia una onda elettromagnetica incidente.

Le forze agenti sull'elettrone sono la forza di Lorenz $ \v{F}_L = q \left( \v{E} + \v{v} \times \v{B} \right) $ dovuta all'onda che incide su $ e^- $, 
la forza di Coulomb $ \v{F}_C $ dovuta al protone di carica opposta e una forza che possiamo esprimere come un attrito dinamico dovuto all'irraggiamento 
dell'elettrone stesso. Il motivo per cui imponiamo che l'irraggiamento dell'elettrone si manifesti come una forza dissipativa è che irraggiando perde 
energia, quindi è del tutto plausibile che sia simile a un attrito, inoltre se non inserissimo questo termine dovremmo farlo uscire dai campi 
elettromagnetici totali e quindi affrontare il problema della \textit{self force}. Questo problema è notoriamente complicato sia classicamente che andando 
ad avanzare teorie più sofisticate che tentano di spiegare cosa sia effettivamente l'elettrone, quindi imporre questa forma è più che conveniente. 
Scriviamo quindi le equazioni del moto
\begin{align}
  m_e \ddot{\v{r}} = -e\left( \v{E} + \v{v} \times \v{B} \right) - m_e \omega_0^2 \v{r} - m_e \gamma \dot{\v{r}}
\end{align}
dove abbiamo già espresso la forza di Coulomb in forma di piccole oscillazioni attorno alla posizione di equilibrio. Per velocità non relativistiche 
possiamo trascurare la forza magnetica dato che rispetto alla forza elettrica ha un fattore $ v/c $. Considerando un campo elettrico con una dipendenza del 
tipo $ \v{E} = E_0 e^{-i\omega t} \vers{z} $ abbiamo che il sistema è un oscillatore armonico forzato cioè
\begin{align}
  \ddot{z} + \gamma \dot{z} + \omega_0^2 z = - \frac{eE_0}{m_e} e^{-i\omega t}.
\end{align}
Questa equazione ha come soluzione omogenea quella dell'oscillatore smorzato con tempo caratteristico $ \uptau = 2/\gamma $, che rappresenta il tempo 
transitorio iniziale. Per $ t \gg \uptau $ rimane soltanto la soluzione particolare che supponiamo essere della forma della forzante:
\begin{align}
  \left[ \der[2]{}{t} + \gamma \der{}{t} + \omega_0^2 \right] \left( z_0 e^{-i\omega t} \right) 
  &= \left( -\omega^2 z_0 -i \omega \gamma z_0 + \omega_0^2 z_0 \right) e^{-i\omega t} \overset{!}{=} - \frac{eE_0}{m_e} e^{-i\omega t}
  \\
  \implies z_0 &= \frac{{eE_0}/{m_e}}{\omega^2 -\omega_0^2 + i\omega \gamma}
\end{align}
e quindi la forma completa è
\begin{align}
  z_{p}(t) &= \frac{{eE_0}/{m_e}}{\omega^2 -\omega_0^2 + i\omega \gamma} \,e^{-i\omega t}.
\end{align}
Ricordiamo che usiamo la notazione in cui $ e^- = \SI{-1.6e-19}{\C} $ mentre $ e = |e^-| $. 
Da questa forma si nota chiaramente il fenomeno della risonanza dato che per $ \omega = \omega_0 $ c'è un massimo nel modulo dell'ampiezza $ |z_0| $ e 
soprattutto abbiamo che $ z_p $ è in \textit{quadratura} con la forzante, infatti abbiamo che $ z_p \propto \sin(\omega t) $, risulta quindi sfasata di 
$ \pi/2 $ rispetto a $ E_z $ che quando andiamo a vedere la potenza $ W = F \cdot \v \propto \cos^2 (\omega t) $. Risulta quindi che la potenza assorbita 
dall'elettrone alla frequenza di risonanza è massima. Notiamo che il termine dissipativo dovuto all'irraggiamento svolge un ruolo fondamentale che 
legittimizza ancora di più la sua presenza: impedisce al denominatore di divergere, il che avrebbe portato a dire che l'elettrone assorbe energia infinita 
che è chiaramente assurdo. D'ora trascureremo l'attrito ponendo $ \gamma = 0 $, più avanti verificheremo quando è possibile farlo e cosa succede nel 
considerarlo. Per trovare la potenza irraggiata usiamo la relazione di Larmor trovata nella sezione precedente:
\begin{align}
  \langle W \rangle_{\textit{Larmor}} = \frac{k_e \,\omega ^4 |\v{p}_0|^2}{3\, c^3} &= \frac{k_e \,\omega ^4 |e\,\ddot{\v{z}}_p|^2}{3\, c^3} 
  = \frac{k_e \, e^4 E_0^2}{3\,c^3\,m_e^2} \frac{\omega^4}{(\omega^2 -\omega_0^2)^2}
  \\
  &= \frac{k_e \, e^4 E_0^2}{3\,c^3\,m_e^2} \, F(\omega).
\end{align}
Adesso definiamo una grandezza chiamata \textit{sezione d'urto}, che ha le dimensioni di una superficie, come il rapporto fra la potenza emessa dalla 
carica eccitata e l'intensità dell'onda incidente, cioè 
\begin{align}
  \sigma = \frac{\langle W \rangle_{\text{diff}}}{I_{\text{inc}}}
\end{align}
e siccome l'intensità di un onda piana vale 
\begin{align}
  I_{\text{piana}} = \langle |\v{S}| \rangle = \eps_0 c^2\langle \frac{ E_0^2}{c} \cos^2(kx - \omega t) \rangle = \frac{\eps_0 \, c \,E_0^2}{2}
\end{align}
quindi la sezione d'urto vale
\begin{align}
  \sigma = \frac{2 \, k_e \, e^4}{3\,\eps_0\,c^4 \,m_e^2}\, F(\omega) = \frac{8\pi}{3} \left( \frac{k_e\,e^2}{m_e^2\,c^2} \right)^2 F(\omega) 
  \eqcolon \frac{8\pi}{3}\, R_e^2 \,F(\omega)
\end{align}
dove abbiamo definito il raggio classico dell'elettrone $ R_e \coloneq \frac{k_e\,e^2}{m_e\,c^2} $ che si ottiene eguagliando l'energia a riposo 
con l'energia coulombiana. Analizziamo la sezione d'urto nei due limiti: 
\begin{align}
  \left.1\right)\qquad &\omega \gg \omega_0 \implies F(\omega) \simeq 1
  \\
  \left.2\right)\qquad &\omega \ll \omega_0 \implies F(\omega) \simeq \left( \frac{\omega}{\omega_0} \right)^4.
\end{align}
Il primo caso è chiamato \textit{scattering Thomson} o scattering di elettrone libero, questo perché il moto di rivoluzione è "invisibile" rispetto alla 
frequenza con cui l'onda elettromagnetica incide. Il secondo caso è detto \textit{scattering Rayleigh} è interessante per le interazioni luce-materia, 
dato che l'onda incidente non ha abbastanza energia per alterare la struttura atomica. Lo scattering Rayleigh è il motivo per il quale il cielo appare di 
colore blu: apparte il disco solare, quando guardiamo il cielo stiamo vedendo le regioni più prossime di un cono di diffusione che si crea quando la 
radiazione solare incontra l'atmosfera terrestre. Il colore che domina in questo modello è quello con la frequenza più alta, tuttavia è non può essere il 
viola per via dello spettro di emissione del Sole, dunque rimane il blu. Al tramonto la radiazione ha un percorso maggiore su cui diffondere e dunque 
rimangono i colori che hanno diffuso di meno e cioè il rosso.

\chapter{Esercizi elettrodinamica}
\section{Circuiti}
\subsection{Circuito con ramo mobile}
Sia un circuito piano rettangolare costituito da: una \textit{ddp} $ V_0 $, una resistenza ohmica $ R $, un ramo mobile di lunghezza $ h $.
Il circuito è sottoposto a un campo magnetico entrante $ \v{ B } = B_0 $ costante. Considerando il sistema inizialmente in quiete, trovare la 
velocità in funzione del tempo $ v(t) $ del ramo mobile e l'energia dissipata per $ t \to \infty $.

\textit{Soluzione} 
Sul ramo mobile agisce la forza di Lorentz,

\subsection{Circuito L-C infinito}
Si consideri la ripetizione orizzontale infinita di una maglia quadrata formata da: un condensatore $ C $, una induttanza $ L $, 
un condensatore $ C $ e l'ultimo ramo vuoto, e si indichi il lato con $ \Delta x $, che è messo a terra, come rappresentato in figura 
\ref{figure:_circuito_LC_infinito}. Si vuole trovare una espressione per la corrente.
\begin{figure}[h]
    \center
    \begin{circuitikz}
        \def \x {3}
        \draw
        ({-2*\x},\x) to[american inductor = $ L $] (-\x,\x) to[american inductor = $ L $] (0,\x) to[american inductor = $ L $] (\x,\x);

        \draw (-\x,\x) to[C = $ C $] (-\x,0);
        \draw (0,\x) to[C = $ C $] (0,0);

        \draw ({-2*\x},0) -- (\x,0);
        \node at ({-2*\x -1},{\x/2}) {\dots};
        \node at ({\x + 1},{\x/2}) {\dots};

        % AGGIUNGERE GENERATORE DI CORRENTE




    \end{circuitikz}
    \caption{Circuito LC infinito}
    \label{figure:_circuito_LC_infinito}
\end{figure}

\textit{Soluzione} Per risolvere questi problemi è utile trovare una relazione ricorsiva fra le varie quantità.
Applichiamo la legge dei nodi al nodo $ n $-esimo: la corrente $ I_{n-1} $ si separa nella corrente $ I_n $ e in quella che va a modificare la carica sull
armatura del condensatore $ n $-esimo cioè
\begin{align}
  \frac{dQ_n}{dt} = I_{n-1} - I_n
\end{align}
Una relazione per l'induttanza è data dalla differenza di potenziale ai sui capi:
\begin{align}
  \frac{1}{C} \left( Q_n - Q_{n+1} \right) = V_n - V_{n+1} = L \frac{dI_n}{dt}
\end{align}
e derivando rispetto al tempo otteniamo
\begin{align}
  L\frac{d^2I_n}{dt^2} = \frac{1}{C} \left( \frac{dQ_n}{dt} - \frac{dQ_{n+1}}{dt} \right) = \frac{1}{C} \left( I_{n-1} -2I_n + I_{n+1} \right)
\end{align}
Definiamo $ \omega_0^2 = \frac{1}{LC} $, in analogia con l'equazione delle onde, questo perché l'equazione che abbiamo ottenuto è ciò che si ottiene per 
esempio nel problema delle infinite masse collegate l'una con la successiva da una molla di uguale costante elastica $ k $. In quel problema si vede che è 
possibile propagare una onda che dipende da $ m $ e $ k $. Supponiamo che il generatore di corrente produca una $ I_G(t) = I_0 e^{-i\omega t} $, e questo 
possiamo farlo senza perdita di generalità dato che per la sovrapposizione di Fourier possiamo ottenere qualsiasi funzione come valore della corrente. 
La coordinata del nodo $ n $-esimo è $ x_n = na $, quindi vogliamo provare una soluzione del tipo
\begin{align}
  I_n(t) = I_0 e^{ikx_n - i\omega t}
\end{align}
e quindi proviamo a sostituire questa espressione nell'equazione ricorsiva che abbiamo ottenuto:
\begin{align}
  \frac{d^2I_n}{dt^2} &= - \omega^2 I_n = \omega_0^2 I_0 e^{-i\omega t} \left( e^{ik(n-1)a} -2e^{ikna}+ e^{ik(n+1)a}\right) 
  \\
  &= \omega_0^2 I_0 e^{ikna -i\omega t} \left( e^{-ika} -2 + e^{ika}\right) = \omega_0^2 I_n \left( -4 \sin^2\left( \frac{ka}{2} \right) \right)
\end{align}
dunque la corrente si semplifica e otteniamo la relazione di dispersione
\begin{align}
  \omega = 2\omega_0 \left|\sin \left( \frac{ka}{2} \right)\right|.
\end{align}
Studiando questa relazione di dispersione si nota che per valori piccoli di $ k $ diventa lineare e dunque la velocità di gruppo è ben definita, questo 
significa che è possibile trasmettere segnali in modo consistente. In questo caso la velocità di fase vale
\begin{align}
  v = \frac{\omega}{k} = \omega_0a
\end{align}

\section{Cavità cilindrica} \label{section:_esercizio_cavità_cilindrica}
Vogliamo trovare le soluzioni alle equazioni di Maxwell in una cavità cilindrica limitata. 

Supponiamo i campi abbiano una dipendenza del tipo
\begin{align}
  \v{ E }(\v{ r }, t) = \mathcal{E}(\rho)e^{-i\omega t} \vers{z}
  \\
  \v{ B }(\v{ r }, t) = \mathcal{B}(\rho)e^{-i\omega t} \vers{\varphi}
\end{align}
in questo modo abbiamo che non dipendono da $ z $, hanno simmetria cilindrica e possiamo imporre successivamente l'annullamento ai bordi. Per 
risolvere le equazioni ci concentriamo sul trovare prima il campo elettrico e per trovare il campo magnetico basterà calcolarne il rotore e dividere per 
il fattore $ i\omega $ dovuto alla derivata rispetto al tempo:
\begin{align}
  \nabla\!\times\!\v{ E } = i\omega \v{ B }.
\end{align}
Scriviamo il rotore di rotore di $ \v{ E } $
\begin{align}
  \nabla \times \left( \nabla \times \v{ E } \right) = - \nabla^2 \v{ E } = \frac{i\omega}{c^2} \der{\v{ E }}{t} 
  = \frac{\omega^2}{c^2} \v{ E }
\end{align}
quindi rimane la componente $ z $ del laplaciano in coordinate cilindriche
\begin{align}
  -\frac{1}{\rho} \der{}{\rho} \left( \rho \der{E_z}{\rho} \right) = \frac{\omega^2}{c^2} E_z
\end{align}
L'equazione di cui dobbiamo trovare gli zeri è 
\begin{align}
  \frac{1}{\rho} \der{}{\rho} \left( \rho \der{E_z}{\rho} \right) + \frac{\omega^2}{c^2} E_z = 0
  \\
  \implies \frac{1}{\rho}\der{E_z}{\rho} + \der[2]{E_z}{\rho} + \frac{\omega^2}{c^2} E_z = 0
  \\
  \implies \rho^2\der[2]{E_z}{\rho} + \rho \der{E_z}{\rho} + \frac{\rho^2\omega^2}{c^2} E_z = 0
\end{align}
che è una equazione riconducibile a quella di Bessel tramite la sostituzione $ x = \frac{\rho \omega}{c} $
\begin{align}
  x^2 \der[2]{E_z}{x} + x \der{E_z}{x} + x^2 = 0
\end{align}
dove abbiamo con il cambio di variabile abbiamo cambiato anche il coefficiente davanti le due derivate, in modo da far sparire 
i vari $ \frac{\omega}{c} $. Le funzioni di Bessel sono le $ J_{\alpha} $ che risolvono l'equazione
\begin{align}
  x^2 \frac{d^2J_\alpha}{dx^2} + x \frac{dJ_\alpha}{dx} + (x^2 - \alpha^2) J_\alpha = 0
\end{align}
che sono tabulate e conosciamo anche gli zeri, e questo ci è utile in modo da riscalare i punti in cui si annullano in base alle 
nostre condizioni al bordo, cioè che il campo sia nullo per $ \rho = a $ cioè $ x = \frac{a\omega}{c} $. Noi siamo interessati 
al caso in cui $ \alpha = 0 $ quindi $ J_0 $ che è
\begin{align}
  J_0(x)=\sum_{m=0}^{\infty} \frac{(-1)^m}{(m!)^2}\left(\frac{x}{2}\right)^{2m} = J_0(x)= 1 - \frac{x^2}{4} + \frac{x^4}{64} - \frac{x^6}{2304} + \dots
\end{align}
Quindi il campo lo possiamo scrivere come 
\begin{align}
  E_z = \mathcal{E}_0 J_0\left( \frac{\rho\omega}{c} \right) e^{-i\omega t} 
  = \mathcal{E}_0 e^{-i\omega t} \left( 1 - \frac{\rho^2 \omega^2}{4 c^2} \right) + \dots
\end{align}
e di conseguenza il campo magnetico vale
\begin{align}
  \v{ B } = \frac{1}{i\omega} \nabla\!\times\!\v{ E } = -\frac{1}{i\omega} \der{E_z}{\rho} \vers{\varphi}
  = - \frac{i\omega}{2 c^2}\, \mathcal{E}_0\, \rho \, e^{-i\omega t} + \dots
\end{align}

\section{Piano di corrente}
Sia la seguente distribuzione di corrente 
\begin{align}
  \v{ J }(\v{ r },t) = J_0 \, \delta(z) \, \Theta(t) \, \vers{x}
\end{align}
dove $ \Theta(t) $ è la funzione di Heaviside che fa $ 1 $ per $ t \geq 0 $, $ 0 $ altrimenti. La $ \v{ J } $ rappresenta una corrente superficiale 
sul piano $ z= 0 $ diretta lungo le $ x $ e che inizia a scorrere al tempo $ t = 0 $. Si trovino i campi elettrico e magnetico tramite i potenziali 
ritardati. \\

\textit{Soluzione}. Prima di risolvere vediamo che se non ci fosse la dipendenza dal tempo allora sarebbe il problema classico del piano di corrente, 
il quale genera un campo magnetico costante e uniforme orientato lungo $ \pm y $. Ci aspettiamo che i campi abbiano una forma simile ma che siano anche 
delle onde. Sappiamo che in presenza di sorgenti il campo magnetico si può ottenere come rotore del potenziale vettore ritardato che si esprime come
\begin{align}
  \v{ A }(\v{ r },t) 
  = \frac{\mu_0}{4\pi} \int \frac{\v{ J }\left( \v{ r' }, t - \frac{|\v{ r-r' }|}{c} \right)}{|\v{ r-r' }|} \,d^3r'
\end{align}
e grazie alla $ \delta(z) $ si trasforma in un integrale di superficie sul piano $ z= 0 $, con il vincolo che la sorgente parta a $ t = 0 $ quindi 
\begin{align}
  t - \frac{|\v{ r-r' }|}{c} > 0 \implies c\,t < |\v{ r-r' }| = \sqrt{\rho^2 + z^2} \eqcolon R
\end{align}
dove abbiamo scritto in coordinate cilindriche la distanza $ |\v{ r-r' }| $ che abbiamo chiamo $ R $. Il fatto fondamentale è che il vincolo ci 
limita il nostro integrale fino a un certo raggio, altrimenti ci sarebbe stato il problema della convergenza dell'integrale. Quindi il potenziale vale
\begin{align}
  \v{ A }(\v{ r },t) 
  &= \vers{x}\, \frac{\mu_0 J_0}{4\pi}\int_0^{2\pi} d\theta \int_0^{\sqrt{c^2 t^2 - z^2}}\frac{\rho }{\sqrt{\rho^2 + z^2}} \,d\rho
  \\
  &= \vers{x} \frac{\mu_0J_0}{2} \int_z^{ct} \frac{1}{R}\, R\,dR = \frac{\mu_0 J_0}{2} \left( ct - z \right) \vers{x}
\end{align}
dove abbiamo fatto la sostituzione che abbiamo definito prima con $ R $. Ricordiamo che il potenziale vettore ottenuto vale $ 0 $  per $ z > c\,t $. 
Ne facciamo il rotore e otteniamo il campo magnetico 
\begin{align}
  \v{ B } = 
  \begin{cases}
    - \frac{\mu_0J_0}{2} \, \vers{y} \quad \text{per} \quad 0 < z < c\,t  \\
    \frac{\mu_0J_0}{2} \, \vers{y} \quad \text{per} \quad -c\,t< z < 0 \\
    0 \qquad \text{altrimenti}
  \end{cases}
\end{align}
che rappresenta un onda quadra che si propaga nello spazio, coerente con quello che ci aspettavamo. Rimane da trovare il campo elettrico e sfruttiamo il 
fatto che nella gauge di Lorenz abbiamo 
\begin{align}
  \v{ E } = - \nabla V - \der{\v{ A }}{t}
\end{align}
ma dato che non ci sono sorgenti del campo elettrico, il potenziale scalare è nullo (o costante). Questo può essere giustificato anche dalla condizione 
che lega i potenziali nella gauge di Lorenz ovvero
\begin{align}
  \nabla \cdot \v{ A } + \frac{1}{c^2} \der{V}{t} = 0
\end{align}
ed è immediato vedere che il potenziale vettore ottenuto ha divergenza nulla. Di conseguenza il campo elettrico vale
\begin{align}
  \v{ E } = -\der{\v{ A }}{t} = -\frac{\mu_0J_0}{2}\, c \,\vers{x} \quad \text{per} \quad 0 < z < ct
\end{align}
Notiamo che le relazioni di ortogonalità dei campi e che $ E = c B $ sono rispettate. Per le $ z < 0  $ il campo elettrico è lo stesso dato che è un 
vettore polare (cioè un vettore \enquote{vero}).

\section{Dipoli oscillanti in fase}

\chapter{Elettromagnetismo in forma covariante}
Fin ora abbiamo sempre assunto che i portatori di carica avessero delle velocità non relativistiche, adesso considereremo il caso più generale possibile 
in cui la velocità può essere un qualunque valore compreso fra $ 0 $ e $ c $.
% ##########
\section{Potenziali di Liénard-Wiechert}
Vogliamo inizialmente trovare i campi generati da una carica in moto, partiamo scrivendo i potenziali ritardati:
\begin{align}
  V(\v{ x },t) = k_e \int \frac{\rho(\v{ x }', t - \frac{|\v{ x-x' }|}{c})}{|\v{ x-x' }|}\,d^3x'
\end{align}
e indichiamo con $ \v{ r } = \v{ x - x' } $. Siccome abbiamo una carica puntiforme, la cui posizione indichiamo con $ \v{r}_q(t) $, riscriviamo le 
sorgenti come
\begin{align}
  \rho(\v{r}, t) &= q \, \delta(\v{ r' } - \v{ r }_q(t))
  \\
  \v{J}(\v{r}, t) &= q \v{v}_q(t) \, \delta(\v{ r' } - \v{ r }_q(t))
\end{align}
e usiamo una proprietà della delta di Dirac ovvero che
\begin{align}
  \int \delta(f(x)) \,dx = \int \delta(y) \, \frac{1}{|\frac{df}{dx}(y)|} \,dy = \frac{1}{|\frac{df}{dx}(y = 0)|}
\end{align}
In questo modo possiamo svolgere l'integrale del potenziale ritardato
\begin{align}
  V(\v{ x },t) = k_e q \int \delta(\v{ x' } - \v{ X }(t - t - \frac{|\v{ x-x' }|}{c})) \frac{1}{|\v{ x-x' }|} \,d^3x'
  = \frac{q}{4\pi\eps} \int \delta(y) \, \frac{1}{|1 - \frac{\v{ v }\cdot \vers{r} }{c}|} \,dy = \frac{q}{4\pi\eps} 
\end{align}
\dots
\begin{align}
  V(\v{r},t) &= k_e \, q \left[ \frac{1}{R - \vec{\beta} \cdot \v{R}} \right]_{t_r} \mathbf{A}
  \\
  \v{A}(\v{r},t) &= k_m \, q \left[ \frac{\v{v}}{R - \vec{\beta} \cdot \v{R}} \right]_{t_r}
\end{align}
e questi sono detti \textit{potenziali di Liénard-Wiechert}, dove $ \v{R} = \v{r} - \v{r}_q(t_r)$ e $ R = |\v{R}| = c (t - t_r) $. 
Calcoliamo i potenziali nel caso del moto rettilineo uniforme lungo $ x $:

\section{Quadricorrente}
Per arrivare a costruire le equazioni di Maxwell nella cosiddetta \textit{formulazione covariante} è necessario applicare il formalismo dei quadrivettori 
agli oggetti dell'elettrodinamica come le sorgenti e i potenziali.

Partiamo dal fatto che la carica elettrica deve essere necessariamente un invariante di Lorentz, altrimenti si dovrebbe osservare che oggetti carichi 
modificano le loro proprietà elettriche in base alla velocità a cui si muovono, cosa che non osserviamo. Consideriamo quindi la densità di carica 
$ \rho(\v{r}) $, per sapere come trasforma notiamo che essendo una carica diviso un volume sappiamo che i volumi si contraggono nella direzione del boost 
quindi nel sistema di riferimento $ O' $ che si muove a velocità costante $ V $ rispetto a $ O $ si ha che la densità di carica vale
\begin{align}
  \rho'(\v{r}) = \gamma \, \rho(\v{r})
\end{align}
Quella che abbiamo espresso non è la legge di trasformazione completa, dato che può esistere una densità di corrente che in un certo sistema di riferimento 
diventa una densità di carica. Per trovare la densità di corrente dovuta a una $ \rho $ che viene messa in moto con una velocità $ \v{v} $ applichiamo il 
modello microscopico dei portatori di carica ottenendo
\begin{align}
  \v{J'}(\v{r}, t) = \rho'(\v{r}, t) \, \v{v}(\v{r},t) = \gamma \, \rho(\v{r}, t) \, \v{v}(\v{r},t)
\end{align}
quindi è facile notare che $ c\cdot \rho $ e $ \v{J} $ sono proporzionali alle componenti della quadrivelocità $ u^\mu $. Questo ci suggerisce che moltiplicando 
per $ c $ la densità di carica possiamo costruire un quadrivettore con le dimensioni di una corrente che chiameremo \textit{quadricorrente} $ J^\mu $:
\begin{align}
  J^\mu \coloneq \rho(x) \, u^\mu (x) = \left( c \rho , \v{J}\right)
\end{align}
dove $ x $ è il quadrivettore posizione. Sfruttiamo subito la quadricorrente per verificare l'equazione di continuità con questo nuovo formalismo: 
calcolando la quadridivergenza di $ J^\mu $ si ha
\begin{align}
  \partial_\mu J^\mu  = \left( \frac{1}{c} \der{}{t}, \nabla \right) \cdot J = \der{\rho}{t} + \nabla \cdot \v{J} = 0
\end{align}
che è proprio l'equazione di continuità.

\section{Formulazione covariante delle equazioni di Maxwell}
Vogliamo prima di tutto dimostrare che anche i potenziali $ V $ e $ \v{A} $ formano un quadrivettore, e questo si può fare prendendo l'equazione delle 
onde che i potenziali rispettano in gauge di Lorenz:
\begin{align}
  \begin{cases}
    \DS -\nabla^2V + \frac{1}{c^2} \der[2]{V}{t} = \frac{\rho}{\eps_0} \\
    \DS -\nabla^2\v{ A } + \frac{1}{c^2} \der[2]{\v{ A }}{t} = \mu_0\v{ J }
  \end{cases}
  \implies \Box A^\mu = \frac{J^\mu}{\eps_0 c}
\end{align}
dove $ A^\mu $ rappresenta il quadripotenziale definito come \footnoteH{Può essere definito anche come $ A^\mu = (V/c, \v{A}) $ con una conseguente 
ridefinizione delle successive quantità ed equazioni.}
\begin{align}
  A^\mu \coloneq \left( V, c \v{A} \right).
\end{align}
Abbiamo tutti gli strumenti per scrivere le equazioni di Maxwell in forma tensoriale. Per farlo dobbiamo trovare un tensore a due indici che sia 
gauge-invariante e le sue componenti siano i campi. Il candidato più semplice è 
\begin{align}
  T^{\mu \nu} = \partial^\mu A^\nu
\end{align}
ma questo oggetto ha 16 componenti e non è gauge-invariante perché
\begin{align}
  T^{\mu \nu} = \partial^\mu \left( A'^\nu + \partial^\nu(c\,\psi) \right) 
  = \partial^\mu A'^\nu + \partial^\mu\partial^\nu(c\,\psi) \neq \partial^\mu A'^\nu = T'^{\mu \nu}.
\end{align}
Proviamo allora con il cosiddetto \textit{tensore di Faraday} $ F^{\mu \nu} $ definito come
\begin{align}
  F^{\mu \nu} = \partial^\mu A^\nu - \partial^\nu A^\mu.
\end{align}
Notiamo che è antisimmetrico, il che ci dice che ha $ 6 $ componenti indipendenti ed è facile verificare che è gauge-invariante. La forma esplicita del 
tensore di Faraday è 
\begin{align}
  F^{\mu \nu} =
  \begin{pmatrix}
    0 & -E_x & -E_y & -E_z \\
    E_x & 0 & -c B_z & c B_y \\
    E_y & c B_z & 0 & -c B_x \\
    E_z & -c B_y & c B_x & 0 
  \end{pmatrix}
\end{align}
che abbiamo calcolato mediante le solite relazioni
\begin{align}
  \v{E} &= -\nabla V - \der{\v{A}}{t} 
  \\
  \v{B} &= \nabla\!\times\!\v{ A }
\end{align}
Dal tensore di Faraday derivano le equazioni non omogenee di Maxwell che si possono scrivere nella forma compatta
\begin{align}
  \partial_\mu F^{\mu \nu} = \mu_0 J^\nu
\end{align}
che si possono verificare contraendo l'indice cioè applicando la quadridivergenza. Per quanto riguarda le equazioni omogenee cerchiamo un 
tensore che in analogia al primo sia antisimmetrico e che ci esprima
\begin{align}
  &\nabla \cdot \v{B} = \partial_x B_x + \partial_y B_y + \partial_z B_z = 0
  \\
  &\left( \nabla \times \v{E} \right)_i = \eps_{i j k} \partial_j E_k = - \partial_t B_i
\end{align}
Per trovarlo notiamo che le equazioni sono praticamente le stesse che esprime il tensore di Faraday ma con i campi scambiati, quindi come la prima colonna 
$ F^{\mu 0} $ ci dà la divergenza del campo elettrico, imponiamo che il tensore che cerchiamo $ \tilde{F}^{\mu \nu} $ sia del tipo
\begin{align}
  \tilde{F}^{\mu \nu} = 
  \begin{pmatrix}
    0 & -c B_x & -cB_y & -cB_z \\
    c B_x & 0 & \ast & \ast \\
    c B_y & \ast & 0 & \ast \\
    c B_z & \ast & \ast & 0 
  \end{pmatrix}
\end{align}
dove abbiamo imposto l'antisimmetria. In questo modo applicando $ \partial_\mu \tilde{F}^{\mu 0} $ otteniamo la divergenza del campo magnetico che 
imponendola a $ 0 $ ci porta all'equazione
\begin{align}
  \partial_\mu \tilde{F}^{\mu \nu} = 0
\end{align}
dalla quale è effettivamente possibile ottenere anche quella per il rotore di $ \v{E} $ semplicemente aggiungendo al posto degli asterischi le 
componenti giuste del campo elettrico, che risulta essere
\begin{align}
  \tilde{F}^{\mu \nu} = 
  \begin{pmatrix}
    0 & -c B_x & -cB_y & -cB_z \\
    c B_x & 0 & -E_z & E_y \\
    c B_y & E_z & 0 & -E_x \\
    c B_z & -E_y & E_x & 0 
  \end{pmatrix}
\end{align}
noto come \textit{tensore di campo duale}. Notiamo che è possibile ottenere il tensore duale da quello di Faraday tramite il tensore completamente 
antisimmetrico $ \eps_{\mu \nu \rho \sigma} $ che è il tensore di Levi-Civita a 4 indici, definito come 
\begin{align}
  \eps_{\mu \nu \rho \sigma} = 
  \begin{cases}
    +1 &\;\; \textnormal{se $ (\mu \nu \rho \sigma) $ è una permutazione pari di $ (0123) $} \\
    -1 &\;\; \textnormal{se $ (\mu \nu \rho \sigma) $ è una permutazione dispari di $ (0123) $} \\
    0 &\;\; \textnormal{se almeno due indici sono uguali}
  \end{cases}
\end{align}
e per ottenere la versione con gli indici alzati si ha
\begin{align}
  \eps^{\mu \nu \rho \sigma} = g^{\mu \alpha} \, g^{\nu \beta} \, g^{\rho \gamma} \, g^{\sigma \delta} \, \eps_{\alpha \beta \gamma \delta} 
  = \det(g^{\mu \nu})\, \eps_{\mu \nu \rho \sigma} = - \eps^{\mu \nu \rho \sigma}.
\end{align}
In termini di $ \eps^{\mu \nu \rho \sigma} $ il tensore duale si scrive come
\begin{align}
  \tilde{F}^{\mu \nu} = \frac{1}{2} \eps^{\mu \nu \rho \sigma} F_{\rho \sigma}
\end{align}
dove è necessario abbassare gli indici di $ F^{\mu \nu} $ con la metrica
\begin{align}
  F_{\mu \nu} = g_{\mu \alpha} g_{\nu \beta} F^{\alpha \beta}
\end{align}
e il fattore $ 1/2 $ è dovuto all'antisimmetria. In questo modo le equazioni di Maxwell nel formalismo covariante sono espresse dalle relazioni compatte:
\begin{align}
  &\partial_\mu F^{\mu \nu} = \frac{J^\nu}{\eps_0 c} 
  \\
  &\partial_\mu \tilde{F}^{\mu \nu} = 0
\end{align}
In questa forma si nota chiaramente la simmetria dei campi nelle equazioni di Maxwell, che sarebbe completa se esistesse la carica magnetica introducendo 
così una quadricorrente magnetica che farebbe da termine noto alla seconda equazione.

\section{Invarianti elettromagnetici}
Vogliamo trovare delle quantità costruite a partire dai tensori elettromagnetici che siano invarianti per trasformazioni di Lorentz. Un modo per farlo è 
prendere lo scalare associato a un tensore: possiamo fare i calcoli espliciti oppure prendendo la traccia del prodotto fra le due matrici
\begin{gather}
  F^{\mu \nu} F_{\mu \nu} = \text{tr} \left( F^{\mu \nu} \, F_{\mu \nu}{}^t \right)
  \\
  = \text{tr} \left[  
  \begin{pmatrix}
    0 & -E_x & -E_y & -E_z \\
    E_x & 0 & -c B_z & c B_y \\
    E_y & c B_z & 0 & -c B_x \\
    E_z & -c B_y & c B_x & 0 
  \end{pmatrix}
  \begin{pmatrix}
    0 & E_x & E_y & E_z \\
    -E_x & 0 & -c B_z & c B_y \\
    -E_y & c B_z & 0 & -c B_x \\
    -E_z & -c B_y & c B_x & 0 
  \end{pmatrix}^t \,
  \right]
  \\
  = 2 \left( c^2\v{B}^2 - \v{E}^2 \right) 
\end{gather}
Questo primo invariante elettromagnetico contiene già molte informazioni: il segno ci dice che se in un sistema di riferimento abbiamo 
$ \v{E}^2 > c^2\v{B}^2 $ ($ < $) allora questo è vero in ogni sistema di riferimento. 
Il secondo invariante lo troviamo tramite il prodotto
\begin{gather}
  F^{\mu \nu} \tilde{F}_{\mu \nu} = \text{tr} \left[  
  \begin{pmatrix}
    0 & -E_x & -E_y & -E_z \\
    E_x & 0 & -c B_z & c B_y \\
    E_y & c B_z & 0 & -c B_x \\
    E_z & -c B_y & c B_x & 0 
  \end{pmatrix}
  \begin{pmatrix}
    0 & c B_x & cB_y & cB_z \\
    -c B_x & 0 & -E_z & E_y \\
    -c B_y & E_z & 0 & -E_x \\
    -c B_z & -E_y & E_x & 0 
  \end{pmatrix}^t \,
  \right]
  \\
  = - 4 c^2 \, \v{E} \cdot \v{B}
\end{gather}
il quale ci dice che se i campi sono ortogonali in un sistema di riferimento allora lo sono in tutti.

\section{Trasformazioni di Lorentz dei campi}
Dato che abbiamo scritto in maniera compatta il campo elettromagnetico con i tensori elettromagnetici sfruttiamo le leggi di trasformazione dei tensori per 
ottenere come cambiano i campi cambiando sistema di riferimento inerziale. Un tensore a due indici trasforma come
\begin{align}
  F'^{\mu \nu} = \Lambda^\mu{}_\rho \, \Lambda^\nu{}_\sigma \, F^{\rho \sigma} = \left( \Lambda \, F \, \Lambda^T \right)_{\mu \nu}
\end{align}
quindi per trovare i campi trasformati ci basta svolvere questo prodotto. Indicando con $ \parallel $ e $ \perp $ le componenti parallele e 
perpendicolari alla direzione del boost otteniamo 
\begin{align}
  &{E'_\parallel = E_\parallel} \qquad {\v{E'}_\perp = \gamma \left( \v{E}_\perp + \v{V} \times \v{B} \right)}
  \\
  &{B'_\parallel = B_\parallel} \qquad {\v{B'}_\perp = \gamma \left( \v{B}_\perp - \frac{\v{V}}{c^2} \times \v{E} \right)}
\end{align}
dove $ \v{V} $ è la velocità del boost cioè la velocità di $ O' $ rispetto a $ O $.

\section{Effetto Doppler-Fizeau}
Abbiamo visto come i campi trasformano fra i sistemi di riferimento, che è però una relazione locale cioè non ci fornisce ancora informazioni su come 
cambia una soluzione delle equazioni di Maxwell. Sempre per Fourier ci conviene lavorare con le onde piane, vogliamo quindi capire come cambiano i 
suoi parametri $ \omega, \v{k} $ in un sistema di riferimento in moto rispetto alla sorgente dell'onda.

Le relazioni che legano i campi ai relativi trasformati abbiamo detto che sono locali, quindi il campo trasformato dipende dal campo nel sistema di 
riferimento della sorgente (che chiameremo $ O $) nel corrispondente evento, ed è anche una relazione lineare. Queste caratteristiche ci dicono che 
se una soluzione è di tipo oscillante in un sistema, deve necessariamente esserlo anche nell'altro, e dato che assume periodicamente gli stessi 
valori in un certo sottoinsieme di eventi dello spazio-tempo, allora quegli eventi devono avere lo stesso valore anche in $ O' $. Dunque la soluzione 
trasformata deve avere la stessa fase della soluzione originale a meno di una costante arbitraria che scegliamo $ 0 $. Quindi le due fasi rispettano
\begin{align}
  \omega' t' - \v{k'} \cdot \v{x'} = \omega t - \v{k} \cdot \v{x}
\end{align}
riscriviamo adesso questa relazione con il quadrivettore posizione $ x = (c t, \v{x}) $ e un certo $ k \coloneq (\omega/c, \v{k}) $ e otteniamo
\begin{gather}
  k^t g x = k'^t g x'= k'^t g \Lambda x \;\; \forall x \;\; \implies k'^tg \Lambda = k^t g 
  \\
  k'^t = k^t g \Lambda^{-1} g^{-1} = k^t \left( g \Lambda g\right)^{-1} \implies k' = \left( g \Lambda g\right)^{-1}{}^t\, k
\end{gather}
dove abbiamo usato che $ g^{-1} = g $. Dato che il gruppo di Lorentz può essere definito come
\footnoteH{Anche il gruppo simplettico ammette una definizione 
simile con la matrice 
$ \Gamma = \DS 
\begin{pmatrix}
  0 & \text{Id} \\ 
  - \text{Id} & 0
\end{pmatrix} $ al posto della metrica $ g $.} 
\begin{align}
  \Lambda^t g \Lambda = g ,
\end{align}
invertendo questa esplicitando la $ \Lambda^t $ si ottiene proprio la quantità che abbiamo ottenuto nella trasformazione di $ k $, 
quindi abbiamo 
\begin{align}
  k' = \Lambda k
\end{align}
e dunque $ k $ deve trasformare come un quadrivettore. In questo modo possiamo esprimere la fase come contrazione degli indici: 
$ \phi = k_\mu x^\mu $. Grazie a ciò possiamo trovare la frequenza e il vettore d'onda trasformati:
\begin{align}
  \begin{cases}
    \DS \omega' = \gamma\left( \omega - \vec{\beta} \cdot c\, \v{k} \right)
    \\
    \DS \v{k}' = \gamma \left( \v{k} - \vec{\beta} \, \frac{\omega}{c} \right)
  \end{cases}
\end{align}
e se vale la relazione di dispersione del vuoto cioè $ \v{k}^2 = \frac{\omega^2}{c^2} $ allora la frequenza trasformata diventa
\begin{align}
  \omega' = \omega \gamma \left( 1 - \vec{\beta} \cdot \vers{n} \right)
\end{align}
dove $ \vers{n} $ è la direzione di propagazione dell'onda. Nel caso in cui $ \vec{\beta} \parallel  \vers{n} $, 
che è detto \textit{effetto Doppler longitudinale}, l'espressione per $ \omega' $ si semplifica ottenendo
\begin{align}
  \omega' = \omega \gamma \left( 1 - \beta \right) = \omega \sqrt{\frac{1 - \beta}{1 + \beta}}
\end{align}
che l'analogo relativistico del effetto Doppler classico. Se esprimiamo con $ \theta $ l'angolo compreso fra la direzione del boost e la direzione di 
propagazione dell'onda, si ha che $ \vec{\beta} \cdot \vers{n} = \beta \cos\theta $ ottenendo
\begin{align}
  \omega' = \omega \gamma (1 - \beta \cos\theta ) \; \overset{\theta = \frac{\pi}{2}}{\longrightarrow}\; \omega' = \gamma \omega
\end{align}
dove abbiamo esplicitato il caso particolare in cui $ \theta = \frac\pi2 $ che è detto \textit{effetto Doppler trasversale}. Quest'ultimo è particolarmente 
interessante perché non esiste un analogo effetto classico e ciò è dovuto al fattore di Lorentz $ \gamma $ che è presente nella versione relativistica, 
infatti tutto il procedimento che abbiamo fatto comprende anche il caso classico per via della struttura delle trasformazioni di Lorentz che in 
approssimazione non relativistica, $ v \ll c $, diventano le trasformazioni di Galileo.


\chapter{Elettromagnetismo nella materia} \label{chap: elettromagnetismo nella materia}
In questo capitolo studieremo come i campi elettromagnetici interagiscono con la materia e come è possibile modellizzarli. Le equazioni di Maxwell 
che abbiamo ricavato, insieme alla forza di Lorenz, costituiscono una teoria esatta per descrivere i campi elettromagnetici. L'aggettivo esatto è da 
intendersi nel senso che non è una teoria ricavata da una qualche approssimazione di una teoria più completa, come ad esempio la meccanica Newtoniana che 
è il limite di bassa velocità della relatività ristretta. Trattare fenomeni elettromagnetici con questa teoria è relativamente semplice quando il numero di 
cariche interagenti è piccolo, ma nella materia abbiamo un numero di atomi dell'ordine di $ N \sim 10^{24} $ per centimetro cubo. Quindi oltre che 
straordinariamente difficile, è addirittura inutile conoscere in modo esatto tutti gli effetti elettromagnetici in un oggetto macroscopico, dato che è 
possibile ottenere le informazioni a cui siamo interessati da delle grandezze studiabili dal punto di vista macroscopico, come le funzioni di stato.

\section{Equazioni di Maxwell nella materia}
Quello che faremo è costruire una nuova teoria elettromagnetica valida in un certo regime di lunghezze che chiameremo \textit{scala mesoscopica}: ovvero una 
scala che ci permette di non considerare la natura discreta degli reticoli atomici (e soprattutto gli effetti quantistici), ma comunque abbastanza piccola 
rispetto alle scale di variazione spaziale dei campi in modo da essere comunque considerata infinitesima. Per dare dei riferimenti, la scala mesoscopica 
che applicheremo è abbastanza lontana dal raggio di Bohr $ a \sim \SI{e-10}{m} $ e dalla lunghezza d'onda $ \lambda \sim \SI{100}{nm} = \SI{e-7}{m}$ 
dei campi:
\begin{align}
  a \ll R \ll \lambda \tag*{[scala mesoscopica]}
\end{align}
Sfrutteremo la scala mesoscopica per introdurre un approccio comune quando si vuole passare a una trattazione su scala maggiore di una qualche teoria 
microscopica, ovvero quello della teoria di campo medio: sostituiremo i campi microscopici $ \v{E}, \v{B} $ (rispettivamente campo elettrico e magnetico) 
con le loro medie in scala mesoscopica $ \langle \v{E} \rangle, \langle \v{B} \rangle $. La media che useremo è 
\begin{align}
  \langle \v{F}(\v{r},t) \rangle = \frac{1}{\Omega} \int_\Omega \v{F}(\v{r + r'}, t) \, \mu(\v{r'})\,d^3r'
\end{align}
dove $ \Omega $ è l'intorno mesoscopico e $ \mu(\v{r'}) $ è un peso che possiamo aggiungere. La cosa importante è che la media commuta con le derivate, 
possiamo quindi applicarla subito alle equazioni di Maxwell:
\begin{align}
  &\nabla \cdot \langle \v{B} \rangle = 0
  \\
  &\nabla \times \langle \v{E} \rangle = - \der{\langle \v{B} \rangle}{t}
  \\
  &\nabla \cdot \langle \v{E} \rangle = 4\pi k_e {\langle \rho \rangle}
  \\
  &\nabla \times \langle \v{B} \rangle = 4\pi k_m \langle \v{J} \rangle + \frac{k_m}{k_e} \der{\langle \v{E} \rangle}{t}
\end{align}
Per quanto riguarda le sorgenti dobbiamo distinguere fra i contributi esterni, i quali sono inseriti da noi dunque ne conosciamo necessariamente una media 
mesoscopica, e quelli microscopici, quindi
\begin{align}
  \langle \rho \rangle &= \rho_{\text{ext}} + \langle \rho_{\text{risposta}} \rangle
  \\
  \langle \v{J} \rangle &= \v{J}_{\text{ext}} + \langle \v{J}_{\text{risposta}} \rangle
\end{align}
e vogliamo quindi modellizzare questa risposta tramite l'approssimazione di dipolo. Dato che il termine di dipolo del potenziale si scrive come 
$ V_d = k_e \frac{\v{d} \cdot \v{r}}{r^3} $, possiamo scrivere il potenziale dividendo i termini esterni e di risposta come 
\begin{align}
  V(\v{r}) \simeq k_e \int d^3r'\left[ \frac{\rho_{\text{ext}}(\v{r'} )}{|\v{r - r'}|} + \frac{\v{P}(\v{r'}) \cdot (\v{r - r'})}{|\v{r-r'}|^3} \right]
\end{align}
dove abbiamo introdotto $ \v{P} $ che è detto  \textit{vettore di polarizzazione} e nella nostra trattazione corrisponde a una media dei dipoli in un certo 
punto del materiale:
\begin{align}
  \v{P}(\v{r}) = \langle \frac{\v{d}_{\text{indotti}}( \v{r})}{\uptau} \rangle
\end{align}
A questo punto applichiamo la solita identità vettoriale
\begin{align}
  \int d^3r'\; \v{P}(\v{r'}) \cdot \nabla' \left( \frac{1}{|\v{r-r'}|} \right) 
  = \int d^3r'\left[ \nabla' \cdot \left( \frac{\v{P}(\v{r'})}{|\v{r-r'}|} \right) - \frac{\nabla' \cdot \v{P}(\v{r'})}{|\v{r - r'}|} \right]
\end{align}
e il primo termine è nullo applicando il teorema della divergenza e per l'ipotesi di sorgenti localizzate. Dunque il potenziale diventa
\begin{align}
  V(\v{r}) \simeq k_e \int d^3r' \;\frac{1}{\v{|r - r'|}} \left[ \rho_{\text{ext}} - \nabla' \cdot \v{P}(\v{r'}) \right] 
  = k_e \int d^3r' \;\frac{1}{\v{|r - r'|}} \left[ \rho_{\text{ext}} - \rho_{\text{pol}}(\v{r'}) \right] 
\end{align}
dove abbiamo introdotto la \textit{densità di carica di polarizzazione} $ \rho_{\text{pol}} $ definita come 
\begin{align}
  \rho_{\text{pol}}(\v{r}) \coloneq - \nabla \cdot \v{P}(\v{r})
\end{align}
Tramite $ \rho_{\text{pol}} $ l'equazione per la divergenza del campo elettrico si riscrive
\begin{align}
  &\nabla \cdot \langle \v{E} \rangle = 4\pi k_e ({\rho_{\text{ext}}} +{\rho_{\text{pol}}})
  = 4\pi k_e ({\rho_{\text{ext}}} - \nabla \cdot {\v{P}})
  \\
  &\implies \nabla \cdot \left[ \langle \v{E} \rangle + 4\pi k_e {\v{P}} \right] = 4\pi k_e {\rho_{\text{ext}}}
\end{align}
e definiamo il \textit{vettore di induzione elettrica} $ \v{D} $ come
\begin{align}
  \v{D} \coloneq \frac{\langle \v{E} \rangle}{4\pi k_e}  + \v{P} \implies \nabla \cdot \v{D} = \rho_{\text{ext}}. \tag*{[SI]}
\end{align}
Qui in realtà il sistema internazionale (SI) e quello gaussiano (cgs) differiscono nella definizione, cosa che succederà anche fra poco per il vettore 
$ \v{H} $. Il motivo è che si vuole arrivare a scrivere le equazioni di Maxwell in due forme diverse: nel sistema internazionale si vuole scrivere 
dell'equazioni differenziali in funzione delle sorgenti (esterne) senza costanti moltiplicative, mentre nel sistema gaussiano rimangono le medesime 
costanti del caso microscopico, e ciò si traduce nel dividere la sorgente $ \rho_{\text{ext}} $ per $ 4\pi k_e $ oppure no. Da questo si capisce che per 
passare da $ \v{D} $ $ (\v{H}) $ da SI a cgs basta moltiplicare per $ 4\pi k_e $ $ (4\pi k_m) $.

Sfruttiamo lo stesso ragionamento con il potenziale vettore e il dipolo magnetico per arrivare a ottenere una nuova equazione anche per le correnti: dato
che $ \v{A}_\v{m} = k_m \frac{\v{m}\times \v{r}}{r^3} $ scriviamo il potenziale vettore totale come
\begin{align}
  \v{A}(\v{r}) \simeq k_m \int d^3r' \left[ \frac{\v{J}_{\text{ext}}(\v{r'})}{\v{|r -r'|}} 
  + \frac{ \v{M} \times \left( \v{r - r'} \right)}{\v{|r - r'|}^3} \right]
\end{align}
dove analogamente a $ \v{P} $ abbiamo introdotto il \textit{vettore di magnetizzazione} $ \v{M} $ definito come la media dei dipoli magnetici in un certo 
volume
\begin{align}
  \v{M}(\v{r}) \coloneq \langle \frac{\v{m}_{\text{indotti}}(\v{r})}{\uptau} \rangle
\end{align}
Sviluppiamo il termine contenente $ \v{M} $ con l'identità vettoriale del rotore e otteniamo
\begin{align}
  \int d^3r'\; \v{M}(\v{r'}) \times \nabla' \left( \frac{1}{\v{|r - r'|}} \right)
  = \int d^3 r' \left[ \frac{\nabla'\!\times\!\v{ M }(\v{r'})}{\v{|r - r'|}} - \nabla'\!\times\!\left( \v{ \frac{M(r')}{|r - r'|} }\right)\right]
\end{align}
adesso il secondo termine di questo integrale è una divergenza dunque si annulla per la localizzazione delle sorgenti, e per mostrarlo basti vedere che
\begin{align}
  \eps_{i j k} \partial_j M_k = \partial_j \eps_{i j k} M_k = \partial_j \eps_{j k i} M_i = \partial_j a_j
\end{align}
che è appunto una divergenza, dunque possiamo trasformarlo in un integrale di superficie. Il potenziale vettore assume la forma
\begin{align}
  \v{A}(\v{r}) \simeq k_m \int d^3r' \frac{1}{\v{|r - r'|}} \left[ \v{J}_{\text{ext}}(\v{r'}) + \nabla'\!\times\!\v{ \v{M}(r') } \right] 
  = k_m \int d^3r' \frac{1}{\v{|r - r'|}} \left[ \v{J}_{\text{ext}}(\v{r'}) + \v{J}_{\text{pol}}(\v{r'}) \right] 
\end{align}
dove $ \v{J}_{\text{pol}} $ è la \textit{corrente di polarizzazione}. L'equazione magnetostatica per il rotore di $ \v{B} $ diventa
\begin{align}
  &\nabla \times \langle \v{B} \rangle = 4\pi k_m (\v{J}_{\text{ext}} + \v{J}_{\text{pol}}) = 4\pi k_m (\v{J}_{\text{ext}} + \nabla \times \v{M})
  \\
  &\implies \nabla \times \left[ \langle \v{B}\rangle - 4\pi k_m \v{M}  \right] = 4\pi k_m \v{J}_{\text{ext}}
\end{align}
e definiamo il \textit{vettore di induzione magnetica} $ \v{H} $ come
\begin{align}
  \v{H} \coloneq \frac{\langle \v{B} \rangle}{4\pi k_m} - \v{M}
  \implies \nabla \times \v{H} = \v{J}_{\text{ext}}
\end{align}
Le equazioni di Maxwell mesoscopiche statiche che abbiamo ottenuto sono quindi:
\begin{align}
  &\nabla \cdot \v{D} = \rho_{\text{ext}}
  \\
  &\nabla \cdot \langle \v{B} \rangle = 0
  \\
  &\nabla \times \langle \v{E} \rangle = 0
  \\
  &\nabla \times \v{H} = \v{J}_{\text{ext}}
\end{align}

Nel caso dinamico dobbiamo prima ridefinire la corrente di polarizzazione $ \v{J}_{\text{pol}} $ perché dato che il vettore di polarizzazione $ \v{P} $ è 
legato alle cariche dalle relazioni 
\begin{align}
  \rho_{\text{pol}} &= - \nabla \cdot \v{P}
  \\
  \sigma_{\text{pol}} &= \v{P} \cdot \vers{n}
\end{align}
allora è chiaro che una variazione temporale $ \partial_t \v{P} $ deve essere legata in qualche modo alla corrente di polarizzazione, cosa che nel caso 
statico era nulla. Per farlo sfrutteremo l'equazione di continuità con il dettaglio che sorgenti esterne e sorgenti di risposta devono rispettare 
indipendentemente la loro equazione di continuità:
\begin{align}
  \der{\rho_{\text{ext}}}{t} + \nabla \cdot \v{J}_{\text{ext}} &= 0
  \\
  \der{\rho_{\text{pol}}}{t} + \nabla \cdot \v{J}_{\text{pol}} &= 0
\end{align}
che sono una semplice conseguenza del fatto che le equazioni di Maxwell valgono indipendentemente per le due sorgenti. Dalla seconda otteniamo
\begin{align}
  \der{(-\nabla\!\cdot\!\v{ P })}{t} + \nabla\!\cdot\!\v{ J }_{\text{pol}} = 0 
  \implies \v{J}_{\text{pol}} = \der{\v{P}}{t} \;\left( + \;\nabla\!\times\!\v{ M } \right)
\end{align}
quindi abbiamo trovato che alla corrente di polarizzazione trovata prima dobbiamo aggiungere il termine dovuto a $ \partial_t\v{P} $:
\begin{align}
  \v{J}_{\text{pol}} = \nabla \times \v{M} + \partial_t \v{P}
\end{align}
Possiamo quindi riscrivere l'equazione per il rotore di $ \v{H} $ con i termini aggiuntivi
\begin{align}
  \nabla \times \v{H} = \v{J}_{\text{ext}} + \der{\v{P}}{t} + \frac{1}{4\pi k_e} \der{\langle \v{E} \rangle}{t} = \v{J}_{\text{ext}} + \der{\v{D}}{t}
\end{align}
e quindi le equazioni di Maxwell dinamiche in scala mesoscopica diventano
\begin{align}
  &\nabla \cdot \v{D} = \rho_{\text{ext}}
  \\
  &\nabla \cdot \langle \v{B} \rangle = 0
  \\
  &\nabla \times \langle \v{E} \rangle = - \der{\v{B}}{t}
  \\
  &\nabla \times \v{H} = \v{J}_{\text{ext}} + \der{\v{D}}{t}
\end{align}
Le equazioni così ottenute sono ben diverse dalle equazioni microscopiche, tanto per cominciare non sono complete, dato che non abbiamo abbastanza relazioni 
per identificare 4 campi, e inoltre le relazioni che legano i campi di induzione rispetto a quelli microscopici, cioè $ \v{D}(\v{E,B}), \v{H}(\v{E,B}) $ 
sono in generale molto complicate. Dalle equazioni otteniamo però le seguenti leggi di conservazione in assenza di sorgenti esterne:
\begin{align*}
  D_\perp, B_\perp, E_\parallel, H_\parallel \qquad \text{si conservano.}
\end{align*}
Al fine di semplificare la trattazione nel seguito faremo un serie di approssimazioni, la prima consiste nel richieder la linearità della dipendenza dei 
campi di induzione in funzione dei campi microscopici, che è una approssimazione molto forte. La seconda riguarda la separazione della risposta cioè 
richiedere che $ \v{D} = \v{D}(\v{E}) $ e che $ \v{H} = \v{H}(\v{B}) $, che in realtà è una cosa piuttosto comune nei mezzi. Adesso per trattare i vari casi 
necessitiamo soltanto delle \textit{equazioni costitutive}, cioè le effettive relazioni che ci permettono di determinare i campi indotti. In generale 
la polarizzazione in funzione del campo elettrico prende la forma di un integrale 
\begin{align}
  P_i (\v{r},t) = \frac1{4\pi k_e} \int d^3r' \int dt' \chi_{i j}(\v{r, r'}, t, t') E_j (\v{r'}, t')
\end{align}
e in questa forma chiediamo che rispetti l'ipotesi di stazionarietà, cioè che dipenda dalla differenza $ t - t' $ e non in modo indipendente dal tempo, e 
l'ipotesi di mezzo omogeneo che consiste nel richiedere che la dipendenza solo da $ \v{r - r'} $. Abbiamo richiesto quindi l'invarianza traslazionale sia 
nel tempo che nello spazio e l'integrale prende la forma
\begin{align}
  P_i (\v{r},t) = \frac1{4\pi k_e} \int d^3r' \int dt' \chi_{i j}(\v{r - r'}, t - t') E_j (\v{r'}, t')
\end{align}
e a questo punto richiediamo anche altre due ipotesi cioè la località spaziale e temporale, la quale ci dice che la risposta del mezzo è legata a un punto 
preciso e a un istante, cosa che chiaramente non è vera ma è utile per dare una prima approssimazione del mezzo, che ci porta a un semplice coefficiente
\begin{align}
  \chi_{i j}(\v{r - r'}, t- t') &= \chi \; \delta( \v{r - r'}) \, \delta(t - t')
  \\
  \implies P_i (\v{r},t) &= \frac1{4\pi k_e} \chi_{i j} E_j (\v{r}, t)
\end{align}
ma finora $ \chi_{i j} $ è ancora un tensore, che può portare a dei risultati come la polarizzazione orientata in modo diverso rispetto al campo elettrico, 
quindi per ottenere il risultato più semplice in assoluto richiediamo che $ \chi_{i j} = \chi \, \delta_{ij} $ quindi si ottiene
\begin{align} \label{eq: relazione fra P e E}
  \v{P}(\v{r}, t) = \frac1{4\pi k_e} \chi \,\v{E}(\v{r}, t)
\end{align}
dove $ \chi $ è detta \textit{suscettività elettrica}, ed è una grandezza adimensionale. Il vettore di induzione elettrica diventa
\begin{align}
  &\v{D} = \eps_0 \v{E} + \v{P} = \eps_0 (1 + \chi) \v{E} \eqcolon \eps \v{E}  \tag*{[\textnormal{SI}]}
  \\
  &\v{D} = 4\pi \v{P} + \v{E} = \left( 1 + 4\pi \chi \right) \v{E} \eqcolon \eps \v{E} \tag*{[\textnormal{cgs}]}
\end{align}
dove abbiamo definito la \textit{costante dielettrica del mezzo} $ \eps $. Abbiamo inoltre omesso la notazione con il campo medio $ \langle \v{E} \rangle $ 
al fine di snellire la notazione, dunque da ora in poi supporremo sempre questa cosa, specificando invece quando si intende il campo microscopico.
Sempre riguardo la notazione, useremo i pedici $ 1,2 $ per indicare rispettivamente la parte reale e immaginaria di una quantità complessa, ad esempio 
$ \eps = \eps_1 + \eps_2 $, e quando è necessario precisare che una certa grandezza è complessa la indicheremo con una tilde: $ \tilde{\eps} $.

\section{Effetti di campo locale}

\begin{figure}[ht]
\centering
\begin{tikzpicture}[
  >=Stealth,
  dipolo/.style={->, thick, brown!70!black},
  Pvec/.style={->, thin, purple!70!black},
  leader/.style={dashed, gray, thin},
  every node/.style={font=\small}
]

% --- Regione esterna: continuo polarizzato ---
\fill[purple!15, rounded corners=4pt] (-4.2,-3.2) rectangle (4.2,3.2);

% Frecce di P sparse (continuo uniformemente polarizzato)
\draw[Pvec] (-3.6,-2.6) -- ++(0,0.5);
\draw[Pvec] (-3.6,-0.5) -- ++(0,0.5);
\draw[Pvec] (-3.6, 1.8) -- ++(0,0.5);
\draw[Pvec] (-2.2, 2.4) -- ++(0,0.5);
\draw[Pvec] ( 0.0, 2.6) -- ++(0,0.5);
\draw[Pvec] ( 2.2, 2.4) -- ++(0,0.5);
\draw[Pvec] ( 3.6, 1.8) -- ++(0,0.5);
\draw[Pvec] ( 3.6,-0.5) -- ++(0,0.5);
\draw[Pvec] ( 3.6,-2.6) -- ++(0,0.5);
\draw[Pvec] ( 2.2,-2.8) -- ++(0,0.5);
\draw[Pvec] ( 0.0,-2.8) -- ++(0,0.5);
\draw[Pvec] (-2.2,-2.8) -- ++(0,0.5);

% --- Sfera di Lorentz ---
\draw (0,0) circle (1.6);
\draw[dashed, thick] (0,0) circle (1.6);

% Dipoli discreti dentro la sfera (orientazioni varie)
\draw[dipolo] (-0.9, 0.6) -- ++(0.05,0.45);
\draw[dipolo] ( 0.8, 0.7) -- ++(0.10,0.45);
\draw[dipolo] (-0.8,-0.5) -- ++(-0.10,0.45);
\draw[dipolo] ( 0.9,-0.4) -- ++(0.10,0.45);
\draw[dipolo] ( 0.0, 1.0) -- ++(0.05,0.40);
\draw[dipolo] (-1.1, 0.0) -- ++(0.00,0.45);
\draw[dipolo] ( 1.1, 0.0) -- ++(0.00,0.45);
\draw[dipolo] (-0.3,-1.0) -- ++(-0.05,0.45);
\draw[dipolo] ( 0.4,-1.0) -- ++(0.05,0.45);

% --- Punto centrale: la molecola di interesse ---
\fill[red!80!black] (0,0) circle (1.8pt);

% --- Etichette con linee guida ---
\draw[leader] (1.35, 1.0) -- (4.2, 2.6);
\node[anchor=west] at (4.2, 2.6) {Continuo, $ \v{E}_P $};

\draw[leader] (0.7, 0.5) -- (4.2, 1.0);
\node[anchor=west] at (4.2, 1.0) {Dipoli discreti, $\mathbf{E}_{\text{micro}}$};

\draw[leader] (0, 0) -- (4.2,-0.6);
\node[anchor=west] at (4.2,-0.6) {Punto di interesse, $ \v{E}_\textnormal{loc} $};

\end{tikzpicture}
\caption{Sfera mesoscopica con cui modelliamo gli effetti di campo locale: tutto il mezzo è modellizzato tramite polarizzazione dando 
il campo $ \v{E}_P $, mentre dentro abbiamo il contributo dei dipoli che contiamo come $ \v{E}_{\text{micro}} $.}
\label{figure:_effetti_campo_locale}
\end{figure}

Per trovare gli effetti di un campo macroscopico su un punto (atomo) dobbiamo considerare anche gli effetti che la polarizzazione degli altri atomi produce 
sul punto considerato: usiamo il modello descritto in figura \ref{figure:_effetti_campo_locale} dove il campo nel punto verrà indicato 
con $ \v{E}_{\text{locale}} $, e considerando una sfera mesoscopica di raggio $ R $ si ha 
\begin{align}
  \v{E}_{\text{loc}} = \v{E} - \v{E}_P + \v{E}_{\text{micro}}
\end{align}
dove $ \v{E}_{\text{P}} $ è il campo dovuto alla polarizzazione della sfera, $ \v{E}_{\text{micro}} $ è quello dovuto ai dipoli vicini. 
In sostanza togliamo dal campo medio $ \v{E} $ il contributo dei dipoli di tutto il mezzo che descriviamo come un campo generato da una polarizzazione, cioè 
$ \v{E}_P $, riaggiungendo (perché erano compresi in $ \v{E}_P $) il contributo dovuto ai dipoli dentro la sfera cioè $ \v{E}_{\text{micro}} $.
È facile vedere che il contributo dei dipoli vicini si annulla per simmetria, mentre il campo di una sfera polarizzata uniformemente è
\begin{align}
  \v{E}_P = - \frac{4\pi k_e}{3} \v{P}
\end{align}
e il motivo perché consideriamo la sfera con polarizzazione uniforme viene dal risolvere il problema della sfera in campo elettrico uniforme. Quindi 
abbiamo che 
\begin{align}
  \v{E}_{\text{loc}} = \v{E} + \frac{4\pi k_e}{3} \v{P}
\end{align}
che ci permette di trovare il dipolo 
\begin{align}
  \v{d} = \eps_0 \alpha \v{E}_{\text{loc}}
\end{align}
dove $ \alpha $ è la \textit{polarizzabilità atomica} (supponiamo che segua questa relazione lineare). La polarizzazione diventa, con $ N $ la densità di 
elettroni,
\begin{align}
  \v{P} = \eps_0 N \alpha \v{E}_{\text{loc}} = N\alpha (\v{E} + \frac{4\pi k_e}{3} \v{P}) 
  \implies \v{P} = \frac{N \alpha}{1 - \frac{N \alpha}{3} }\, \eps_0 \v{E} \eqcolon \eps_0 \chi \v{E}
\end{align}
Confrontiamo la relazione appena ottenuta con 
\begin{align}
  \v{P} = \eps_0 (\eps_r - 1)\v{E}
\end{align}
e otteniamo che la costante dielettrica relativa rispetta la relazione
\begin{align}
  \chi = \eps_r - 1 = \frac{N \alpha}{ \left( 1 - \frac{ N \alpha}{3} \right)} \implies \frac{\eps_r - 1}{\eps_r + 2} = \frac{n\alpha}{3}
  \tag*{[Clausius-Mossotti]}
\end{align}
nota come \textit{relazione di Clausius-Mossotti}. Questa relazione ci dà una modo operativo per trovare la relazione costitutiva, dovendo soltanto 
trovare $ N $ (relativamente facile) e $ \alpha $, che come vedremo nella sezione \ref{section:_Drude-Lorentz} richiederà un modello per descrivere 
la risposta atomica. L'equazione di Clausius-Mossotti inserisce le correzioni al campo macroscopico dovute a una polarizzazione, dunque il campo locale 
sarà tanto più diverso dal campo macroscopico quanto più sono importanti gli effetti di polarizzazione, che si traduce in tanto più è grande la 
densità di dipoli $ N $. Per questo ragionamento, la relazione che abbiamo ricavato si dice che vale per \textit{mezzi densi}, proprio perché nei mezzi 
rarefatti come i gas la densità $ N $ è molto più piccola rispetto a liquidi e solidi.

\section{Energia e leggi di conservazione}
Il lavoro del campo elettromagnetico sulle cariche è dovuto solo al campo elettrico e vale 
\begin{align}
  W = \int_V \v{J \cdot E} \,d^3x
\end{align}
e, per ottenere una analoga legge di conservazione come quella trovata nel vuoto, sfruttiamo le equazioni di Maxwell macroscopiche per sostituire  $ \v{J} $:
\begin{align}
  W = \int_V \v{J \cdot E} \,d^3x = \int_V \v{E} \cdot \left( \nabla\!\times\!\v{ H } \right) - \v{E} \cdot \der{\v{D}}{t} \, d^3x.
\end{align}
Adesso applichiamo l'identità vettoriale
\begin{align}
  \nabla \cdot \left( \v{E \times H} \right) = \v{H} \cdot \left( \nabla\!\times\!\v{ E } \right)
  - \v{E} \cdot \left( \nabla\!\times\!\v{ H } \right) 
\end{align}
e otteniamo
\begin{align}
  \int_V \v{J \cdot E} \,d^3x &= \int_V \v{H} \cdot \left( \nabla\!\times\!\v{ E } \right) - \nabla \cdot \left( \v{E \times H} \right)
  - \v{E} \cdot \der{\v{D}}{t} \, d^3x 
  \\
  &= - \int_V \v{H} \cdot \der{\v{B}}{t} + \nabla \cdot \left( \v{E \times H} \right)
  + \v{E} \cdot \der{\v{D}}{t} \, d^3x 
\end{align}
dove abbiamo usato l'equazione per il rotore di $ \v{E} $. Per procedere oltre assumiamo che il mezzo sia lineare e che non ci siano effetti di 
assorbimento, in questo modo possiamo riconoscere nell'espressione che abbiamo ottenuto la derivata di una quantità $ u $ che definiremo come la 
\textit{densità di energia elettromagnetica}:
\begin{align}
  u = \frac{1}{2} \left( \v{E \cdot D} + \v{B \cdot H}  \right) \implies \der{u}{t} = \v{E} \cdot \der{\v{D}}{t} + \v{B} \cdot \der{\v{H}}{t} 
\end{align}
e nel caso $ \v{D} = \eps \v{E}, \v{H} = \frac{\v{B}}{\mu} $ allora 
\begin{align}
  u = \frac{\eps}{2}\v{E}^2 + \frac{1}{2\mu} \v{B}^2.
\end{align}
Manca solo da definire il nuovo vettore di Poynting:
\begin{align}
  \v{S} = \v{E \times H}
\end{align}
e quindi la legge di conservazione che avevamo scritto prende il nome di \textit{teorema di Poynting nella materia} ed è in forma globale
\begin{align}
  \int_A \v{S} \cdot d\v{A} + \int_V \der{u}{t}\, d^3x + \int_V \v{J \cdot E} \,d^3x = 0.
\end{align}
Notiamo che il teorema di Poynting possiamo definirlo adesso anche su un altra scala: in quella mesoscopica useremo il nuovo vettore di Poynting 
$ \v{S} = \v{E \times H} $ e la nuova densità di energia e il teorema è lo stesso solo che le quantità in gioco sono una media mesoscopica di quelle 
microscopiche. Tuttavia se volessimo invece definire un teorema di Poynting microscopico, allora sarebbe esattamente quello definito 
\ref{eq:teorema_di_Poynting_locale} e dovremmo in prima istanza considerare tutte le piccolissime correnti, cosa che si può approcciare con il metodo 
sviluppato nella sezione precedente studiando gli effetti locali tramite un modello di una sfera mesoscopica.


\subsection{Atomo di Thomson}


\section{Interazione luce-materia}
Nei paragrafi precedenti abbiamo visto come poter descrivere i campi elettromagnetici su scala mesoscopica, e facendo varie ipotesi semplificative sulla 
natura della relazione costitutiva che lega $ \v{P} $ e $ \v{E} $. Adesso vogliamo capire cosa succede se i campi variano rapidamente nella nostra zona di 
interesse.

Nella relazione 
\begin{align} \label{eq:_integrale_di_convoluzione}
  \v{P}(\v{r}, t) = \eps_0 \int_{-\infty}^{+\infty} \chi( t-t') \v{E}(\v{r}, t') \, dt'
\end{align}
non possiamo più imporre che $ \chi( t - t') = \delta( t - t') $ e quindi che la risposta sia istantanea, perché i campi variano molto rapidamente. Dobbiamo 
quindi sistemare questo integrale, pur mantenendo il principio di causalità, cioè che la risposta può dipendere soltanto dal campo a un tempo 
$ \v{E}( t > t') $, e possiamo farlo imponendo che
\begin{align}
  \chi( t < t') = 0.
\end{align}
Adesso è conveniente passare in trasformata di Fourier, e i motivi sono due: in generale il campo $ \v{E} $ nel dominio del tempo può essere molto 
complicato, passando in trasformata lavoriamo con le sue componenti monocromatiche per poi ricomporre il campo facendo un integrale; il secondo motivo è 
che riconosciamo che l'integrale in \eqref{eq:_integrale_di_convoluzione} è detto un integrale di convoluzione, e la trasformata di Fourier ha la proprietà 
di trasformare integrali di convoluzione in prodotti. Ci aspettiamo di ottenere un risultato del tipo 
\begin{align}
  \v{P}_\omega(\v{r}) = \eps_0 \chi_\omega \v{E}_\omega(\v{r})
\end{align}

\section{Modello di Drude-Lorentz} \label{section:_Drude-Lorentz}
Al fine di poter descrivere le interazioni fra la radiazione elettromagnetica con vari tipi di materia, abbiamo bisogno della relazione fra polarizzazione 
indotta e campo elettrico, cioè la $ g_\omega $ vista precedentemente, che possiamo ricavare dall'esperimento in modo empirico, oppure attraverso una 
modellizzazione dei costituenti della materia, e in particolare degli elettroni.

Il modello di Drude-Lorentz considera l'elettrone come una particella sottoposta a 3 tipi di forze: una di richiamo, una dissipativa e una forzante che per 
noi sarà la forza di Coulomb dovuta al campo elettrico di una onda (piana) incidente. 
Questo si realizza in una equazione del moto per una particella di carica $ q $ e massa $ m $ in una equazione del moto del tipo 
\begin{align}
  m \ddot{\v{r}} = - \omega_0^2 \v{r} - \gamma m \dot{\v{r}} + q \v{E}(\v{r}, t)
\end{align}
Per una onda piana abbiamo $ \v{E}(\v{r} = 0 , t) = E_0 \vers{z} e^{-i\omega t} $ che abbiamo ipotizzato lungo $ z $ (cosa che non cambia la trattazione, 
a regime la direzione delle oscillazioni sarà la stessa del campo elettrico), supponiamo una soluzione dello stesso tipo cioè $ \v{z} = z_0 e^{-i\omega t} $ 
e si ottiene che 
\begin{align}
  z_0 = \frac{ qE_0 / m}{\omega_0^2 - \omega^2 - i \gamma \omega} 
\end{align}
quindi il dipolo risultante sarà
\begin{align}
  \v{d} = d_0 \vers{z} e^{-i\omega t} = q z_0 \vers{z} e^{-i \omega t}
  \\
  \implies \v{P}(\v{r} , t) = \frac{n \eps_0 q^2 /m}{\omega_0^2 - \omega^2 - i \gamma \omega} \v{E}(\v{r} , t)
\end{align}
dove $ n $ è il numero di cariche $ q $ in un volumetto (densità volumetrica di elettroni). È possibile notare che eravamo partiti dal calcolare le 
oscillazioni nell'origine e finiti ottenendo una dipendenza spaziale della polarizzazione, ma questo può essere giustificato dicendo che quello che abbiamo 
ottenuto è in realtà la trasformata di Fourier $ \v{P}_\omega $, cosa che avremmo potuto fare all stesso modo aggiungendo $ \v{P}_{k,\omega} $, ma $ \v{r} $ 
non svolge alcun ruolo nel nostro modello se non quello di selezionare l'atomo che consideriamo. Inoltre per l'ipotesi di isotropia la risposta deve essere 
uguale in ogni punto quindi è normale che la polarizzazione erediti la dipendenza spaziale del campo elettrico. Il nostro $ \chi(\omega) $ è quindi
\begin{align}
  \chi(\omega) = \frac{n q^2 /m}{\omega_0^2 - \omega^2 - i \gamma \omega}
\end{align}
Definiamo la \textit{frequenza di plasma} $ \omega_p $ come
\begin{align}
  \omega^2_p \coloneq \frac{n q^2}{\eps_0 m} \implies \v{P}(\v{r}, t) = \eps_0 \chi(\omega )\v{E}(\v{r}, t)
\end{align}
e quindi la $ \eps(\omega) $ definita come 
\begin{align}
  \eps_r(\omega) = 1 + \frac{g_\omega}{\eps_0} = 1 + \chi_\omega = 1 + \frac{\omega_p^2}{\omega_0^2 - \omega^2 - i \gamma \omega}
\end{align}
Essendo complesso la dividiamo in parte reale e parte immaginaria, che rappresentano rispettivamente il termine dispersivo e il termine dissipativo:
\begin{align}
  \eps_{r1} &= 1 + \frac{\omega_p^2 (\omega_0^2 - \omega^2) }{(\omega_0^2 - \omega^2)^2 + \gamma^2 \omega^2}  \tag*{dispersivo}
  \\
  \eps_{r2} &= \frac{\omega_p^2 \gamma \omega }{(\omega_0^2 - \omega^2)^2 + \gamma^2 \omega^2} \tag*{dissipativo}
\end{align}

\begin{figure}
    \center
    \begin{tikzpicture}
        \def \a {8}
        \def \b {4}

        \def \gamma {0.8}
        \def \omegaO {3}
        \def \omegaP {3}

        \draw[->] (0,-0.5) -- (0,\b) node[left] {$ \eps $};
        \draw[->] (-0.5,0) -- (\a,0) node[below] {$ \omega $};

        \draw[domain=0:\a,samples = 500, variable=\x, thick, blue]
            plot ({\x},
                {1 + \omegaP*\omegaP*(\omegaO*\omegaO - \x*\x)/(((\omegaO*\omegaO - \x*\x)*(\omegaO*\omegaO - \x*\x)) + \gamma*\gamma*\x*\x)});

        \draw[ domain = 0:\a,samples = 500, variable = \x,thick, orange]
            plot ({\x}, 
                {(\omegaP*\omegaP*\gamma*\x)/(((\omegaO*\omegaO - \x*\x)*(\omegaO*\omegaO - \x*\x)) + \gamma*\gamma*\x*\x)});

        \node[ below left ] at (\omegaO,0) {$ \omega_0 $};
        \node[ left ] at (0,1) {$ 1 $};
        \draw[domain= -0.05:\a, dashed, variable = \x] plot ({\x}, 1);

    \end{tikzpicture}
    \caption{Andamento di $ \eps_1 $ e $ \eps_2 $.}
    \label{figure:_dispersione_dissipazione}
\end{figure}

Il $ \chi(\omega) $ che abbiamo appena trovato grazie a Drude-Lorentz è esattamente il coefficiente (dipendente dalla frequenza $\omega$) che ci dà la 
proporzionalità fra il campo elettrico e la polarizzazione indotta, dunque è l'$\alpha$ che ci serviva nella relazione di Clausius-Mossotti.

\section{Zone di frequenza} % da cambiare il titolo
Nella sezione precedente abbiamo ricavato che la funzione dielettrica nel modello Drude-Lorentz ha parte reale e parte immaginaria, ognuna relativa 
a un fenomeno fisico diverso, dispersione e assorbimento, con grafici:
% grafico dispersione e assorbimento
Possiamo distinguere il dominio delle frequenze in tre zone dette I, II, III.

\textbf{Zona I}. In questa zona trascuriamo sia l'assorbimento, sia la dipendenza dalla frequenza di $ \eps(\omega) = \eps \in \R  $, cioè è il limite per 
$ \omega \ll \omega_0 $. Questa zona non presenta differenze con l'elettrodinamica nel vuoto, se non per la relazione di dispersione dovuto al fatto che 
l'equazione di d'Alembert che si ottiene è
\begin{align}
  \nabla^2 \v{E} - {\mu \eps} \der[2]{\v{E}}{t} = 0
\end{align}
dunque abbiamo che il numero d'onda $ k $ vale
\begin{align}
  k = \sqrt{\mu_r \eps_r} \,\frac{\omega}{c}
\end{align}
e definiamo l'indice di rifrazione $ n $ come:
\begin{align}
  n \coloneq \sqrt{\mu_r \eps_r} \implies k = \frac{n \omega}{c}
\end{align}
La velocità di fase e la velocità di gruppo coincidono essendo una relazione di dispersione lineare, ma a differenza del vuoto sono soppresse del fattore 
$ n $:
\begin{align}
  v_f = \frac{\omega}{k} = \frac{c}{n}
\end{align}


\textbf{Zona II}. In questa zona torniamo a considerare la dipendenza dalla frequenza ma trascuriamo sempre l'assorbimento, quindi 
$ \eps(\omega) = \eps_1(\omega) \in \R $, partiamo dalla relazione di dispersione:
\begin{align}
  k = \frac{\omega}{c} \sqrt{\eps_r(\omega)} \simeq \frac{\omega}{c} \sqrt{1 + \frac{\omega_P^2}{\omega_0^2 - \omega^2 }} \coloneq \frac{\omega}{c} n(\omega)
\end{align}
e quindi la velocità di fase è
\begin{align}
  v_f = \frac{c}{n(\omega)} = \frac{c}{\sqrt{1 + \frac{\omega_P^2}{\omega_0^2 - \omega^2 }}}
\end{align}
Notiamo che in questa espressione il denominatore $ \omega_0^2 - \omega^2 $ diventa negativo per $ \omega > \omega_0 $, il che ci porta al risultato 
assurdo di $ v_f > c $, cioè che la velocità dell'onda nel mezzo diventa superluminale e per questo motivo la velocità di fase perde di significato in 
questa zona. Un altro motivo che ci spinge a cercare una nuova velocità da associare alle onde nel mezzo è il fatto che adesso la velocità di fase varia in 
base alla frequenza, dunque ci serve una velocità che descriva un pacchetto di onde che si muove nel mezzo. È chiaro che la velocità di gruppo $ v_g $ 
risolve il problema, dalla definizione abbiamo 
\begin{align}
  v_g = \frac{d\omega}{dk} = \frac{1}{\frac{dk}{d\omega}} = \frac{c}{n(\omega) + \omega \frac{dn}{d\omega}}
\end{align}

\textbf{Zona III}: mezzo assorbente e dispersivo. In questa zona la funzione dielettrica è in generale complessa: 
$ \eps(\omega) = \eps_1(\omega) + i \eps_2(\omega) $, tuttavia consideriamo sempre mezzi non magnetici, dunque 
$ \mu(\omega) = \mu_0 $ e $ \v{B} = \mu_0 \v{H} $. Dalla definizione il vettore di induzione elettrica vale 
$ \v{D}(\omega) = \tilde{\eps}(\omega) \v{E}(\omega) $ e quindi usiamo le equazioni di Maxwell per ricavare la relazione di dispersione:
\begin{gather}
  \nabla \times \nabla \times \v{E} = \nabla \times \left( i\omega \v{B} \right) = i\omega \mu_0 \nabla \times \v{H} 
  = \omega^2 \mu_0 \tilde{\eps}(\omega) \v{E}
  \\
  \implies \nabla^2 \v{E} + \omega^2 \mu_0 \tilde{\eps}(\omega) \v{E} = 0
\end{gather}
considerando un andamento dei campi del tipo $ f(\v{r},t) \sim e^{i(\v{k \cdot r} - \omega t)}$ dalla precedente equazione otteniamo
\begin{align}
  i \v{k} \cdot (i \v{k}) \v{E} + \omega^2 \mu_0 \tilde{\eps}(\omega) \v{E} = 0 \implies \v{k \cdot k} = \omega^2 \mu_0 \tilde{\eps}(\omega) 
  = \frac{\omega^2}{c^2} \tilde{\eps}_r(\omega)
\end{align}
che è la relazione di dispersione che cercavamo. Da questa osserviamo che $ \v{k} $ deve essere necessariamente complesso e lo indicheremo come
$ \v{k} = \v{k}_1 + i\v{k}_2 $ e quindi l'andamento dei campi può essere riscritto nella forma 
$ f(\v{r},t) \sim e^{i(\v{k_1 \cdot r} - \omega t)}\cdot e^{-\v{k}_2 \cdot \v{r}} $. Sviluppiamo la relazione di dispersione:
\begin{align}
  (\v{k}_1 + i\v{k}_2)(\v{k}_1 + i\v{k}_2) &= k_1^2 -k_2^2 + 2 i\v{k}_1  \cdot \v{k}_2 
  = \frac{\omega^2}{c^2} \eps_{r1}(\omega) + i \frac{\omega^2}{c^2} \eps_{r2}(\omega)
  \\
  &\begin{cases}
    k_1^2 -k_2^2 = \frac{\omega^2}{c^2} \eps_{r1}(\omega)
    \\
    2\v{k}_1  \cdot \v{k}_2 = \frac{\omega^2}{c^2} \eps_{r2}(\omega).
  \end{cases}
\end{align}
Dal sistema che abbiamo ottenuto, dato che $ \v{k}_1 \neq 0 $ perché altrimenti non si propagherebbe, osserviamo che se $ \eps_{r2} \neq 0 $ allora 
$ \v{k}_2 \neq 0 $, ma non vale il viceversa dato che per $ \v{k}_2 \perp \v{k}_1 \implies \eps_{r2} = 0 $. Riscrivendo la relazione di dispersione con 
l'indice di rifrazione $ \tilde{n}^2 = n_1 + in_2 = \eps_{r1} + i \eps_{r2} $ si ottiene
\begin{align}
  \begin{cases}
    k_1 = \frac{\omega^2}{c^2} n_1
    \\
    k_2 = \frac{\omega^2}{c^2} n_2
  \end{cases}
  \quad \textnormal{con} \quad 
  \begin{cases}
    n_1^2 -n_2^2 = \eps_{r1}(\omega)
    \\
    2n_1 n_2 = \eps_{r2}(\omega).
  \end{cases}
\end{align}

Precedentemente avevamo definito la polarizzabilità atomica $ \alpha(\omega) $ come il coefficiente di proporzionalità (a meno di un $ \eps_0 $) fra il 
campo elettrico locale e il dipolo indotto, cioè
\begin{align}
  \v{d} = \eps_0 \alpha(\omega) \v{E}_{\text{loc}},
\end{align}
la quale è legata alla suscettibilità elettrica dalla densità volumetrica di dipoli $ N $ come 
\begin{align}
  \chi(\omega) = N \alpha(\omega).
\end{align}
Scriviamo l'indice di rifrazione in termini della polarizzabilità:
\begin{align}
  n = \sqrt{\eps_r} = \sqrt{1 + N\alpha} = \sqrt{1 + \frac{\omega_p^2}{\omega_0^2 - \omega^2 - i \gamma \omega}}.
\end{align}
Definiamo il \textit{coefficiente di assorbimento} $ \beta $ come 
\begin{align}
  \beta \coloneq 2 k_2 = 2 n_2 \frac{\omega}{c}
\end{align}
e in questo modo possiamo esprimere l'onda in un mezzo assorbente che assumiamo propagarsi lungo $ z $ come
\begin{align}
  \v{E} = \v{E}_0 e^{i (k_1 z - \omega t)} e^{ - \beta \frac{z}{2}}
\end{align}

\section{Onde nei conduttori}
Il modello del conduttore consiste nell'eliminare la frequenza di risonanza perché consideriamo gli elettroni non più come legati armonicamente al nucleo 
ma liberi di muoversi all'interno del mezzo, sottoposti soltanto alla dissipazione e alla forza di Lorentz. La dissipazione in questo caso è dovuta, 
oltre che all'irraggiamento, agli urti con gli atomi del conduttore nel muoversi. Proprio questi urti sono il motivo per il quale gli elettroni, in un 
conduttore sottoposto a campo elettrico costante, non subiscono una accelerazione costante ma nel complesso raggiungono una velocità detta 
\textit{velocità di deriva} (\textit{drift velocity}) $ v_d $. Per ricavare tale velocità consideriamo l'equazione del moto (in una dimensione) di un 
elettrone sottoposto a un campo elettrico costante lungo $ x $ che è
\begin{align}
  m_e \ddot{x} = - e E = m_e a_x
\end{align}
e introducendo un parametro $ \tau $ che rappresenta un tempo caratteristico fra un urto e il successivo si ottiene che la velocità di drift è
\begin{align}
  v_d = a_x \tau = - \frac{e \tau E}{m_e}.
\end{align}
Questa affermazione è valida perché: $ v_d $ rappresenta una velocità caratteristica del flusso di elettroni; discretizzando l'equazione fra un urto e 
l'altro, per $ \tau $ abbastanza piccoli (e lo è) si ha che il moto accelerato a tratti degli elettroni può essere descritto da una equazione del moto con 
una dissipazione, tuttavia questo modello non aggiunge nulla alla fisica del sistema e porta allo stesso risultato; l'unica velocità che si poteva ottenere 
dai parametri che avevamo è proprio $ v_d $. La densità di corrente dovuta al moto di  questi elettroni si ricava come
\begin{align}
  J = \frac{N_e \,q\, v_d}{V} = \frac{N e^2 \tau }{m_e} E \eqcolon \sigma_0 E
\end{align}
dove $ N_e $ è il numero di elettroni nel volume $ V $, $ N $ il loro rapporto (densità di elettroni) e abbiamo definito $ \sigma_0 $ come il coefficiente 
fra il campo elettrico e la corrente che genera, che è detta \textit{conducibilità statica}. Quella che abbiamo ricavato è la \textit{legge di Ohm} 
microscopica che lega la densità di corrente al campo elettrico che la induce. Per studiare il caso dinamico applichiamo il modello di Drude-Lorentz 
(in realtà solo di Drude perché la parte della risonanza si riconduce a Lorentz) imponendo $ \omega_0 = 0 $:
\begin{align}
  m \ddot{\v{r}} = - m \gamma \dot{\v{r}} + q \v{E}(\v{r}, t).
\end{align}
Cerchiamo prima una soluzione per la velocità: dato che il campo elettrico ha 
un andamento temporale del tipo $ e^{-i\omega t} $, supponiamo lo stesso andamento per la velocità e otteniamo
\begin{align}
  -i\omega \v{v} = -\gamma \v{v} + \frac{q}{m} \v{E} \implies \v{v} = \frac{q/m}{\gamma - i \omega} \v{E}.
\end{align}
In analogia col caso statico, leghiamo $ \gamma $ al tempo caratteristico come $ \gamma = 1/\tau $, e ci ricaviamo una nuova conducibilità moltiplicando 
per $ N q $, ottenendo 
\begin{align}
  \v{v} = \frac{q \tau/m}{1 - i \omega \tau } \v{E} \quad \textnormal{e} \quad \sigma(\omega) = \frac{{N q^2 \tau}/{m}}{1 - i \omega\tau} 
  = \frac{\sigma_0 }{1 - i \omega \tau}
\end{align}
dove abbiamo riconosciuto la conducibilità statica. Questa nuova conducibilità dipendente dalla frequenza è complessa e quindi la dividiamo in parte reale 
e immaginaria:
\begin{align}
  \sigma = \sigma_1 + i \sigma_2 = \frac{\sigma_0}{1 + \omega^2 \tau^2}\left( 1 + i \omega \tau \right)
\end{align}
dove la parte reale è detta \textit{dissipativa} perché responsabile dell'energia ceduta alle cariche, mentre quella immaginaria è la risposta inerziale 
degli elettroni quindi relativa allo sfasamento.

Risolvendo invece per la posizione $ \v{r} $ otteniamo esattamente il modello di Drude-Lorentz con $ \omega_0 = 0 $, leghiamo con la polarizzazione 
$ \v{P} = N q \v{r} $ e si ottiene la funzione dielettrica relativa $ \eps_r(\omega) $:
\begin{align}
  \eps_r(\omega) = 1 - \frac{\omega_p^2}{\omega^2 + i \gamma \omega} \;\; \textnormal{con} \;  \omega_p^2 = \frac{N e^2 }{\eps_0 m_e } 
\end{align}
che leghiamo alla conducibilità complessa appena definita come 
\begin{gather}
  \eps_r( \omega) = 1 + \frac{i \sigma(\omega)}{\eps_0 \omega}, 
  \\
  \textnormal{con} \;
  \begin{cases}
    \DS \eps_{r,1} = 1 - \frac{\omega_p^2}{\omega^2 + \gamma^2} 
    \\
    \DS \eps_{r,2} = \frac{\omega \gamma \omega_p^2}{\omega^4 + \gamma^2\omega^2}
  \end{cases}  
\end{gather}
dove abbiamo diviso in parte reale e immaginaria. Per comodità riportiamo la relazione tra frequenza di plasma e conducibilità statica:
\begin{align}
  \omega_p^2 = \frac{\sigma_0}{\eps_0 \tau}.
\end{align}
Di seguito analizzeremo alcuni regimi di queste di questo modello.

\textbf{Basse frequenze}. Consideriamo il limite di basse frequenze la condizione 
\begin{align}
  \omega \tau \ll 1
\end{align}
che ci porta a una conducibilità reale (niente sfasamento) pari alla conducibilità nel caso statico $ \sigma_0 $. L'indice di rifrazione in questo 
regime vale 
\begin{align}
  n^2 \simeq 1 - \frac{i \sigma_0}{\eps_0 \omega} 
\end{align}
dove con l'approssimazione di buon conduttore $ \sigma_0 \gg \eps_0 \omega $ l'indice di rifrazione è quali solo immaginario, il che ci porta a un vettore 
d'onda pari a 
\begin{align}
  \v{k}^2 = \frac{\omega^2}{c^2} \, i \frac{\sigma_0}{\eps_0 \omega} \implies |\v{k}| = \sqrt{\frac{\omega \mu_0 \sigma_0}{2}} (1 + i)
\end{align}
che è proprio il risultato a cui eravamo arrivati nello studiare i le cavità in regime di campi lentamente variabili. Infatti la lunghezza pelle è di nuovo 
\begin{align}
  \delta = \sqrt{\frac{2}{\omega \mu_0 \sigma_0}}, 
\end{align}
che rappresenta la lunghezza caratteristica con la quale penetrano i campi, infatti all'interno abbiamo che 
\begin{align}
  \v{E} = \v{E}_0 e^{i (\v{k \cdot r} - \omega t)} \simeq \v{E}_0 e^{i (k x - \omega t)} \cdot e^{-x/\delta}.
\end{align}
Questo range di frequenze corrisponde nel rame a $ f \ll \SI{4e12}{\Hz} = \SI{4}{\THz} $ e quindi copre dalla continua fino all'infrarosso lontano, 
diminuendo via via la penetrazione (aumenta l'effetto pelle) all'aumentare della frequenza.

\textbf{Alte frequenze}. Stavolta la condizione è opposta, cioè
\begin{align}
  \omega \tau \gg 1
\end{align}
che ci porta una conducibilità puramente immaginaria:
\begin{align}
  \sigma(\omega) = i \frac{\sigma_0}{\omega \tau} \implies n^2 = 1  - \frac{\sigma_0}{\eps_0 \tau \omega^2} = 1 - \frac{\omega_p^2}{\omega^2}.
\end{align}
Da questo risultato possiamo distinguere altri due regimi: quando la frequenza è minore della frequenza di plasma, $ \omega < \omega_p $, $ \eps_r $ è 
negativo che ci porta a un vettore d'onda $ \v{k} $ puramente immaginario, di conseguenza il campo all'interno del mezzo decade soltanto senza fare 
oscillazioni come
\begin{gather}
  k_2 = \frac{1}{c} \sqrt{\omega_p^2 - \omega^2} \coloneq \frac{1}{\delta}
  \\
  \v{E} = \v{E}_0 e^{-i\omega t} \cdot e^{-x/\delta}.
\end{gather}
In questo caso l'onda è detta \textit{evanescente} e siamo in un regime di riflessione totale (il coefficiente di riflettanza che definiamo nella sezione 
\ref{sec: onde incidenti} è uguale a $ 1 $), le frequenze di cui stiamo parlando (sempre nel rame) sono comprese fra 
$ \SI{4e12}{\Hz} \ll f \ll \SI{2.6e15}{\Hz} $ coprendo quindi l'infrarosso medio, il visibile e l'ultravioletto vicino. Proprio per questo motivo la 
lucentezza è tipica dei metalli (nel visibile).

Invece nel caso in cui $ \omega > \omega_p $ abbiamo che l'indice di rifrazione si avvicina sempre più ad $ 1 $, e il vettore d'onda è reale. Dunque 
l'onda si propaga senza attenuazione dove si dice che il metallo diventa trasparente. Se non si trascura completamente lo smorzamento, in realtà 
anche in questo regime compare una componente dissipativa al secondo ordine in $ \gamma $: sviluppiamo in Taylor la funzione dielettrica per 
$ \gamma \ll \omega $ e otteniamo
\begin{align}
  \eps_r(\omega) = 1 - \frac{\omega_p^2}{\omega^2} + i \, \gamma \,\frac{ \omega_p^2 }{\omega^3} + O(\gamma^2)
\end{align}
che ci porta a un vettore d'onda che vale
\begin{align}
  k_1 &\simeq \frac\omega{c} \sqrt{\eps_1} = \frac{1}{c} \sqrt{\omega^2 - \omega_p^2}
  \\
  k_2 &\simeq \frac\omega{c} \sqrt{\frac{\eps_2}{4\eps_1}} = \frac{\omega_p^2 \gamma }{2c\omega \sqrt{\omega^2 - \omega_p^2}}
\end{align}
Questo regime di frequenze corrisponde a $ f \gg \SI{2.6e15}{\Hz} = \SI{2.6}{\peta \hertz} $ cioè $ \lambda \ll \SI{114}{\nm} $, che corrispondono 
all'ultravioletto estremo. Il limite del modello di Drude per i metalli è quando la frequenza si avvicina all'energia necessaria per la promozione 
degli elettroni di valenza, che per Einstein risulta essere 
\begin{align}
  E = \hbar \omega = h f.
\end{align}

Un esempio di applicazione del modello fuori dall'ambito dei conduttori è quello della ionosfera: modellizziamo la ionosfera come un gas trasparente 
di elettroni liberi con la differenza rispetto ai metalli riguardo la densità di portatori di carica $ N $ che arriva a 
$ \SIrange{e11}{e12}{\per\cubic\metre} $ contro i $ \SIrange{e28}{e29}{\per\cubic\metre} $ dei metalli. Questo fa sì che la frequenza di plasma 
della ionosfera arrivi nelle onde radio:
\begin{align}
  {f_p}_{\text{iono}} = \SIrange{3}{9}{\MHz}
\end{align}
e quindi le onde con una frequenza inferiore a quella di plasma vengono completamente riflesse.

\section{Relazioni di Kramers-Kronig}
In questa sezione introdurremo dei risultati di Analisi Complessa al fine di ricavare delle relazioni che legano la parte reale del $ \chi $ in 
equazione \eqref{eq: relazione fra P e E} con la sua parte immaginaria. Tuttavia queste relazioni non esistono "a priori" perché senza applicare 
particolari ipotesi, $ \chi_1 $ e $ \chi_2 $ sono indipendenti. Non lo sono più quando abbiamo imposto l'ipotesi di causalità cioè che la risposta 
deve essere successiva alla causa, che abbiamo tradotto in $ \chi(t - t') = 0 $ per $ t < t' $. La trasformata di Fourier della $ \chi(t) $ è 
(con la convezione di dividere successivamente per $ 2\pi $)
\begin{align}
  \chi(\omega) = \int_{-\infty}^{+\infty} \chi(\tau) e^{i\omega \tau} \,d\tau
\end{align}
e se scriviamo l'esponenziale in forma trigonometrica si ottiene
\begin{align}
  \chi(\omega) = \int_{-\infty}^{+\infty} \chi(\tau) \cos(\omega \tau) \,d\tau + i \int_{-\infty}^{+\infty} \chi(\tau) \sin(\omega \tau) \,d\tau 
  = \chi_p(\omega) + i \chi_d(\omega)
\end{align}
dove abbiamo riconosciuto la parte pari e dispari dovuta all'integrazione per le funzioni trigonometriche e alla unicità della decomposizione. Imponendo 
che $ \chi(\tau) $ sia nulla per $ \tau < 0 $, si ha che $ \chi_p(\tau) = - \chi_d(\tau) $ per $ \tau < 0 $, mentre $ \chi_p(\tau) = \chi_d (\tau) $ per 
$ \tau > 0 $. Dobbiamo tradurre questa relazione in frequenza, che sarebbe diretto se applicassimo la \textit{trasformata di Hilbert}, tuttavia è fuori dai 
nostri scopi introdurre un'altra trasformata, e quindi ci viene in aiuto l'analisi complessa. 
Scriviamo la frequenza (che fino a ora era reale) come numero complesso:
\begin{align}
  \omega = \omega_1 + i \omega_2.
\end{align}
Sostituiamo questa nuova frequenza nell'integrale ottenendo 
\begin{align}
  \chi(\omega) = \int_{0}^{+\infty} \chi(\tau) e^{i\omega_1 \tau} \cdot e^{-\omega_2 \tau} \,d\tau
\end{align}
che nel semipiano positivo del piano complesso ($ \omega_2 > 0 $) è assolutamente convergente, e per la convergenza dominata anche derivabile, per queste 
ragioni la funzione $ \chi(\omega) $ è \textit{olomorfa}. \begin{figure}[ht]
\usetikzlibrary{decorations.markings, arrows.meta}
\definecolor{kkcontour}{RGB}{216,90,48}
\definecolor{kkpole}{RGB}{226,75,74}
\definecolor{kkshade}{RGB}{55,138,221}
  \centering
  \begin{tikzpicture}[>=Stealth, scale=0.95,
      contourarrow/.style={postaction={decorate},
        decoration={markings,
          mark=at position 0.13 with {\arrow{Stealth[length=2.4mm]}},
          mark=at position 0.30 with {\arrow{Stealth[length=2.0mm]}},
          mark=at position 0.36 with {\arrow{Stealth[length=2.4mm]}},
          mark=at position 0.70 with {\arrow{Stealth[length=2.4mm]}}}}]

    % regione di analiticità (semidisco superiore)
    \fill[kkshade!8] (4,0) arc (0:180:4) -- cycle;

    % assi
    \draw[->, gray!60, thin] (-4.9,0) -- (5.0,0) node[below right, black]{$\operatorname{Re}\omega$};
    \draw[->, gray!60, thin] (0,-0.7) -- (0,5.0) node[left, black]{$\operatorname{Im}\omega$};

    % contorno di integrazione (asse reale + indentazione + arco grande), orientato in senso antiorario
    \draw[very thick, kkcontour, contourarrow]
          (-4,0) -- (1.65,0)
          arc (180:0:0.35)      % indentazione sopra il polo, da sinistra a destra
          -- (4,0)
          arc (0:180:4);        % semicerchio grande, da destra a sinistra

    % polo
    \fill[kkpole] (2,0) circle (2.2pt);
    \node[below=2pt] at (2,0) {$\omega_0$};
    \node[below left=-1pt] at (0,0) {$O$};

    % etichette
    \node[kkcontour] at (3.25,3.05) {$\Gamma$};
    \node[kkcontour, right] at (2.25,0.62) {$C_\eps$};
    \node[align=center] at (-1.7,2.0)
          {$\chi(\omega)$ olomorfa};

  \end{tikzpicture}
  \caption{Dominio di integrazione di $\chi(\omega)/(\omega-\omega_0)$ per le relazioni di
  Kramers--Kronig. Il contorno percorre l'asse reale, aggira il polo semplice in $\omega_0$
  con la semicirconferenza $C_\eps$ e si chiude con il semicerchio $\Gamma$ nel
  semipiano superiore, dove $\chi$ è olomorfa per causalità. L'orientazione è antioraria.}
  \label{fig:contorno-kk}
\end{figure}

Fissato $ \omega_0 \in \R $, vogliamo valutare il seguente integrale:
\begin{align}
  \oint_\gamma \frac{\chi(\omega)}{\omega - \omega_0} \,d\omega
\end{align}
con $ \gamma $ una curva semplice, chiusa, rappresentata in figura \ref{fig:contorno-kk}, formata dall'unione di $ C $, $ \Gamma $ e la semiretta reale 
(escluso il punto $ \omega_0 $), con $ \Gamma $ la semicirconferenza di raggio $ R $ centrata nell'origine del piano complesso. 
Spezziamo l'integrale sulla curva $ \gamma $ nelle tre curve che la compongono:
\begin{align}
  \oint_\gamma  = \int_{C_\eps} + \int_\Gamma + \int_{-\infty}^{+\infty} 
\end{align}
e notiamo subito che membro di sinistra (l'integrale originale) è nullo per il teorema di Cauchy-Goursat, così come l'integrale sulla semicirconferenza 
$ \Gamma $ per il decadimento all'infinito (dopo ritorneremo su questo punto). Gli altri due termini sono non nulli, e sviluppiamo il primo ottenendo
\begin{align}
  \int_{C_\eps} \frac{\chi(\omega)}{\omega - \omega_0} \,d\omega 
  = - \int_0^\pi \frac{\chi(\omega_0 + \rho e^{i\theta})}{\rho e^{i\theta}} i \rho e^{i \theta} \,d\theta
  \overset{\rho \to 0}{=} - i \pi \chi(\omega_0) 
\end{align}
dove abbiamo usato la sostituzione $ \omega - \omega_0 = \tilde{\omega} = \rho e^{i \theta} $ e il fatto che il raggio del cerchietto va a $ 0 $. 
Abbiamo quindi che 
\begin{align}
  \int_{-\infty}^{+\infty} \frac{\chi(\omega)}{\omega - \omega_0} \,d\omega 
  \coloneq \fint_{-\infty}^{+\infty} \frac{\chi(\omega)}{\omega - \omega_0} \,d\omega 
  = i \pi \chi(\omega_0) = \pi \left( -\chi_2(\omega_0) + i \chi_1(\omega_0) \right)
\end{align}
dove abbiamo anche definito il simbolo di "integrale parte principale" con $ \fint $. Non resta che esplicitare le relazioni eguagliando parte reale e 
parte immaginaria:
\begin{align}
  \begin{cases}
    \DS \chi_1(\omega_0) =\frac1\pi \fint_{-\infty}^{+\infty} \frac{\chi_2(\omega)}{\omega - \omega_0} \,d\omega 
    \\
    \DS \chi_2(\omega_0) = - \frac1\pi \fint_{-\infty}^{+\infty} \frac{\chi_1(\omega)}{\omega - \omega_0} \,d\omega 
  \end{cases} \tag*{[Kramers-Kronig]}
\end{align}
che sono dette \textit{relazioni di Kramers-Kronig}. A questo punto usiamo il fatto che la funzione del tempo $ \chi(\tau) $ era reale, e che passando in 
trasformata ci diventa
\begin{align}
  \chi^*(\omega) = \int\chi^*(\tau) e^{-i\omega \tau} \,d\tau = \int \chi(\tau) e^{-i\omega \tau} \,d\tau = \chi(-\omega)
\end{align}
ma allora possiamo usare questa per rielaborare le relazioni di Kramers-Kronig nel seguente modo:
\begin{align}
  \chi_1(\omega_0) &= \frac1\pi \left[ \fint_{0}^{+\infty} \frac{\chi_2(\omega)}{\omega - \omega_0} \,d\omega 
  + \fint_{-\infty}^0 \frac{\chi_2(\omega)}{\omega - \omega_0} \,d\omega \right]
  \\
  &= \fint_{0}^{+\infty} \frac{\chi_2(\omega)}{\omega - \omega_0} \,d\omega 
  + \fint_{0}^{+\infty} \frac{\chi_2(\omega)}{\omega + \omega_0} \,d\omega
  \\
  &= \frac2\pi \fint_{0}^{+\infty} \frac{ \omega\, \chi_2(\omega)}{\omega^2 - \omega_0^2} \,d\omega
\end{align}
e in modo analogo per la componente immaginaria che ci porta alle relazioni
\begin{align}
  \begin{cases}
    \DS \chi_1(\omega_0) =\frac2\pi \fint_{0}^{+\infty} \frac{\omega \,\chi_2(\omega)}{\omega^2 - \omega_0^2} \,d\omega 
    \\
    \DS \chi_2(\omega_0) = - \frac{2\omega_0}\pi \fint_{0}^{+\infty} \frac{\chi_1(\omega)}{\omega^2 - \omega_0^2} \,d\omega.
  \end{cases}
\end{align}
Ricondurci a questa forma ci ha permesso di esprimere le relazioni soltanto rispetto alle frequenze "fisiche", cioè quelle positive.

Nel limite di alte frequenze, come abbiamo visto nei conduttori, ci aspettiamo che ogni elettrone risponda come se fosse libero, dunque per 
$ \omega $ molto più grande di tutte le scale energetiche la funzione di dielettrica di risposta $ \eps \to 1 - \frac{\omega_p^2}{\omega^2} $. Facendo 
il limite nelle relazioni di Kramers-Kronig si ottiene che $ \omega'^2 - \omega^2 \to - \omega^2 $ e quindi
\begin{gather}
  - \frac{\omega_p^2}{\omega^2} = \eps_1(\omega) - 1 = - \frac{2}{\pi \omega^2} \int_{0}^{+\infty} \omega' \eps_2(\omega) \,d\omega' 
  \\ 
  \implies \int_{0}^{+\infty} \omega' \eps_2(\omega) \,d\omega' = \frac\pi{2} {\omega_p^2} = \frac{ \pi N e^2 }{2 \eps_0 m_e}
\end{gather}
dove si evince che l'assorbimento totale pesato per le frequenze è dettato dalla densità dei portatori di carica. 
Per quanto riguarda la parte immaginaria $ \eps_2 $ c'è da aggiungere un termine nel caso dei metalli, questo perché in generale la $ \eps $ codificava 
la corrente di spostamento in termini del campo elettrico, ma nei metalli abbiamo anche una corrente di conduzione sempre dovuta al campo elettrico, cosa 
che si vede nell'equazione di Maxwell per il rotore di $ \v{H} $:
\begin{align}
  \nabla \times \v{H} = \v{J} + \partial_t \v{D} = \sigma(\omega) \v{E} - i\omega \eps(\omega) \v{E}.
\end{align}
Può essere utile raccogliere queste due risposte in una unica, ridefinendo la $ \eps_{\text{leg}}(\omega) $ come la risposta dielettrica dovuta alla 
polarizzazione si ottiene
\begin{align}
  \eps_r(\omega) = \eps_{r, \text{leg}}(\omega) + \frac{i \sigma(\omega)}{\eps_0 \omega}.
\end{align}
In modo analogo, si ricava che nel limite di alta frequenza, utilizzando il limite della conducibilità dal modello di Drude, si ricava 
una regola si somma anche per la reale:
\begin{align}
  \int_{0}^{+\infty} (\eps_1(\omega') - 1) \,d\omega' = - \frac{\pi \sigma_0}{2 \eps_0 }.
\end{align}
Queste due identità che determinano la somma delle proprietà ottiche di un mezzo materiale costituiscono un controllo sulla validità di dati acquisiti, 
poiché dipendono da dalle caratteristiche costitutive del materiale.



\section{Onde incidenti} \label{sec: onde incidenti}
Vogliamo studiare come si comportano i mezzi quando facciamo incidere un onda all'interfaccia, il più semplice setup sperimentale è quello realizzato in 
figura \ref{figure:_onda_incidente}. \begin{figure}
    \center
    \begin{tikzpicture}
        
        \draw[->] (0,-4) -- (0,4) node[left] {$ z $};
        \draw[->] (-5,0) -- (5,0) node[below] {$ x $};

        \draw[ domain = -3:0, variable = \x, ->,thick] plot ({\x}, {\x*2/3});

        \draw[ domain = 0:2, variable = \x, ->,thick, orange] plot ({\x}, {\x*2});

        \draw[ domain = 0:3, variable = \x, ->,thick, blue] plot ({\x}, {-\x*2/3});

        \node at (-5,4) {Mezzo: $ \tilde{\eps}_b(\omega) $};

        \node at (-5,-4) {Vuoto: $ \eps_a $};
    \end{tikzpicture}
    \caption{Onda che incide all'interfaccia di un mezzo.}
    \label{figure:_onda_incidente}
\end{figure} 
Consideriamo il "mezzo a" come un mezzo semplice, cioè non assorbente, mentre il mezzo di trasmissione "mezzo b" lo consideriamo non magnetico quindi con 
$\mu=\mu_0 $ inoltre lavoriamo come sempre con onda piana incidente indicando con $ \v{k}_i $ il suo vettore d'onda e con $ \theta_i $ il suo angolo 
di incidenza rispetto alla normale. Definiamo il piano di incidenza come il piano generato dai vettori $ \v{k}_i, \vers{n} $, cioè lo span del vettore 
d'onda incidente e del versore normale alla superficie.

Facendo incidere una onda piana sul mezzo si generano due onde, questo perché devono essere rispettate le condizioni di raccordo all'interfaccia, cioè 
devono conservarsi le componenti $ E_\parallel, \mu_0 H_\parallel = B_\parallel $. Siccome devono essere verificate a tutti i tempi, ricaviamo subito che 
la frequenza $ \omega $ dell'onda incidente deve essere la stessa dell'onda riflessa e trasmessa. Per l'invarianza traslazionale sulle $ y $ abbiamo che 
anche $ k_{y} $ si conserva dunque è $ 0 $ per ogni onda. Se nell'onda trasmessa è presente la componente immaginaria del vettore d'onda, ci aspettiamo 
che sia solamente lungo $ z $ cioè $ \v{k}_2 = k_2 \vers{z} $, questo perché ci aspettiamo che l'onda si attenui man mano che procede all'interno del mezzo, 
mentre sempre per l'invarianza traslazionale la componente $ k_x $ deve conservarsi, cioè $ k_{ix} = k_{rx} = k_{tx} $. Una visualizzazione più matematica 
di ciò si ha considerando che le tre onde, incidente, riflessa e trasmessa, sono onde piane cioè $ \v{E} e^{i(\v{k\cdot r } - \omega t)} $, e le condizioni 
di raccordo sono relazioni lineari fra queste, che significa che avremo per ogni condizione una equazione del tipo:
\begin{align}
  A e^{i(\v{k}_i \cdot \v{r} - \omega t)} + B e^{i(\v{k}_r \cdot \v{r} - \omega t)} = C e^{i(\v{k}_t \cdot \v{r} - \omega t)}
\end{align} 
tutte da valutare alla superficie di incidenza che nel nostro caso è $ z = 0 $. Queste equazioni ammettono dei coefficienti costanti come soluzione 
soltanto nel caso in cui gli esponenziali sono uguali e cioè la fase $ \phi = (\v{k\cdot r } - \omega t) $ deve essere la stessa per tutti, 
ma questo significa proprio
\begin{align}
  \omega, k_x, k_y \;\; \textnormal{conservati.}
\end{align}


Adesso abbiamo abbastanza informazioni per 
imporre la relazione di dispersione che è 
\begin{align}
  \v{k} \cdot \v{k} = \tilde{\eps}\, \frac{\omega^2}{c^2}
\end{align}
e abbiamo quindi per l'onda riflessa che 
\begin{align}
  k_x^2 + k_{iz}^2 = \eps_a \frac{\omega^2}{c^2} = k_x^2 + k_{rz}^2  \implies k_{rz} = \pm k_{iz}
\end{align}
e il segno che dobbiamo scegliere è il $ - $ poiché deve andare nel verso opposto (delle $ z $) dell'onda incidente, quindi $ k_{rz} = -k_{iz} $. 
Da questa informazione ricaviamo che l'angolo dell'onda riflessa è lo stesso dell'onda incidente: $ \theta_i = \theta_r $.
Per l'onda trasmessa abbiamo
\begin{align}
  k_x^2 + k_{tz}^2 = \tilde{\eps}_b \frac{\omega^2}{c^2} \implies k_{tz} = \pm \sqrt{\tilde{\eps}_b \frac{\omega^2}{c^2} - k_x^2} \in \mathbb{C}
\end{align}
e qui invece scegliamo il segno $ + $ perché vogliamo che allontanandosi decada, quindi $ k_{tz} = \sqrt{\tilde{\eps}_b \frac{\omega^2}{c^2} - k_x^2} $.

\textbf{Caso di mezzo trasparente} ($ \eps_b \in \mathbb{R} $): l'onda trasmessa è omogenea nel suo mezzo quindi definiamo un angolo $ \theta_t $ e 
imponiamo la conservazione di $ k_x $:
\begin{align}
  k_{ix} = k_{tx} \implies n_a \frac{\omega}{c} \sin\theta_i &= n_b \frac{\omega}{c} \sin\theta_t 
  \\
  \implies n_a \sin\theta_i &= n_b \sin\theta_t \tag*{[Legge di Snell]}
\end{align}
dove abbiamo ottenuto la nota legge di Snell-Cartesio. Sostituiamo in $ k_{tz} $ e abbiamo
\begin{align}
  k_{tz} = \sqrt{{\eps}_b \frac{\omega^2}{c^2} - \eps_a \frac{\omega^2}{c^2}\sin^2\theta_i}= \frac{\omega}{c}\sqrt{{\eps}_b - \eps_a\sin^2\theta_i } 
\end{align}
e in questa forma notiamo che se $ n_a > n_b $, per $ \sin\theta_i > \frac{n_b}{n_a} $ abbiamo che $ k_{tz} \in \mathbb{C} $ puramente immaginario, dunque 
si verifica un fenomeno detto \textit{riflessione totale interna}, dove anche se non è presente alcun processo dissipativo l'onda trasmessa è soppressa 
esponenzialmente.

Abbiamo caratterizzato tutti i gradi di libertà cinematici, per ottenere le ampiezze delle onde abbiamo come gradi di libertà l'ampiezza dell'onda incidente 
e la polarizzazione, che diremo essere: TE quando $ \v{E}_i $ è $ \perp $ al piano di incidenza; TM quando $ \v{B}_i $ è $ \perp $ al piano di incidenza. 
Il motivo per cui dividiamo in onde TE e TM è perché ogni onda è possibile scriverla come somma di questi due modi (sono una base delle polarizzazioni), 
e soprattutto perché le onde risultanti mantengono la polarizzazione di partenza, e questo non è vero in generale perché mandando un'onda che è una somma 
delle due polarizzazioni, benché inizialmente lineare, le onde risultanti otterranno in generale una polarizzazione ellittica.

\textbf{Polarizzazione TE}: $ \v{E} = E_y \vers{y} $, le condizioni di raccordo sono che $ E_\parallel,B_\parallel $ devono essere continue, è opportuno 
scriverci $ B_\parallel $ in funzione di $ \v{E} $:
\begin{align}
  i\omega \v{B} &= \nabla \times \v{E} = \partial_x E_y \vers{z} - \partial_z E_y \vers{x} = i k_x E_y \vers{z} - ik_z E_y \vers{x}
  \\
  &\implies \v{B} = \frac{k_z}{\omega} E \vers{x} - \frac{k_x}{\omega} E \vers{z}
\end{align}
l'onda incidente è della forma
\begin{align}
  E_{iy} = E_i \, \underline{e^{ik_x x}}\, e^{ik_z z}\, \underline{e^{-i\omega t} }
\end{align}
dove i termini sottolineati sono già conservati, mentre per il termine che dipende da $ z $ lo valutiamo all'interfaccia quindi in $ z = 0 $. 
Imponiamo le leggi di conservazione che abbiamo trovato:
\begin{align}
  \begin{cases}
    E_i + E_r = E_t  \\
    B_i + B_r = B_t \implies k_{iz} (E_i - E_r) = k_{tz} E_t
  \end{cases}
\end{align}
risolvendo otteniamo 
\begin{align}
  E_r = \frac{k_{iz} - k_{tz}}{k_{iz} + k_{tz}}\, E_i
  \\
  E_t = \frac{2k_{iz}}{k_{iz} + k_{tz}}\, E_i
\end{align}
che sono detti \textit{coefficienti di Fresnel}. Li riscriviamo in termini degli indici di rifrazione e degli angoli delle onde:
\begin{align}
  E_r = \frac{n_a \cos\theta_i - n_b \cos\theta_t}{n_a \cos\theta_i + n_b \cos\theta_t}\, E_i \eqcolon r_s \,E_i
  \\
  E_t = \frac{2 n_a \cos \theta_i}{n_a \cos\theta_i + n_b \cos\theta_t}\, E_i \eqcolon t_s \,E_i
\end{align}

\textbf{Polarizzazione TM}: $ \v{B} = B_y \vers{y} $, ci scriviamo $ \v{E} $ in funzione di $ \v{B} $:
\begin{align}
  -i \frac{\omega \, n^2}{c^2} \v{E} &= \nabla \times \v{B} = \partial_x B_y \vers{z} - \partial_z B_y \vers{x}
  \\
  \implies \v{E} &= \frac{c^2}{\omega \, n^2} \left( k_z B \vers{x} - k_x B \vers{z} \right) \quad \textnormal{con} \quad E = \frac{c}{n} B
\end{align}
e imponiamo le stesse condizioni di raccordo
\begin{align}
  \begin{cases}
    E_{i\parallel} + E_{r\parallel} = E_{t\parallel} \implies \frac{k_{iz}}{n^2_a} \left( B_i - B_r \right) = \frac{k_{tz}}{n^2_b} B_t
    \\
    B_i + B_r = B_t 
  \end{cases}
\end{align}
che risolvendo ci porta a 
\begin{align}
  B_r = \frac{n_b \cos\theta_i - n_a \cos \theta_t }{n_b \cos\theta_i + n_a \cos \theta_t}\, B_i
  \\
  B_t = \frac{2 n_b \cos\theta_i}{n_b \cos\theta_i + n_a \cos \theta_t}\, B_i
\end{align}
dove abbiamo sfruttato la relazione $ k_{jz} = k_j \cos\theta_j = n_j \frac{\omega}{c} \cos \theta_j  $. Per ricavare il campo elettrico usiamo la relazione 
fra $ E $ e $ B $ ricavata sopra ottenendo: 
\begin{align}
  E_r = \frac{n_b \cos\theta_i - n_a \cos \theta_t }{n_b \cos\theta_i + n_a \cos \theta_t}\, E_i \eqcolon r_p \, E_i
  \\
  E_t = \frac{2 n_a \cos\theta_i}{n_b \cos\theta_i + n_a \cos \theta_t}\, E_i \eqcolon t_p \,E_i
\end{align}
Di seguito riportiamo i 4 coefficienti che abbiamo trovato,
\begin{align}
  {r_s = \frac{n_a \cos\theta_i - n_b \cos\theta_t}{n_a \cos\theta_i + n_b \cos\theta_t}} 
  \quad &{t_s = \frac{2 n_a \cos \theta_i}{n_a \cos\theta_i + n_b \cos\theta_t}}
  \\
  {r_p = \frac{n_b \cos\theta_i - n_a \cos \theta_t }{n_b \cos\theta_i + n_a \cos \theta_t}} 
  \quad &{t_p = \frac{2 n_a \cos\theta_i}{n_b \cos\theta_i + n_a \cos \theta_t}}
\end{align}
I coefficienti di Fresnel sono definiti come rapporti fra le intensità dei campi elettrici, per ottenere i rispettivi campi magnetici è opportuno usare la 
relazione $ E = \frac{c}{n} B $ nel relativo mezzo.

\subsection{Incidenza normale}
Consideriamo adesso il caso in cui $ \theta_i = 0 $, cioè l'onda incidente è ortogonale alla superficie. In questa situazione abbiamo che 
$ \v{k}_i = k_{iz} \vers{z} $ e dalla relazione di dispersione vale $ k_{iz} = \sqrt{\eps_a} \frac{\omega}{c} = n_a \frac{\omega}{c} $, e 
$ k_t = \frac{\omega}{c} \sqrt{\tilde{\eps}_b} = \frac{\omega}{c} \tilde{n}_b $. I coefficienti di Fresnel diventano
\begin{align}
  E_t = \frac{2n_a}{n_a + \tilde{n}_b}\, E_i
  \\
  E_r = \frac{n_a - \tilde{n}_b}{n_a + \tilde{n}_b}\, E_i
\end{align}
L'intensità di onda piana in un mezzo abbiamo visto essere
\begin{align}
  I = \langle |\v{S}| \rangle = \frac{n_1}{2 c \mu_0} E^2 = \frac{n_1}{2 Z_0} E_0^2
\end{align}
dove abbiamo definito l'impedenza di vuoto $ Z_0 \coloneq \sqrt{\frac{\mu_0}{\eps_0}} \simeq \SI{377}{\ohm} $, 
e da questa intensità definiamo due coefficienti detti \textit{coefficiente di riflessione} $ R $ (\textit{riflettanza}) e 
\textit{coefficiente di trasmissione} $ T $ (\textit{trasmittanza}):
\begin{align}
  R \coloneq \frac{I_r}{I_i} \quad T \coloneq \frac{I_t}{I_i}.
\end{align}
A questo punto utilizziamo le relazioni di Fresnel per scrivere i due coefficienti in funzione degli indici di rifrazione:
\begin{align}
  R &= \left| \frac{E_r}{E_i} \right|^2 = |r|^2 = \frac{|n_a - \tilde{n}_b|^2}{|n_a + \tilde{n}_b|^2}
  \\
  T &= \frac{n_{bR}}{n_a} \left| \frac{E_t}{E_i} \right|^2 = \frac{n_{bR}}{n_a} |t|^2 = \frac{n_{bR}}{n_a} \frac{(2n_a)^2}{|n_a + \tilde{n}_b|^2} 
  = \frac{4 n_a n_{bR}}{|n_a + \tilde{n}_b|^2}
\end{align}
È interessante notare che la somma dei due coefficienti è $ R + T = 1 $, dato che il mezzo è non assorbente, e infatti è cioè che si ottiene provando a 
sommarli.

% -----------------------------------------------------------------------------------------------------------------
% --------------------------------------------------- APPENDICE ---------------------------------------------------
% -----------------------------------------------------------------------------------------------------------------
\appendix

\chapter{Identità vettoriali}

\section{Operazioni generiche}
Di seguito alcune delle identità vettoriali utilizzate nel testo.\\
Triplo prodotto:
\begin{itemize}
  \item $ \v{ a \times (b \times c) } = \v{ b \cdot (a \cdot c) } - \v{ c\cdot (a \cdot b) } $
  \item $ \v{ a \cdot \left( b \times c \right) } = \v{ b \cdot \left( c \times a \right) } = \v{ c \cdot \left( a \times b \right) } $
\end{itemize}

\section{Operatori differenziali}
\begin{align}
  \label{identity: div(A times B)}
  &\nabla \cdot \v{ (A \times B) } = \v{ B \cdot (\nabla \times A) } - \v{ A \cdot (\nabla \times B) }
  \\
%  \label{identity: rot(A times B)}
%  &\nabla \times \v{ (A \times B) } =  
%  \\
  \label{identity: rot(rot A)}
  &\nabla \times \v{ (\nabla \times \v{ A }) } = \nabla (\nabla \cdot \v{ A }) - \nabla ^2 \v{ A }
\end{align}
\begin{align}
  [\nabla  \cdot \left( \v{ A \times B } \right)]_i &= \partial_i \left( \v{ A \times B } \right)_i = \partial_i \left( \eps_{i j k} A_j B_k \right) 
  = \eps_{i j k} \partial_i (A_j B_k) 
  \\
  &= \eps_{i j k} \left( B_k \partial_i A_j + A_j \partial_i B_k\right) = B_k \, \eps_{k i j} \, \partial_i A_j - A_j \, \eps_{j i k} \, \partial_i B_k 
  \\
  &= [\v{ B } \cdot (\nabla\!\times\!\v{ A })]_k - [\v{ A } \cdot (\nabla\!\times\!\v{ B })]_j
\end{align}
dove abbiamo sfruttato le proprietà di permutazione del tensore di Levi-Civita e la convenzione sugli indici ripetuti.

\textbf{Divergenza di scalare per vettore}:
\begin{align} \label{identity:_divergenza(scalare_per_vettore)}
  \nabla \cdot \left( \alpha \, \v{V} \right) = (\nabla \alpha) \cdot \v{V} + \alpha (\nabla\!\cdot\!\v{ V })
\end{align}


\chapter{Operatori differenziali in vari sistemi di coordinate}
Di seguito elenchiamo le espressioni dei più comuni operatori differenziali esprimibili con $ \nabla $ nei sistemi di coordinate cartesiane, cilindriche e 
sferiche.

\begin{table}[h]
\centering
% \footnotesize
\renewcommand{\arraystretch}{2.3}

\begin{tabular}{>{\raggedright\arraybackslash}p{2.2cm} 
                >{\raggedright\arraybackslash}p{4.2cm} 
                >{\raggedright\arraybackslash}p{4.6cm} 
                >{\raggedright\arraybackslash}p{4.6cm}}

\toprule
Operatore & Cartesiane $(x,y,z)$ & Cilindriche $(\rho,\phi,z)$ & Sferiche $(r,\theta,\phi)$ \\
\midrule

Gradiente $\nabla f$

& $\partial_x f\,\hat{x}+\partial_y f\,\hat{y}+\partial_z f\,\hat{z}$

& $\partial_\rho f\,\hat{\rho}
+\frac{1}{\rho}\partial_\phi f\,\hat{\phi}
+\partial_z f\,\hat{z}$

& $\partial_r f\,\hat{r}
+\frac{1}{r}\partial_\theta f\,\hat{\theta}
+\frac{1}{r\sin\theta}\partial_\phi f\,\hat{\phi}$

\\

Divergenza $\nabla\!\cdot\!\v A$

& $\partial_x A_x+\partial_y A_y+\partial_z A_z$

& $\frac{1}{\rho}\partial_\rho(\rho A_\rho)
+\frac{1}{\rho}\partial_\phi A_\phi
+\partial_z A_z$

& $\frac{1}{r^2}\partial_r(r^2A_r)
+\frac{1}{r\sin\theta}\partial_\theta(A_\theta\sin\theta)
+\frac{1}{r\sin\theta}\partial_\phi A_\phi$

\\

Rotore $\nabla\times\v A$

& $\begin{aligned}
(\partial_yA_z-\partial_zA_y)\hat{x} \\
+(\partial_zA_x-\partial_xA_z)\hat{y} \\
+(\partial_xA_y-\partial_yA_x)\hat{z}
\end{aligned}$

& $\begin{aligned}
(\tfrac{1}{\rho}\partial_\phi A_z-\partial_z A_\phi)\hat{\rho}\\
+(\partial_zA_\rho-\partial_\rho A_z)\hat{\phi}\\
+\tfrac{1}{\rho}(\partial_\rho(\rho A_\phi)-\partial_\phi A_\rho)\hat{z}
\end{aligned}$

& $\begin{aligned}
\tfrac{1}{r\sin\theta}(\partial_\theta(A_\phi\sin\theta)-\partial_\phi A_\theta)\hat r\\
+\tfrac{1}{r}(\tfrac{1}{\sin\theta}\partial_\phi A_r-\partial_r(rA_\phi))\hat\theta\\
+\tfrac{1}{r}(\partial_r(rA_\theta)-\partial_\theta A_r)\hat\phi
\end{aligned}$

\\

Laplaciano $\nabla^2 f$

& $\partial_x^2f+\partial_y^2f+\partial_z^2f$

& $\frac{1}{\rho}\partial_\rho(\rho\partial_\rho f)
+\frac{1}{\rho^2}\partial_\phi^2 f
+\partial_z^2 f$

& $\frac{1}{r^2}\partial_r(r^2\partial_r f)
+\frac{1}{r^2\sin\theta}\partial_\theta(\sin\theta\,\partial_\theta f)
+\frac{1}{r^2\sin^2\theta}\partial_\phi^2 f$

\\

\bottomrule
\end{tabular}

\end{table}

\chapter{Decomposizione di Helmholtz}
La decomposizione di Helmholtz può essere vista come un caso particolare della decomposizione di Hodge (che non tratteremo). 
Questa decomposizione ci dice che un campo sufficientemente regolare può essere scomposto come somma di due campi, 
uno a rotore nullo e uno a divergenza nulla, cioè come somma di un campo solenoidale e un campo conservativo.

Sia $ \v{ F } \in C^2, \quad \v{ F }: \mathbb{R}^3 \to \mathbb{R}^3$, allora $ F $ è tale che:
\begin{equation}
  \v{ F } = \v{\nabla}\times\v{ A } - \nabla \varphi
\end{equation}


\chapter{
  \texorpdfstring
  {Sviluppo di $ \frac{1}{\v{ |r-r'| }}$}
  {Sviluppo di 1 su r - r'}}
Vogliamo sviluppare questo oggetto in coordinate sferiche perché:
\begin{itemize}
  \item rispetto ad altri sistemi di coordinate e soprattutto rispetto alle cartesiane, gli sviluppi successivi facilitano i conti;
  \item è una funzione armonica quasi dappertutto, perché
  \begin{equation}
    \nabla^2 \left( \frac{1}{\v{ |r-r'| }} \right) = -4\pi \delta(\v{ r-r' })
  \end{equation}
  cioè è nullo (il suo laplaciano) $ \forall \v{ r } \not= \v{ r' } $.
\end{itemize}
Scegliamo un sistema di coordinate dove $ \v{ r' } $ giace sull'asse polare, in questo modo siamo nel caso di simmetria azimutale e quindi il problema 
è notevolmente più semplice e soprattutto non richiede di invocare lo sviluppo in armoniche sferiche.
Dallo studio dell'equazione di Laplace con simmetria azimutale ($ m = 0 $) sappiamo che ogni funzione armonica si può esprimere come
\begin{equation}
  F (\v{ r }) = \sum_{l=0}^{\infty} \left( A_l \, r^l + \frac{B_l}{r^{l+1}} \right) P_l\left( \cos \theta \right)
\end{equation}
Svolgo i limiti
\begin{itemize}
  \item $\displaystyle \lim_{r \to \infty} \frac{1}{\v{ |r-r'| }} = 0 \overset{!}{=} \lim_{r \to \infty} F(\v{ r }) \quad \implies A_l = 0 \quad \text{per} \quad r>r'$
  \item $\displaystyle \lim_{r \to 0} \frac{1}{\v{ |r-r'| }} = \frac{1}{|r'|} \overset{!}{=} \lim_{r \to 0} F(\v{ r }) \quad  \implies B_l = 0 \quad \text{per} \quad r<r' $
\end{itemize}
Sfruttando lo sviluppo di Taylor (con $ \v{ r } $ e $ \v{ r' } $ allineati) si ottiene che
\begin{equation}
  \frac{1}{ r_a - r_b } = \frac{1}{r_a} \frac{1}{1- \frac{r_b}{r_a}} = \frac{1}{r_a} \sum_{n = 0}^{\infty} \left( \frac{r_b}{r_a} \right)^n =  \frac{1}{r_a} \left( 1 + \frac{r_b}{r_a} + \left( \frac{r_b}{r_a} \right)^2 + ... \right)
\end{equation}
dove $ r_a \; (r_b) $ è il maggiore (minore) fra $ r $ e $ r' $.\\
Per $ r>r' $:

\begin{align}
  \sum_{l=0}^{\infty} \left( \frac{B_l}{r^{l+1}} \right) & \overset{!}{=} \frac{1}{r} \sum_{n = 0}^{\infty} \left( \frac{r'}{r} \right)^n \quad \implies B_l = \left( r' \right)^l \\
  F(\v{ r }) &= \sum_{l=0}^{\infty} \frac{1}{r} \left( \frac{r'}{r} \right) ^l P_l \left( \cos \theta \right)
\end{align}

Per $ r<r' $:

\begin{align}
  \sum_{l=0}^{\infty} \left( A_l \; r^l \right) & \overset{!}{=} \frac{1}{r'} \sum_{n = 0}^{\infty} \left( \frac{r}{r'} \right)^n \quad \implies A_l = \left( \frac{1}{r'} \right)^{l+1} \\
  F(\v{ r }) &= \sum_{l=0}^{\infty} \frac{1}{r'} \left( \frac{r}{r'} \right) ^l P_l \left( \cos \theta \right)
\end{align}

Queste due espressioni possono essere compattate nella forma:
\begin{equation} \label{sviluppo_1/R_compatto}
  F(\v{ r }) = \sum_{l=0}^{\infty} \frac{r_{<}^{l}}{r_{>}^{l+1}} \; P_l \left( \cos \theta \right)
\end{equation}
che va intesa come: $ r_< = \min\{r, r'\} , \quad r_> = \max\{r, r'\} $

\subsection{Alcuni termini}
Sviluppo fino al 2° ordine della \eqref{sviluppo_1/R_compatto}:

\begin{align}
  F(\v{ r }) &= \frac{1}{r} + \frac{r'}{r^2} \cos \theta + \frac{r'^2}{r^3} \left( \frac{3\cos ^2 \theta -1}{2} \right) + ...\\
  &= \frac{1}{r} + \frac{\v{ r' \cdot r }}{r^3} + \frac{r'^2}{r^3} \left( \frac{3\cos ^2 \theta -1}{2} \right) + ...
\end{align}


\chapter{Pseudo-vettori}
A differenza dei vettori "veri" che trasformano normalmente
\begin{equation}
  \v{ V' } = T\v{ V } \implies \; V'_i = T_{i j}V_j
\end{equation}
dove nell'ultima relazione è implicita la sommatoria sull'indice ripetuto, gli pseudo-vettori trasformano le loro componenti in base segno della 
trasformazione applicata:
\begin{equation}
  \v{ A' } = \det \!T \left( T\v{ A } \right) \implies \; A'_i = \det\!T \, T_{i j} A_j
\end{equation}
Questo è molto utile per escludere fin dal principio alcune componenti di vettori e pseudo-vettori:
ad esempio se il sistema è pari sotto lo scambio di una coordinata (come $ x \to -x $), che è una riflessione quindi con determinante $ -1 $, 

\chapter{Richiami di relatività}

\section{Covarianza e controvarianza}
Non tutti gli oggetti matematici trasformano allo stesso modo; per convincerci di ciò consideriamo un vettore $ \v{v} $ scritto in una base 
$ \v{e}_1, \v{e}_2 $:
\begin{align}
  \v{v} = v^1\v{e}_1 + v^2 \v{e}_2
\end{align}
e supponiamo di fare un riscalamento dei vettori di base nel seguente modo: $ \v{e}_1 \mapsto \v{e'}_1 = 2\v{e}_1, \v{e}_2 \mapsto \v{e'}_2 = 5\v{e}_2  $. 
Dato che abbiamo aumentato i vettori di base allora il vettore $ \v{v'} = \left. \v{v} \right|_{\{ \v{e'}_1, \v{e'}_2\}} $ deve trasformate le sue 
componenti affinché rimanga lo stesso vettore
\begin{align}
  &\v{v'} = v'^1 \v{e'}_1 + v'^2 \v{e'}_2 = v^1\v{e}_1 + v^2 \v{e}_2
  \\
  &\implies v'^1 = \frac{v^1}{2} \quad v'^2 = \frac{v^2}{5}
\end{align}
È facile osservare che il vettore trasforma in modo inverso rispetto a come trasformano gli elementi della base. Il vettore trasforma secondo l'espressione 
\begin{align} \label{eq:_trasformazione_controvariante}
  {v'}^i = \sum_j \der{{v'}^i}{{v^j}} \, v^j = \sum_j J_{i j} \, v^j
\end{align}
dove $ J_{i j} $ è lo jacobiano della trasformazione, mentre la base $ \{ e_i \} $ trasforma nel modo inverso cioè
\begin{align} \label{eq:_trasformazione_covariante}
  {e'}_i = \sum_j \der{x^i}{x'^j} \, e_i = \sum_j {J^{-1}}^t_{i j} \, e_j
\end{align}
Definiamo \textit{controvariante} gli oggetti che trasformano come in \eqref{eq:_trasformazione_controvariante} e indicheremo le sue componenti con un 
indice in alto. Definiamo allo stesso modo \textit{covariante} gli oggetti che trasformano come in \eqref{eq:_trasformazione_covariante} e indicheremo le 
componenti con un indice in basso. Vediamo come trasforma il gradiente di una funzione scalare:
\begin{align}
  \nabla'_i f = \der{f}{x'^i} = \der{x^j}{x'^i} \der{f}{x^j} = \left( J^-1 \right)^t_{i j} \nabla_j f
\end{align}

\section{Quadrivettori}
Il quadrivettore posizione è definito come
\begin{align}
  x^\mu = (ct, \v{x}) = (ct, x, y, z)
\end{align}
definiamo la quadrivelocità $ u^\mu $ come
\begin{align}
  u^\mu = \frac{dx^\mu}{d\uptau} = (\gamma c, \gamma \v{ v })
\end{align}
e moltiplicando per la massa a riposo otteniamo il quadrimpulso $ p^\mu $ 
\begin{align}
  p^\mu = m_0 u^\mu = (m_0 \gamma c, m_0 \gamma \v{ v })
\end{align}
I quadrivettori hanno la norma invariante per trasformazioni di Lorentz, ma la norma non è quella euclidea bensì quella di Lorentz:
\begin{align}
  \Lambda = 
  \begin{pmatrix}
    1 & 0 & 0 & 0 \\
    0 & -1 & 0 & 0 \\
    0 & 0 & -1 & 0 \\
    0 & 0 & 0 & -1 
  \end{pmatrix}
\end{align}
Definiamo ora il quadrigradiente di uno scalare come 
\begin{align}
  \partial_\mu \phi = a_\mu = \left( \frac{1}{c} \der{}{t}, \nabla  \right) \phi
\end{align}
allo stesso modo troviamo la quadridivergenza di un quadrivettore
\begin{align}
  \partial_\mu b^\mu = \frac{1}{c} \der{b^0}{t} + \nabla \cdot \v{ b } = \frac{1}{c} \der{b^0}{t} + \nabla \cdot \v{ b }
\end{align}
Possiamo notare come effettivamente il quadrigradiente e la quadridivergenza siano rappresentati dallo stesso simbolo in questa notazione, l'unica 
differenza rimane chiaramente l'oggetto a cui si applicano. Il d'Alembertiano applicato ai quadrivettori diventa
\begin{align}
  \Box \coloneq \frac{1}{c^2} \der[2]{}{t} - \der[2]{}{x} - \der[2]{}{y} - \der[2]{}{z} = g^{\mu \nu} \partial_\mu \partial_\nu = \partial^\mu \partial_\mu
\end{align}

\section{Operazioni con gli indici}
Nella sezione precedente, nell'applicare la quadridivergenza a un quadrivettore abbiamo ottenuto uno scalare, ma più in generale si dice che abbiamo 
\textit{contratto gli indici}. La contrazione di un indice è una somma implicita su quell'indice e si può fare quando l'indice compare due volte, una 
come controvariante e una come covariante:
\begin{align}
  A^\mu B_\mu = \sum_{\mu = 0}^3 A^0B_0 + A^1B_1 + A^2B_2 + A^3B_3.
\end{align}
Questa operazione ha un suo significato, dato che quando moltiplichiamo una quantità covariante con una controvariante, le due matrici di trasformazione 
sono una l'inversa dell'altra, dunque moltiplicate danno l'identità e otteniamo quindi una quantità invariante, cioè uno scalare.

Altra operazione importante è l'alzare o abbassare gli indici: per farlo occorre la metrica $ g_{\mu \nu} $ e serve per quando abbiamo una quantità 
covariante (controvariante) e vogliamo alzare (abbassare) l'indice per farla trasformare nel modo opposto. È definita nel modo seguente:
\begin{align}
  A^\mu &\coloneq g^{\mu \nu} A_\nu
  \\
  A_\mu &\coloneq g_{\mu \nu} A^\nu
\end{align}
mentre per quanto riguarda la metrica la sua versione con gli indici alti è l'inversa cioè
\begin{align}
  g^{\mu \nu} g_{\nu \rho} = \delta^\mu{}_\rho
\end{align}
dove $ {\delta^\mu}_\rho $ è la delta di Kronecker. In relatività la $ \delta $ è definita come quella a cui siamo abituati solo quando ha un indice alto 
e uno basso, cioè
\begin{align}
  \delta^\mu{}_\nu = 
  \begin{cases}
    \DS 1 \quad \textnormal{se} \; \mu = \nu \\
    \DS 0 \quad \textnormal{se} \; \mu \neq \nu \\
  \end{cases}
\end{align}
mentre se abbassiamo (o alziamo) un indice otteniamo di nuovo la metrica:
\begin{align}
  \delta_{\mu \nu} = g_{\mu \rho}\, \delta^\rho{}_\nu = g_{\mu \rho} \, g^{\rho \alpha} \, g_{\alpha \nu} = g_{\mu \nu}
\end{align}

\section{Metrica di Minkowski}
Si dice \textit{metrica di Minkowski} il tensore metrico costante e uniforme $ g_{\mu \nu} = \diag(+1, -1, -1, -1) $ cioè
\begin{align}
  g_{\mu \nu} = 
  \begin{pmatrix}
    +1 & 0 & 0 & 0 \\
    0 & -1 & 0 & 0 \\
    0 & 0 & -1 & 0 \\
    0 & 0 & 0 & -1
  \end{pmatrix},
\end{align}
tuttavia alcuni testi invertono i segni. Questa è la metrica più semplice, sopratutto per due proprietà molto importanti: la prima è che il tensore con 
gli indici alzati è sempre $ g^{\mu \nu} = \diag(+1, -1, -1, -1) $, e la seconda è che nell'alzare o abbassare un indice cambia soltanto il segno 
della parte spaziale ovvero
\begin{align}
  A^\mu = g^{\mu \nu} A_\nu = (A_0, -\v{A}).
\end{align}

\end{document}